\documentclass{article}
\usepackage[utf8]{inputenc}
\usepackage[T1]{fontenc}
\usepackage{geometry}
\usepackage{hyperref}
\geometry{margin=1in}

\title{多空转折 手抓}
\date{}

\begin{document}
\maketitle
《跟大师学投资》丛书

蔡森薯

SPM 南方出版传媒 广东任济女版社 ·广州·

\clearpage

\section{图书在版编目（CIP）数据}

多空转折一手抓／蔡森著.一广州：广东经济出版社，2016.9 ISBN978-7-5454-4764-4

I.①多..Ⅱ.①蔡Ⅲ.①股票投资-基本知识IV.① F830.91

\section{中国版本图书馆CIP数据核字（2016）第199320号}

出版人：姚丹林 责任编辑：陈念庄罗振文 责任技编：许伟斌

\section{出版广东经济出版社（广州市环市东路水荫路11号11\textasciitilde{}12楼）}

发行 经销全国新华书店 广州家联印刷有限公司 印刷 （广州市天河区东圃镇吉山村坑尾路3-2号） 开本730毫米×1020毫米1/16 印张15.5 字数349000字 版次2016年9月第1版 印次2016年9月第1次 印数1\textasciitilde{}5000 书号ISBN978-7-5454-4764-4 定价68.00元

如发现印装质量问题，影响阅读，请与承印厂联系调换。 发行部地址：广州市环市东路水荫路11号11楼 电话：（020）3830605537601950邮政编码：510075 邮购地址：广州市环市东路水荫路11号11楼 电话：（020）37601980营销网址：http：//www.gebook.com 广东经济出版社新浪官方微博：http：//e.weibo.com/gebook 广东经济出版社常年法律顾问：何剑桥律师 ·版权所有翻印必究·

\clearpage

\section{推荐序}

精通量价观念 进出股市将如探囊取物

“天下武功出少林，少林武功出佛宗”。在技术分析领域，所有指标（KD、 RSI、MACD、DMI）皆来自于K线图，一如“天下武功出少林”，懂得K线图，便 已然具备进出股市不败境界。量是因，价为果，量价是因果，量价是技术指标中 的精髓，其意义正如“少林武功出佛宗”。精通量价观念，进出股市将如探囊取 物，无所不能。

对于技术分析，我的见解是“钻进去，跳出来，不断实证，始终在悟字上 下功夫。”持之以恒，悟的东西多了，就可以跳出书中理论，有所创新，自成一 家了。操作很显然的，在师其法而不泥其方，灵活变化而不逾越规矩的技术分析 领域中，蔡森已成一方大家。 “外行看热闹，内行看门道”。坊间技术分析书籍林林总总，几乎都在指 标上下功夫，花俏有余实用不足，而对量价深入的探讨者寥寥可数。能以台湾地 区股市作为实例，循序渐进由浅而深，一步一步的诠释前因（量）后果（价）的

\clearpage

多空转折一手抓

技术分析书籍，蔡森算是开路先锋。

蔡森在技术分析领域里，是一位彻底的专注（Focus）者，正因为这种专 注，让蔡森有了扎实的专业发展，历经了台湾地区股市三次的万点崩盘多空大循 环，累积了20多年的实践经验与有形的财富，过着闲云野鹤的半退休生活。一 朝佛心来著，起心动念将他20多年的心血结晶集结成册，嘉惠技术分析后进， 书中对于“用”的突破，和对“表达”的突破，都是具体而微。个人乔为技术分 析领域的同好，十分认同蔡森理念、作法，将此书推荐给读者，相信此书定能为 技术分析的学习者奠立良好的基础与完整的操作逻辑。

知名专职投资人萧明道

\clearpage

\section{自序}

量价做指标靠严谨纪律拉长红

技术分析的指标繁多，自接触至今约29年来的实战操作，发现唯有量价关 系才是所有指标的最先行指标。商品的涨跌都是市场主力作量作价的结果，因此 无论底部承接布局、攻击拉拾、盘头出货、高档压低出货等，大多可以从量价结 构产生的型态来做判断，例如主力进货盘底完成后，出现突破底部整理区的攻击 量后，所有短中长期指标才会陆续出现翻多信号。既然量价结构为最先行指标， 何必浪费时间去学习其他一堆落后指标，毕竟现实上钱（量）才能真正决定商品 的价格（价），因此价格的方向是由量来决定的。除了量价，其他一切都是空 谈，懂得量价结构产生的型态，便能尽量做到先知先觉，领先市场察觉主力的动 向随势而为，也才能掌握多空方向的波段利润。出书的目的是希望以自身的经验 帮助投资人能以正确的方式尽快了解市场。当然成功没有捷径，但正确的学习方 式才能避免走冤枉路，打好基础才能发展出适合自已操作的成功模式。

这本书是本人多年来大赚小赔经验累积的操作模式，重点在从量价结构所 产生的型态来判断多空趋势：“随势操作善设停损，多头看支撑，空头看压力， 严格执行停损才能小赔出场。反之，达设定的涨跌幅满足，则波段获利了结，以 达大赚小赔的目的。”当然成功率影响操作绩效，除了学习外，经验绝对是对量 价结构了解的关键，历史是学习最好的教材，现在软件普及方便取得，多研究以 往的量价结构产生的型态后的发展有绝对的帮助，等累积一定的基础再开始操 作，实战经验一多，便能逐渐体会了解其中秘诀，当你顿悟时便会有把握，一旦 有把握，便有机会掌握时机，届时就是你大展身手之时。

\clearpage

多空转折一手抓

当你有把握，成功率便会大大提高，如一档股票底部完成出现买进信号时， 多单进场停损设约在5\%～7\%以内，有底部的型态支撑保护多头，经验上跌破支 撑停损的机率约在两至三成，相对的获利机率约有七至八成，而一个波段至少约 20\%或以上的获利机率（通常主力作价一波至少20\%），反之，头部完成时的做 空亦然，如此以较小的机率小赔，较大的机率大赚，自然长期操作下来便容易大 赚小赔。

操作不要急躁，行情是等来的，当机会来时胆大心细，如台北股市即使一 年才出现一次大多头买进时机或大空头放空时机，只要掌握住一两次就够了，就 会让你的资金快速累积，记得达预估的波段满足就好，莫再贪心才能持盈保泰， 多数人都败在贪心，以致于前功尽弃。时机来之前只需做观望等待，如此波段操 作轻松惬意。

气，长期操作短线必败，能在市场上数十年长期成功的经验都是波段操作，千万 别相信短短几年就有机会短线操作快速成功的经验，这些少数人再过几年的考验 几乎也都消失了，要记得那句常见的标语“滚石不生苔，短线不聚财”。

本书由简入繁，第一部分介绍一些常见的基本型态，并配合实际的例证， 第二部分则是实际盘势训练，看似为复杂的延伸，其实只是一般基本型态的变 换，经过第一部分的基础学习后便不难理解，对已有些程度者，更容易上手。

\clearpage

\section{前言}

关于技术分析日线、周线、月线甚至短至小时线、五分钟线选择，以及型 态的判别，尤其是量价结构，是学习的重点。有基本的认识后，当然还需要一段 时间的经验累积，从基础到进阶，除了研读历史外，还要不断累积实务经验，才 能慢慢研究出适合自已的成功模式。

长久以来，我坚守的操作方法是“严设停损，资金控管，波段随势操作”。 波段操作是以长线为基础，时间短至数月，长至数年，拉长格局来看才知道空间 有多大，迷失在短线当中如井底之蛙，则永远无法进步，跳出井口才能知道世界 有多广。

短线的格局太短，会陷入主力设局的陷阱框框中，无法了解真正的趋势， 当然容易任人宰割赚小钱赔大钱。假若时间拉长来看，主力便是在你的眼下，变 成是在你的框框当中，相对容易掌握趋势，再搭配适时的停损，犯的只会是小 错，对的时候大多是波段的大利润，便能实现大赚小赔。

既然是波段操作，当然就是尽量使用半年以上的日线或周线来做趋势上的 判断，当型态上的结构有机会到达转折时，小时线或五分钟线便是掌握先机的重 要参考指标，所以看五分钟线或小时线并不是为了要作短线，而是因为判断长线 的趋势，即将出现转折的时候，可利用较短的五分钟线或小时线，抢得先机作积 极的介入。而此时的五分钟线或小时线，也大多是拉长至两三个月来看量价变化 是否出现契机，例如K线型态上作出约三个月的底部。当来到上缘颈线量缩关前 整理判断时机已成熟时，一旦五分钟线或小时线带量突破颈线，便作积极介入同 时设好停损即可，主要看的还是波段的大利润。

\clearpage

多空转折一手抓

无论基础到进阶，历史线图的型态与量价结构所产生的变化是绝对要研读 的，尤其是初学者。这是历史演进的活教材，主力盘底进货、拉抬作价、盘头出 货凡走过必留下痕迹。每个时期（最好还能配合当时的事件）主力的战术运用， 作价模式，除了大盘外，同时期的各类股及个股琳琅满目非常丰富，虽然不外乎 进货、拉抬、出货，但过程千变万化，多练习、多了解便能增进盘感，逐渐熟 悉，至精益求精后终究能够化繁为简，进货期（买进）、拉抬期（初期加码后坐 轿）、出货期（卖出甚至翻空操作），简单的操作却是长期学习的经验及磨练才 能了解的成果。

建议一段时间的学习稍有把握后，开始以少量资金作实务上的操作累积经 验，精确的设定停损固然重要，但是能够“严格果断”执行停损又是一回事。学 习一段时间后，在有把握的一次机会中，拟定好策略积极投入资金并获得成功的 波段操作后，恭喜你的模式已进入成功的开始，这是你应得的，但往后的路还是 要持续坚持在你的成功模式中精进，甚至演化配合相关的衍生性金融商品，莫让 一次的迟疑毁了你长久的努力，严设停损是长期成功的根本，资金控管及波段操 作则是成功的体现。

1.本书的个股技术分析图均为填权息图。 2.本技术分析皆假设任何商品可以交易，且不受任何影响的情况下（如无信用 交易或融券强迫回补等）。 3.本书使用的是自创的4色K线软件，但量价结构与型态完全相同，没有影 响，读者不必在意颜色的不同。

\clearpage

\subsection{目录}

\subsubsection{技术篇}

、W底涨幅满足计算：多头买进信号 二、破底翻：多头买进信号 三、破底翻（W底）：多头买进信号 （一）2001矽统日线 （二）2008-2009爱之味日线10 （三）2008-2009新光金日线

四、波段涨幅满足计算一一下飘旗形多头买进信号…·14 （一）2001-2004中华（中华车）周线16 （二）2005-2006台股日线 （三）2008-2011佳格周线9 （四）2006-2007机电类日线 （五）2009运输类日线 （六）2012广丰周线

五、波段跌幅满足计算上飘旗形：空单进场信号24 （一）2000台苯日线26 （二）2000台达化日线 （三）2008远东新周线28

\clearpage

多空转折一手抓

（四）2011亚聚日线 （五）2010-2011台股小时趋势线30

\section{六、M头跌幅满足计算：空单进场信号32}

（一）1993-1996华纸周线 （二）1998-2000仁宝周线 （三）2006-2009中钢周线36 （四）2009-2011永光周线38

七、假突破：空单进场信号40

（一）2007-2008中纤日线 （二）2007-2008运输类指数周线 （三）2007-2008苹果日线 （四）1997日月光日线48 （五）2011正峰新日线50 （六）2011晟铭电日线 （七）2010-2011晶电日线56 （八）1993-2011微软月线58

八、头肩顶跌幅满足计算：空单进场信号60

（一）1996-1997日月光日线62 （二）2009-2011中釉周线63 （三）1995-1999营建类指数周线64

九、假突破（头肩顶）：空单进场信号66

（一）2000-2001华邦电日线68

\clearpage

（二）2008大统益日线 （三）2010-2011塑胶类日线

十、头肩底涨幅满足计算：多头买进信号74

（一）1998-2000东元日线76 （二）2008-2009蓝天日线 （三）2010扬明光日线80 （四）2011-2012广宇日线81

十一、收敛三角形头部跌幅满足计算：空单进场信号82

（一）2002仁宝日线84 （二）2007-2008台玻日线85 （三）2011扬明光日线86 （四）2009-2010所罗门日线88 （五）2011华硕日线89

十二、收敛三角形底部涨幅满足计算：多头买进信号

（一）1990-1994荣化周线· （二）2008-2009联电日线 （三）2009中航周线 （四）2011-2012新兴日线98

附注：时间波100

（一）2013华邦电日线100 （二）2011-2012大盘日线102 （三）2013法国日线103

\clearpage

多空转折一手抓

（四）2011大盘周线104 （五）2011-2012耀华日线105

\subsection{进阶篇}

一、大盘108

（一）1997-1998台股日线108 （二）1997-1999台股周线110 （三）1998台股日线112 （四）1999台股日线114 （五）1999-2001台股周线116 （六）2001台股日线118 （七）2002台股日线120 （八）2002-2004台股周线· 122 （九）2004台股日线124 （十）2005台股日线126 （十一）2005-2006台股日线128 （十二）2006台股日线130 （十三）2006-2007台股周线132 （十四）2006-2007台股日线134 （十五）2007台股日线0 136 （十六）2007台股日线2 138 （十七）2007-2008台股日线140 （十八）2008台股日线142 （十九）2008-2009台股日线· 144

\clearpage

（二十）2007-2009台股周线146 （二十一）2009-2010台股日线148 （二十二）2009-2011台股周线152 （二十三）2011-2012台股周线· 154

、个股157

（一）1997-1998华泰日线158 （二）1997-1998联电日线162

（三）1998-2000友讯周线164

（四）1998-2001茂矽周线·· 166 （五）1998-2003台积电周线168

（六）1999-2002金宝周线· 172

（七）2002-2004友达周线174 （八）2003-2008光磊周线176 （九）2003-2013宏碁日线178 （十）2004-2012台达电周线· 180

（十一）2006-2008神达日线182

（十二）2006-2009友讯周线184 （十三）2007-2008鸿准周线186

（十四）2007-2011台泥周线188

（十五）2008-2009彩晶日线190

（十六）2008-2011台化周线194 （十七）2008-2012元太日线196 （十八）2009裕日车日线200 （十九）2009-2012宏达电周线202 （二十）2010苹果日线204

\clearpage

多空转折一手抓

（二十一）2012苹果日线206

\subsubsection{附录：}

认识投机股结构2012普格日线208

压低出货陷阱1997-2001思科日线210

崛起的大国中国大陆股市212 （一）上证月线212 （二）上证周线213 （三）上证日线214

\section{四、关于A股的看法216}

\section{后记230}

\clearpage

\section{1}

\section{技术篇}

109

\subsection{76.00}

跌幅满

深森指機 成交量 价涨量增量价背离

\clearpage

多空转折一手抓

\section{一、W底涨幅满足计算：多头买进信号}

28

23

20颈线 19 18 突破

\subsubsection{双底14}

\section{13}

内容说明：

这一个多头型态，如上图，可看到当突破18颈线位置后，确立W双底的型 态，便可计算涨幅满足。技术面颈线突破，便不能再跌破颈线之下，因此突破颈 线多单进场，停损设在跌破颈线时，反之，颈线不破，看涨幅满足附近。这就是 波段操作获利大于停损风险的大赚小赔方式。

\clearpage

技术篇

计算方式：

①底部至颈线的距离=突破颈线后的等幅距离。 ②颈线中间的值约19，双底中间的值约14（中间值会有点误差但只要线画 得精准些差异不大），19-14=距离5。 ③突破颈线的位置约在18+距离5=等幅计算的涨幅满足23。

通常到达涨幅满足附近大多会进入整理，然后再看当时的情况，若多头 结构较强，则在整理过后再创波段新高时，再计算第二波涨幅满足。

④第二波的涨幅满足直接从第一波涨幅满足再加等幅距离，第一波涨幅满足 23+等幅距离5=等幅计算的第二波涨幅满足28。

操作方式：

突破颈线18双底确立多单进场，如买在18.2附近，停损设5\%～7\%以内， 也是跌落颈线之下时，颈线支撑不破，看一至两波涨幅满足附近，多单获利了 结，约26\%\textasciitilde{}53\%。

通常无论底部或头部型态，达两波段满足大多会做获利了结的动作。

\clearpage

多空转折一手抓

二、破底翻：多头买进信号

22突破

站回颈线之上 破底翻

\subsubsection{20破底}

18.5

内容说明：

如上图，股价在盘跌一段进入筑底的过程中，出现跌破底部区下缘颈线，随 后又翻升站回颈线，这是属于破底翻的结构。技术面出现破底翻时，大多是相当 难得的多头企稳买进信号。

出现这种情况大多是主力在低档筑底进货期，刻意或藉由基本面或消息面的 利空随势破底，但既然是主力进货，一旦破底，便很容易拉升站回颈线（股价经 过盘跌筑底后，筹码已清），这就是市场所谓的“甩轿”清洗浮额动作的一种。

\clearpage

技术篇

可想而知，有甩轿动作的破底翻，大多表示有主力在场，随主力进场获利的机 率，自然大幅提高。

图示经过盘跌后进入20～22整理，跌破整理区下缘颈线20后，再拉升站回 颈线20，便属破底翻。

操作方式：

0破底翻后随之买进，破底翻后的支撑仍是在整理区下缘颈线20，因此买进 后的停损价设在20附近，如果低点不远，属可承受的风险范围，则可设 在低点18.5跌破后停损，增加成功率，待突破整理区上缘颈线22做最后 加码。 ②整理区上缘突破，表示主力攻击发起，通常都是价位波动最快的时候，因 此整理区上缘22，突破后便不能再拉回跌破。 3一旦跌破22，则加码多单停损及破底翻布局的多单停利，反之，多头看支 撑，支撑不破续抱波段，如此便能低档布局，以有限的风险获取大波段利 润。

\section{MEMO}

\clearpage

多空转折一手抓

\section{三、破底翻（W底）：多头买进信号}

颈线45 突破颈线 加码信号

3232.1

\section{28.5}

内容说明：

这是另一种底部型态的破底翻W底，可藉由破底翻的结构，掌握底部 的相对低点，介入布局波段多单。

图示为W底型态，当第二只脚在30～32区间震荡时，突然出现破底动作， 随后又拉回站上30颈线之上，便属破底翻买进信号，尤其是突破起跌前高32时， 为破底翻确立，拉回皆低档布局买点。当破底翻出现后，支撑仍是颈线30附近 而非最低28.5，这是因为30才是主力成本区的下缘。操作上若买进的成本在30 附近距最低点28.5风险有限，可将停损设在破28.5时；若破底的距离太远（非图 示的28.5），便可将停损设在跌破30时（停损点太近则设在约7\%的距离，避免

\clearpage

技术篇

被洗掉），如此便能以有限的风险，获取最大的波段利润。

?操作方式：

技术面当突破颈线45时，为波段加码信号，颈线突破变支撑，表示压回跌 破颈线为停损利点（破底翻低档买进的多单获利了结，突破颈线加码的多单停 损），颈线不破，看W底等幅计算的涨幅满足附近。

\section{MEMO}

一般技术分析大多在突破颈线时，确立多头攻击信号为买进最佳时机，破 底翻则是更进阶更安全的底部布局方式，通常等待的时间较长，相对获利的空 间更大。

\clearpage

多空转折一手抓

\subsection{（一）2001矽统日线}

2363（日） 55.20商57.1.20收5.12 15956 465.460.0）甘期2001/7/11

\section{2回伊森指成交量假突破104.0 10848}

\subsubsection{104.6, 104.14}

96.29

88.43

\section{[9008] 81.8 79.18 -80.58}

79.01]

\section{73.20下飘旗形72.86}

\section{颈线59.5 68.28 -65.00}

59,22 59.72

\subsection{10/25突破颈线W底确立 -49.43 41.57}

\subsection{33.72 W底33.72}

\section{101110 75833 50555 25278 0.000}

\subsection{重点提示：W底、涨幅满足、假突破。}

\section{2001年10月25日矽统日线长红突破颈线W底确立，便能计算涨幅满足。}

\section{观察重点一：}

\section{0双底的均值约在36.7元至颈线约在59.5元，距离约22.8元。 ②突破颈线的位置约59元+22.8元=等幅计算的涨幅满足81.8元。}

\section{矽统达涨幅满足附近（最高79.18元）拉回作高档震荡整理，至突破下飘旗}

\subsubsection{形上缘颈线时，进入第二波涨势。}

\clearpage

技术篇

观察重点二：

第二波涨幅满足 第一波涨幅满足81.8元+22.8元=等幅计算的涨幅满足104.6元。 矽统达第二波涨幅满足附近（最高104.01元），便出现假突破波段卖出信号 盘头反转。

操作方式：

2001年10月25日突破颈线W底确立多单进场，约可买在59.8元（越过前高 59.72元确定突破时），停损设57元附近，也就是跌落颈线下时，风险约4.7\%， 达两波段涨幅满足，尤其是出现假突破卖出信号时，多单获利了结约60\%。

大师语录 “你的投资要集中。要是你有四十个小老婆，那么你对谁都 不会深入了解。

沃伦·巴菲特（WarrenBuffett）

明。

\clearpage

多空转折一手抓

\subsubsection{（二）2008-2009爱之味日线}

1217乘之味（日） M.57 5

\section{R 14.57}

\subsubsection{[13.90] 13.92}

\subsection{12.84}

11.76

\subsection{10.71}

\subsection{[10.45]}

9.628

\subsubsection{颈线8.551}

\subsection{7.6762 4/122}

\subsection{7.3\textasciitilde{}2009/3/30跳空7.492}

\subsection{突破颈线W底确立6.415}

\subsubsection{2009/3/5破底翻5.337}

\subsection{4.260}

\section{42829 32223 21618 11013 0.000}

\subsubsection{→重点提示：破底翻W底、涨幅满足。}

爱之味日线2009年3月5日长红站回颈线，破底翻确立，3月30日跳空涨停

\subsubsection{突破颈线，双底（W底）确立，便可计算涨幅满足。}

\subsection{观察重点一：}

\section{0W底型颈线中间位置约7.6元，至双底连线的位置约4.65元，距离2.95元。}

\section{②突破颈线的位置约7.3元+2.95元=等幅计算的涨幅满足10.25元}

\subsubsection{爱之味达涨幅满足附近（10.65元）拉回。}

10

\clearpage

技术篇

以当时爱之味约半年的底部规模来看，通常不会只涨不到一个月，便一波结 束，因此拉回不破4月1日大量低点8.01元支撑，续看第二波涨幅满足。

观察重点二：

第一波涨幅满足10.25元+2.95元=等幅计算的第二波涨幅满足13.2元 爱之味达第二波涨幅满足附近（最高13.9元），便盘头压回。

操作方式：

2009年3月5日破底翻多单进场，约可买在5.9元附近，停损可设在5.6元， 也就是破底前的前低附近，风险约5\%，3月30日突破颈线为多单加码信号，但 因跳空涨停应难介入，隔日续跳空涨停虽盘中打开，但风险已高（若盘中打开买 进多单，则停损设前一根涨停价7.49元风险约7\%）。通常不建议加码，续抱破 底翻时买入的多单即可，颈线不破，看两波段涨幅满足附近，多单获利了结约 123\%

大师语录 “很多交易员碰上商品（或股票）仓位下跌时，就采取对冲 避险来保护自己，也就是说抛空某些商品（或股票）以弥补亏损。 这实在是大错特错。

威廉·江恩（W.D.Gann）

\clearpage

多空转折一手抓

\subsubsection{（三）2008-2009新光金日线}

2888新光金（日） 7619.40低19:00收19.10 2462840720.50（-2.551\%）期2008/7714

\subsubsection{30.82}

29.19

\subsection{24.10 21.59 5/20 19.05 17.00 涨幅满足16.50}

\subsubsection{2008/9/19爆量压力区13.99}

\subsection{12.40] 11.45}

\subsection{9.56}

-8.898

\subsubsection{6.60] 7.01] 3/13破底翻-6.350}

\subsection{带量突破颈线208185262233}

\section{104092 52046 0.000}

\subsubsection{→重点提示：破底翻、涨幅满足。}

\section{新光金日线在2009在3月13日破底翻，确立首见买进信号；3月25日带量}

\section{突破颈线多头加码信号，盘涨至4月达前一年2008年9月19日爆量区（当日高点}

\section{12.3元低点11.3元）关前震荡整理约一个月。2009年5月5日突破前高12.15元，}

\subsubsection{便可计算涨幅满足。}

\subsection{观察重点：}

\subsubsection{涨幅满足计算（波段满足计算请参考P28）}

\subsubsection{1自低点6.35元涨至12.15元高点，共涨5.8元 ②自9.9元低点+5.8元=等幅计算的涨幅满足15.7元。}

12

\clearpage

技术篇

新光金5月20日爆量进入涨幅满足附近（技术面爆量还会有高点）后盘头而 下，结束一波涨幅逾倍的多头行情。

操作方式：

2009年3月13日破底翻确立多单进场，约可买在7.7元附近，停损设在7.15 元也就是跌落线下转弱时，风险约7.2\%。3月25日带量突破颈线多单加码，约 可买在9.2元，停损设8.8元，风险约5\%，达涨幅满足附近，多单获利了结约 70\%\textasciitilde{}104\%

小叮：

量价结构所产生的型态是主要观察与判断的依据，如新光金低档震荡整理约 五个月后，3月25日的带量突破颈线，便是技术面所谓的量先价行。随后便突破 整理区高点9.77元，到达2008年9月19日的爆量套牢压力区，以及大量后大多 还有高点等等。这些量价关系都是在研究历史图形时要特别观察的，有助于对未 来的判断。

\section{MEMO}

13

\clearpage

多空转折一手抓

四、波段涨幅满足计算一一下飘旗形：多头买进信号

45

\section{突破颈线 35突破前高35才能 确定30是低点也 才能计算涨幅满足}

15下飘旗形

30

20

内容说明：

首先要了解波段涨幅满足不是从图示20开始起涨后就能计算，而是做对波 段方向涨完一波回档整理后，再突破前高35，才能确定回档的低点是30，方能 计算波段涨幅满足。

如图示，自波段高点35拉回整理一段时期，回升再突破前高35，便能确定 30为回档的低点，此时才能计算波段涨幅满足。

\clearpage

技术篇

波段涨幅满足计算方法：

0自20低点涨至35高点，共涨15。 ②自30低点+15=等幅计算的涨幅满足45。

操作方法：

型态上自35高点回档整理属于下飘旗形。通常下飘旗形为多头波段的中继 站结构，当突破下飘旗形上缘颈线时，便可视为结束整理的多头再度攻击信号， 当然颈线突破就不能再跌破。因此若照图示突破颈线上缘的位置是33，约可买在 33.5，设5\%～7\%停损，也就是跌落线下时，表示当越过前高35，可以计算涨幅 满足后33便是多头支撑，也就可以设停损利点在跌破33时，操作上就是支撑33 不破，看涨幅满足45附近，获利约34\%，这便是大赚小赔的操作方式。

\section{MEMO}

多头波与空头波的整理中继站，还有矩形、三角形（正三角形、倒三角 形）等，这里仅以上下飘旗形举例说明。

15

\clearpage

多空转折一手抓

\subsubsection{（一）2001-2004中华（中华车）周线}

2204中（通）高13.65 12.50收1281 15526第

\section{保森指程成交挑量51.70 54.73}

\subsubsection{假突破51.70 51.78}

\subsection{47.89 47,19}

\subsection{42.55}

\subsubsection{38.01}

\subsection{38.17}

\subsection{33.39}

\subsubsection{24.24}

\section{17.487/21 19.62}

\subsection{13.3513.47}

\subsection{12/25逾一年的大量突破}

\subsection{12.5站上颈线确定底部完成买点10.39}

\subsection{167863}

\section{101231 67914 34598 0000}

\subsubsection{→重点提示：底形、涨幅满足。}

2001年7月21日中华（中华车）周线带量长黑破底，4周后翻红突破长黑高 点反压进入打底（长黑高点13.35元，翻红高点13.47元），同年12月25日逾一 年的大量突破站上颈线为确定底部完成的买点，技术面带量翻红前的起涨低点 12.5元为买进后的停损点。

\subsection{观察重点：}

中华（中华车）首波涨至29.37元压回整理两个多月后再突破高点时，便可 计算涨幅满足。 0自低点10.39元涨至29.37元，共涨18.98元。

16

\clearpage

技术篇

②自24.59元低点+18.98元=等幅计算的涨幅满足43.57元。

中华（中华车）达大波段涨幅满足附近后，长期大区间整理，超过一年，直 到2004年2月7日当周才带量突破整理区，6周后出现假突破的长线卖出信号， 同年4月24日当周跌破带量长红低点41.2元，翻空确立。

操作方式：

012月25日带量突破颈线买进，约可买在13.9元附近，停损设12.5元，风 险太高达10\%，可设跌破13.2元停损也就是跌落线下时风险约5\%，达大波 段涨幅满足43.64元附近，或假突破确立，多单获利了结，获利约200\%。 2上下飘旗形、距形、收敛三角形、正三角形、倒三角形、上升楔形、下降 楔形等等都是中继整理的一种，中华（中华车）周线为收敛三角形整理型 态。

收敛三角形

\clearpage

多空转折一手抓

\subsubsection{（二）2005-2006台股日线}

0000大始乘数（日） 10520.30

加1179=等幅7479.55 涨幅满足7523 7270.88

7065.68

\section{涨1179 6797,20 6747.43 4/7突破前高6857.01}

6651.82

\subsubsection{2006/3/31突破颈线28V999}

6446.63

\subsection{6373.63下飘旗形5344.73 5268.9-623796}

603276

5827.57

5618.90] 5618.90

保森指

\section{3876878 0.000}

\subsubsection{重点提示：涨幅满足、下飘旗形。}

2006年3月31日大盘日线突破颈线，下飘旗形确立。4月8日突破前高6797 点，确立6344为低点，便能计算涨幅满足。

\subsubsection{观察重点：}

0自5618低点涨至6797高点，共涨1179点。 ②自6344低点+1179点=等幅计算的波段涨幅满足7523点。

\subsubsection{操作方式：}

2006年3月31日突破颈线多单进场，约可买在6630点附近，停损设在前两 日开始带量上涨的低点6504点附近，也是跌落颈线下时，风险约2\%，达涨幅满 足附近，多单获利了结约10\%

18

\clearpage

\subsubsection{技术篇}

\subsection{（三）2008-2011佳格周线}

1227格（） 1284高13.0 8收2.31 11150 3924涨0.554.305\%

\section{特②体森指78.99668L}

78.99

\subsubsection{涨幅满足77.01}

71.01

\subsection{62.88}

54.75

\subsection{53.76}

\subsection{3/19 46.76}

\subsection{2011/4/30 38.63}

\subsection{36.35长红突破前高}

\section{40.66下飘旗形30.51}

22.52

\subsection{[13,97 -14.39}

\subsection{9.064}

\subsection{6.261}

\subsection{42936}

\section{25893 17371 8849.35 0.000}

\subsubsection{重点提示：波段涨幅满足、下飘旗形（倾向下降楔形）。}

佳格周线2008年11月创波段新低6.26元后，日线便出现破底翻买进信号。 佳格2011年3月19日当周带量突破整理区颈线，完成多头中继站的下飘旗形，4 月30日当周长红，突破前高46.92元后，便可计算大波段涨幅满足。

\subsection{观察重点：}

\subsubsection{06.26元低点涨至46.92元高点，共涨40.66元。 ②36.35元+40.66元=等幅计算的波段涨幅满足77.01元。}

\subsubsection{操作方式：}

\section{2011年3月19日当周带量突破下飘旗形颈线为多单买进时机，约可买在40.5}

\section{元附近，停损设38元，也就是跌落颈线之下时，风险约6\%，达大波段涨幅满足}

\subsubsection{附近，获利了结约90\%}

19

\clearpage

多空转折一手抓

\subsubsection{（四）2006-2007机电类日线}

1500（扫）期94.13西9493.9收组155450

\section{弹森指楼[109.71]109.71}

109.78

\subsubsection{加23.14=等幅105.31}

涨幅满足107.45

100.90

96.50

\subsection{93.56 94.11}

\subsection{91.75]92.09}

\subsection{-87.61}

\subsection{3/20突破颈线}

\subsubsection{下飘旗形8431 -83.21}

\subsection{78.80}

09/21破底翻74.40

\subsection{69.92}

文量69.92 464774

\section{278864 185910 92955}

→重点提示：破底翻、波段涨幅满足、下飘旗形。

机电类日线2006年9月21日破底翻波段买进信号，2007年3月20日突破颈

\section{线下飘旗形确立，4月2日越过前高93.06点，确立84.31点为回档低点，便能计}

\subsubsection{算波段涨幅满足。}

\subsection{观察重点：}

\subsubsection{0自69.92点低点涨至93.06点高点，共涨23.14点。}

\section{机电类指数达涨幅满足附近（最高109.71点）后大幅回跌。}

20

\clearpage

技术篇

操作方式（假设此商品可交易）：

02006年9月21日破底翻多单进场，约可买在74点附近，停损设72.3点附 近，颈线跌破时，风险约2.3\%。 22007年3月20日突破颈线下飘旗形确立，多单加码，约可买在90点，停 损设88点，也就是跌落线下时，风险约2.2\%。达涨幅满足附近，多单获 利了结约19\%（突破下飘旗形加码多单）～45\%（破底翻买单）。

小叮：

0经验上当多头强势上涨时，可以三日低点做短线强弱的判断依据（目前交 易日往前推三个交易日）做为高档减码的参考。 ②大量后虽大多有高点，但也表示风险已高，尤其是到达满足的时候。

大师语录

“我很早就发现华尔街没有什么新东西，也不可能有什么新

东西，因为股票投机历史悠久，今天在股市发生的一切都在以前 发生过，也将在未来不断地再发生。”

杰西·利佛摩尔（JesseL.Livermore）

21

\clearpage

多空转折一手抓

\subsubsection{（五）2009运输类日线}

2600输加（扫） 512毫6146454收45327274 0.39-050） E期20031

\section{26889.27}

\subsubsection{加17.88=等幅89.27}

\subsection{涨幅满足89.56 85.35 85.90}

82.53

\subsection{79.16}

\section{涨17.88 75.80}

\subsection{4/2突破颈线下飘旗形确立73.46] 72.37}

\subsubsection{71.68且突破前高便能计算涨幅满足}

69.00

\subsection{63.69 65.63}

6227

\subsection{61.09 51.10 3/4破底翻2209}

\subsection{58.84}

\section{705111 528833 352556 176278 0.000}

→重点提示：破底翻、下飘旗形、波段涨幅满足。

2009年3月4日运输类指数破底翻进入多头格局，4月2日突破颈线，下飘旗 形确立，且突破前高便能计算波段涨幅满足。

\subsection{观察重点：}

\subsubsection{0自58.84点低点涨至76.72高点，共涨17.88点。 ②自拉回低点71.68点+17.88点=等幅计算的涨幅满足89.56点。}

\subsubsection{操作方式（假设此商品可交易）：}

\section{2009年3月4日破底翻为买进信号多单进场，约可买在62点附近，停损设当 日带量翻红低点60.13点，风险约3\%，达涨幅满足，多单获利约43\%。}

\clearpage

技术篇

\subsubsection{（六）2012广丰周线}

1416量（）附1810高18251775收1亿5童53第1569肤4.3133\%日期20134 安成交

\section{指[0661] 19.92}

[19,10]

\subsection{加6.95=等幅18.30] 18.00 18.63}

涨幅满足19.15

17.29

16.85

\section{涨6.95 [15.25] 16.04}

[15.60] 14,30] 8/18突破前高14.75

\section{[3.50] 13.48}

\subsection{7/28突破颈线1218}

\section{[12.30]下飘旗形}

10.89

\subsection{9.594}

-8.300 -41510 33465 25099 16733 0.000

\subsubsection{→重点提示：波段涨幅满足、下飘旗形。}

广丰2012年7月28日突破颈线，下飘旗形确立，8月18日突破前高15.25元， 确立12.2元为低点，便可计算涨幅满足。

\subsection{观察重点：}

\subsubsection{0自8.3元低点涨至15.25元高点，共涨6.95元。 ②自回档低点12.2元+6.95元=等幅计算的波段涨幅满足19.15元。}

\subsection{操作方式：}

2012年7月28日当周突破颈线多单进场，约可买在13.5元，停损设13元，

也就是跌回线下时，风险约4\%，达涨幅满足附近，多单获利了结约42\%。

23

\clearpage

多空转折一手抓

五、波段跌幅满足计算一一上飘旗形：空单进场信号

06

\section{75}

上飘旗形 跌破颈线 跌破前低70才能 确定75是高点也 70才能计算跌幅满足

20

55

内容说明：

波段跌幅满足计算与波段涨幅满足计算为反向关系。

如图示，自波段低点70反弹整理一段时间，回跌再破前低70，便能确定75 为反弹的高点，此时才能计算跌幅满足。

7

\clearpage

技术篇

波段跌幅满足计算方法：

①自90高点跌至70低点，共跌20。 ②自75高点-20=等幅计算的跌幅满足55。

?操作方式：

型态上自70反弹的整理属上飘旗形，其为下飘旗形的反向关系。通常上飘 旗形为空头的中继站结构，当跌破上飘旗形的下缘颈线，便可视为结束整理的空 头再度攻击信号，技术面颈线跌破，就不能再站回。因此若跌破颈线的位置在73 附近，表示当跌破前低可以计算跌幅满足后73便是空头反压，空头看反压，操 作上停损利点设在站回反压73时，如空在72，停损设74也就是站回线上时，风 险约3\%，反压站不回看跌幅满足55，获利约24\%，这便是大赚小赔的操作方式。

\section{MEMO}

25

\clearpage

多空转折一手抓

\subsection{（一）2000台苯日线}

1310会等（田)

\section{45.30]假突破45.35}

\subsection{42.30 6/29 42.16}

\subsection{36.40}

\subsection{35.78}

跌18.8 9/19 32.59

破线上飘旗形确立

\subsubsection{跌破前低26.27}

\subsection{减18.8=17.6 19.50}

16.70

\section{36769 24513 2256}

\subsubsection{重点提示：假突破、波段跌幅满足、上飘旗形。}

\section{（假突破内容请见“七、假突破”）}

\section{台苯日线2000年6月29日跌落线下，假突破确立卖出信号，盘跌一波反弹}

\section{后9月19日破线上飘旗形确立，随后跌破前低26.5元，便能计算波段跌幅满足。}

\subsection{观察重点：}

0自45.3元高点跌至26.5元低点，共跌18.8元。 ②自反弹高点36.4元-18.8元=等幅计算的跌幅满足17.6元。 台苯跌至波段跌幅满足附近（最低16.7元）才企稳打底，后大幅反弹。

\subsection{操作方式：}

\section{2000年6月29日跌落线下，假突破确立，空单进场，约可空在41.8元附近，}

\section{停损设当日高点反压44元，也是站回线上时，风险约5\%。达跌幅满足空单获利}

\subsubsection{了结，约58\%。跌破上飘旗形时的空单操作亦同。}

26

\clearpage

\subsubsection{技术篇}

\subsection{（二）2000台达化日线}

1309占在（日3门位收红1

\section{22,15 22.15}

-20.15 18.10

\section{上飘旗形16.06}

\section{14.91跌10.07 14.70 7/10前低12.08确定14.02}

6/2 12.01

9.968

\subsection{926}

减10.07=等幅 跌幅满足5.11

3.841

\subsubsection{219522.714}

\section{13171 8780.74 4390.37 000}

\subsubsection{重点提示：波段跌幅满足、上飘旗形。}

\section{2000年6月2日台达化日线，破线上飘旗成形，7月10日跌破前低12.08元，}

\subsubsection{确定15.18元是高点，便能计算跌幅满足。}

\subsection{观察重点：}

\subsection{0自22.15元高点至12.08元低点，共跌10.07元。 ②自反弹高点15.18元-10.07元=等幅计算的跌幅满足5.11元。}

\section{台达化跌至跌幅满足附近（最低3.84元），因当时政治事件等连续利空，在}

\section{10至12月间短短3个月连续实施三次股市跌幅减半，但是虽有超跌，仍因达跌幅}

\subsubsection{满足后企稳打底后大幅反弹。}

\section{操作方式：}

\section{台达化6月2日破线，空单进场，约可空在13.9元附近，停损设14.5元压力}

\section{附近，也是站回颈线时，风险约4\%，达跌幅满足附近，空单获利了结约63\%。}

\clearpage

多空转折一手抓

\subsection{（三）2008远东新周线}

1402新（星） 2910高刀25.96收00.09 12383633691.11.729）日2000/2/2

\subsubsection{指费成汉直}

\section{39.1439.95}

\subsection{36.03}

-32.92

\subsection{28.73上飘旗形29.80}

跌15.12 8/16破线26.69 25.96 8/23跌破前低23.63

\subsection{18.13}

\subsection{减15.12=等幅17.40}

跌幅满足13.61 14.29 13.76]

[11.18] 11.18

\section{140937 105703 70468 35234 0.000}

\subsubsection{→重点提示：波段跌幅满足、上飘旗形。}

\section{远东新2008年8月16日当周跌破颈线，上飘旗形确立，隔周8月23日当周}

跌破前低23.97元，确立28.73元为高点，便能计算跌幅满足。

\subsection{观察重点：}

0自39.09元高点跌至23.97元低点，共跌15.12元。 ②自反弹28.73元高点-15.12元=等幅计算的波段跌幅满足13.61元。 远东新跌至跌幅满足附近（最低11.18元），才企稳打底回升。

\subsection{操作方式：}

2008年8月16日当周跌破颈线，上飘旗形确立，空单进场，约可空在25.4元 附近，停损设在长黑高点27.1元，也就是站回线上时，风险约6.6\%，达跌幅满足 附近，空单获利了结约46\%。

28

\clearpage

\subsubsection{技术篇}

\subsubsection{（四）2011亚聚日线}

1301聚（甘星） 32商30.9725.59收30.72量5506里（470.00.000\%）日期2011/11/11 爱润量

\section{2保森指39.98 40.01}

37.98

35.95

\subsection{33.48 33.92}

\subsection{32.28 31.89}

\section{跌12.84 29.90}

27.87

\subsection{25.84}

23.81

减12.84=等幅

\subsubsection{跌幅满足20.64 21.78}

\subsubsection{11/10 9157.6411799}

\subsection{4578.82 2289.41 0.000}

\subsubsection{→重点提示：跌幅满足、上飘旗形。}

2001年11月10日亚聚日线在整理区已率先出现不正常大量下跌的现象，11 月15日破线上飘旗形确立，当股价跌破前低27.14元时，便能计算跌幅满足。

\subsection{观察重点：}

039.98元高点至27.14元低点，共跌12.84元。 ②自反弹高点33.48元-12.84元=等幅计算的跌幅满足20.64元。 亚聚跌至跌幅满足附近（最低21.78元），才出现破底翻企稳。

\subsection{操作方式：}

\section{11月15日跌破颈线，上飘旗形确立，空单进场，约可空在30.6元附近，停}

损设11月10日大量高点31.89元，也是站回颈线之上，风险约4\%，达跌幅满足 附近或破底翻时，空单回补获利了结，约24\%。

29

\clearpage

多空转折一手抓

（五）2010-2011台股小时趋势线

0000大做（60分星棚） 901051

\subsubsection{9220.692/8带量不涨隔日翻黑929190}

\subsubsection{/跌回颈线之下假突破确立9220.69}

9046.18] 9093.37

\subsection{8874.50}

3/10 8838.74 8754

\subsection{8711.42}

-8581.95

\subsection{8364.95}

8327.31

\subsection{跌幅满足}

\subsection{8070.52}

8070.52

→重点提示：跌幅满足。

\section{当大盘达涨幅满足附近高风险区，此时5分钟线及60分钟的小时线，便可更}

精确的帮助判断反转时机，也就是所谓的转折点，尤其是小时线确立转折时趋势

\subsubsection{形成，因此我习惯把60分钟线称做小时趋势线。}

\section{观察上图2010-2011台股小时线，2011年2月8日台股在周线上已达涨幅满}

足附近，当天开盘第一个小时出现带量不涨，隔日开盘立即翻黑跌落原本突破的

\section{整理区上缘颈线之下，技术面假突破卖出信号确立，台股进入一波逾干点的大幅}

\section{修正。同年3月10日翻黑破线则是加码放空点，停损为当日开盘小时线，黑K高}

点8754点，站不回，看跌幅满足附近。

30

\clearpage

技术篇

这波头部区震荡的盘跌中，两个反弹高点8874点及8829点，在跌破前低后， 皆可计算跌幅满足。

观察重点：

绿色虚线1部分 0自9220高点跌至8565低点，共跌655点。 ②自8874高点-655点=等幅计算的跌幅满足8219点。

绿色虚线2部分 0自9220高点跌至8470低点，共跌750点。 ②自8829高点-750点=等幅计算的跌幅满足8079点。

台股同年3月中旬，达两个高点计算的跌幅满足附近（最低8070点）企稳反 弹，弹回前高9220点（周线第一个M头高点）附近，完成M头的第二个高点。

操作方式：

台股小时线2011年2月8日带量不涨，隔日2月9日开盘翻黑跌回颈线之下， 假突破确立，空单进场，约可空在9080点附近，停损设在压力9159点，也是站 回线上时风险约0.9\%，3月10日翻黑破线空单加码，约可空在8690点，停损设 黑K高点8754点，也就是站回线上时风险约0.74\%，达跌幅满足附近，空单获利 回补约5.4\%～11\%。（期指属高杠杆操作）

31

\clearpage

多空转折一手抓

\section{六、M头跌幅满足计算：空单进场信号}

\section{105 100}

破线

80

55

30

内容说明：

M头为W底的反向，计算的方式是相同的。 图示当跌破80颈线位置后，确立M头的型态，便可计算跌幅满足。技术面 颈线跌破，便不能再站回颈线之上，因此跌破颈线空单进场，停损设在站回颈线 时，反之，颈线站不回，看跌幅满足附近。

32

\clearpage

技术篇

\section{M头计算方式：}

①头部至颈线的距离=跌破颈线后的等幅距离。 ②头部的中间值约105，颈线假设为平行线80，105-80=距离25。 ③跌破颈线的位置80-距离25=等幅计算的跌幅满足55。

若需再计算第二波跌幅满足， ①第二波的跌幅满足直接从第一波跌幅满足再减等幅距离。 ②第一波跌幅满足55-等幅距离25=等幅计算的跌幅满足30。

操作方式：

跌破颈线80双头确立空单进场，如空在79.5附近，设5\%～7\%停损，也是 站回颈线之上时，颈线压力站不回，看一至两波跌幅满足附近，空单获利了结， 约30\%\textasciitilde{}62\%

\section{MEMO}

33

\clearpage

\section{www.guziyuan.cn}

多空转折一手抓

\subsubsection{（一）1993-1996华纸周线}

1905纸（迪莲）限10：11高10.34低9.73收.5 103802

\section{真成交32.91 34.37}

\subsection{32.97}

\section{29.96}

26.95

\subsection{23.99}

7/22 19.89] -20.98

\section{[699] 18.1 19.8420.9 17.97}

\subsection{16.64] 15.01}

13.49

\section{跌幅满足10 12.00}

\subsection{9.73 0668}

\subsection{5.98 -5.980}

\section{保森指成交量5.98 264112}

\section{158467 105645 52822 0.000}

\subsubsection{重点提示：M头、跌幅满足。}

\section{1995年华纸周线当周长黑损破颈线M头确立，便可计算跌幅满足。}

\subsection{观察重点：}

1双头中间值约29元至颈线位置约18.1元，距离10.9元。 ②跌破颈线的位置约20.9元-10.9元=等幅计算的跌幅满足10元。 华纸盘跌至隔年2月，达跌幅满足附近（最低9.73元）企稳打底回升。

\subsection{操作方式：}

\section{跌破颈线时，空单进场约可空在20.5元附近，停损设约5\%在21.5元附近，}

\section{也就是站回颈线反压时，颈线站不回，看跌幅满足附近，空单回补获利约50\%。}

\clearpage

\subsubsection{技术篇}

\subsection{（二）1998-2000仁宝周线}

2324仁营（烫层号） 12高.3 012.2收证77

\section{26.85]26.85 26.5 26.37 26.892763}

\section{M头}

24.94

22.99

\subsection{2000/8/5 21.07}

\section{17,92 19.12 19.12}

\section{[17,15] 16.80] 18.0颈线18.117.43 16.38 16.38 17.17}

\subsection{15.25}

\subsection{破底翻 [1.46] 11.35}

\section{.40 817 9.400}

\section{266068 199551 133034 66517 0000}

重点提示：M头、跌幅满足。

仁宝周线2000年8月5日跌破颈线M头确立，便能计算跌幅满足。

\subsection{?观察重点：}

\section{0M头两个高点中间值约26.5元，垂直至颈线位置约18.1元，距离8.4元。}

\section{②颈线为平行线，因此跌破颈线的位置18.1元-8.4元=等幅计算的跌幅满}

\subsubsection{足9.7元。 仁宝同年10月中旬，跌至跌幅满足附近（最低10.17元）企稳反弹。}

\subsection{操作方式：}

\section{跌破颈线时空单进场，约可空在17.5元附近，停损设在18.5元附近，也就是站}

\section{回颈线反压风险约6\%，反压站不回，看跌幅满足附近，空单回补，获利约40\%。}

35

\clearpage

多空转折一手抓

\subsubsection{（三）2006-2009中钢周线}

2002中（） 16.53高17.07 16.63收16.98 9829 21504 成交室量量

\section{33.9 34.7334.73 M头-34.7336.44}

\subsection{-32.16}

29.54

\section{26.97}

\subsection{24.35}

\subsection{-21.73}

19.16

\subsection{16.55}

\section{14.56跌幅满足13.6 [14.99] 13.93}

13.70

\subsection{11.31}

\section{荣森指標成交量1170579}

\section{705922 473593 241264 0.000}

\subsubsection{重点提示：M头、跌幅满足。}

\section{2008年8月9日中钢周线，当周长黑损破颈线，逾一年的M头确立，便可计}

\subsection{算跌幅满足。}

\section{观察重点一：}

\section{0双头中间值约33.9元至颈线位置约在26.2元，距离7.7元}

\section{②跌破颈线的位置约29元-7.7元=等幅计算的跌幅满足21.3元。}

\section{当年受2008年金融海啸影响，中钢达第一波跌幅满足仍无企稳迹象，直}

\subsubsection{接损杀两波。}

36

\clearpage

技术篇

观察重点二：

第二波跌幅满足 第一波跌幅满足21.3元-7.7元=等幅计算的跌幅满足13.6元。 中钢两波崩跌至跌幅满足附近（最低13.7元），才企稳打底回升。

操作方式：

跌破颈线时，空单进场约可空在28元附近，停损设在日线带量长黑破线高 点反压29.94元，也是站回颈线反压之时，风险约7\%，颈线反压站不回，看两大 波一次到底的跌幅满足附近，空单回补获利约50\%。

小叮：

涨跌幅满足看一波或两波，看头部或底部的规模及进行的时间做对比，也可 做判断，如中钢周线M型头部的规模达逾一年，破颈线后才约一个半月便达第 一波跌幅满足，因此可以再观察一段时间整理的结构。当然中钢直接跌两波是做 空最好的结果，若要再算第三波对稳健绩优的中钢而言，价位已属不合理，这些 都是考量的因素。

大师语录

“有钱的人可以投机；钱少的人不可以投机；根本没钱的人

必须投机。

安德烈·科斯托兰尼（AndréKostolany）

37

\clearpage

多空转折一手抓

\subsection{（四）2009-2011永光周线}

1711永光（）高

\subsubsection{34.19}

\subsection{M头32.31}

\subsection{27,68 29.42}

2/19 26.48

\subsection{21,79 23.58}

20.15] 20.64

\subsection{19.9[19.66] 17.311 17.70}

15.7]

\subsubsection{跌幅满足14.4 14.80}

\subsection{[13.13] [11.93] 11.86}

-8.915

\subsection{5.971}

森指程129194

\section{76925 52289 26628 000}

\subsubsection{重点提示：M头、跌幅满足。}

\section{永光周线2011年2月19日跌破颈线M头型态确立，便可计算跌幅满足}

\subsection{观察重点：}

\section{0M头两个高点的中间值约29.9元至颈线的位置约19.9元，距离10元。 ②跌破颈线的位置约24.4元-10元=等幅计算的等幅满足14.4元。}

\section{永光盘跌至2011年9月下旬进入跌幅满足附近，便盘出双底止跌回升。}

38

\clearpage

技术篇

操作方式：

跌破颈线后，空单进场约可空在24元附近，站上颈线停损，停损可设在25.5 元附近风险约6\%，空头看反压，反压站不回看满足，达跌幅满足附近，空单回 补，获利约40\%，大赚小赔。

小叮：

技术分析画线的时候，最常遇到的问题，就是上下影线如何取舍。记住一个 观念点多为主，越多点所连起来的一直线，精准度越高，结构也越完整。这样一 来，上下影线就自然知道如何取舍。如永光周线，M头的两个高点，取上影线， 也可以取实体线上缘，一样是两个点，虽然影响到一些满足点的幅度，但不影响 到整个结构。下缘颈线相同，取多点为主，如此一些下影线便自然舍去，因此， 可以得到一个重点，就是以多点连成的直线为主就对了。

大师语录 “投机价值的升跌，有客观的因素，是可以通过分析，计算 出来的，大家都看得中，就大家一起赚，大家都看错，只有你一 人看中，你也一样可以大赚。譬如某种股票，在没有人看好之前 你先入市，随后才有人大量入市，就会把价位堆高，你先人一步， 就可以坐享其成。

安德烈·科斯托兰尼（AndréKostolany）

39

\clearpage

多空转折一手抓

七、假突破：空单进场信号

\section{跌落颈线之下}

\subsubsection{80突破假突破确立}

70破线

内容说明：

假突破是破底翻的反向关系，股价在盘涨一段进入整理的过程中，出现突破 整理区上缘颈线，随后又跌落颈线之下是属于假突破的结构，技术面出现假突破 是转弱确立翻空的做空时机。

筹码面主力在拉高股价做高档整理震荡出货期间，刻意作价拉过新高不久， 便回落颈线之下，这就是技术面的假突破。假突破就是一种骗线，制造转强的假 象，诱使多头追价，以达其积极出货目的，随后便是将剩余持股在高档做最后的

0

\clearpage

技术篇

压低出货动作，既然是主力出货型态，当然随势做空，才会有较大的获利机会与 空间。

如图示，经过盘涨后进入70～80的高档整理，突破整理区上缘颈线80后， 再回落颈线80之下，便属假突破。操作上假突破后随势做空，假突破后80颈线 仍是反压，也是空单停损点（可稍设高些），再跌破整理区下缘颈线70，表示主 力已出货完毕或接近尾声，通常都会开始加速下跌为加码做空时机。

颈线70跌破支撑变反压，空头看反压，站回反压加码的空单停损及假突破 布空的空单停利，反之则空单续抱波段，如此便是大赚小赔，以有限的风险获取 大波段利润的操作方式。

操作方式：

跌落颈线之下，假突破确立波段多单卖出信号，空单进场，停损设5\%～ 7\%，也是站稳颈线之上时，跌破下缘颈线70空单加码，停损同样设5\%～7\%以 内，也是站回颈线之上时，颈线反压站不回，达当时空头结构计算的跌幅满足 时，空单获利了结。

\section{MEMO}

\clearpage

多空转折一手抓

\subsection{（一）2007-2008中纤日线}

1718中编（甘运）第11.2 11.55 1.01收1134量2103540.080.749\%

\subsubsection{A 15.31}

\subsection{11/8翻黑破线假突破确立 13.27}

\subsection{12.01}

\subsubsection{5/26破线假突破确立}

\subsubsection{10.78}

\section{110/8大量后9.78 9.520}

\section{大多还8.263}

\subsubsection{7.96有高点}

7.027

5.770

4.513

3.256

价涨量缩量价背离128366 102889 77411 51934 26457 0.000

\subsubsection{重点提示：量价背离、假突破、跌幅满足。}

\section{中纤2007年下半年大涨至10月8日爆量，技术面大量后大多还会有高点，}

\section{量价结构上再创新高时，已明显价涨量缩出现量价背离，11月8日翻黑破线且跌}

\section{破爆量红K低点，假突破确立，翻空盘跌。股价跌深反弹后，再破反弹低点9.78}

\subsubsection{元时，便可计算跌幅满足。}

\subsection{观察重点一：}

\subsubsection{绿色虚线部分}

\subsubsection{①自14.51元高点跌至9.78元低点，共跌4.73元 ②自11.47元高点-4.73元=等幅计算的跌幅满足6.74元。}

\clearpage

技术篇

随后盘跌过程中的小波段反弹后再破底，也可计算跌幅满足。 0自11.47元高点跌至8.27元低点，共跌3.2元。 ②自10.29元高点-3.2元=等幅计算的跌幅满足7.09元。

两个计算结果的位置都差不多。

中纤盘跌至隔年2008年1月下旬，达跌幅满足附近（最低6.93元），企 稳大幅反弹。5月26日破线且跌破5月14日长红低点假突破确立结束反弹波， 盘跌至跌破反弹波低点6.93元时，便可计算大波段跌幅满足。

观察重点二：

黑色虚线部分 0自14.51元高点跌至6.93元低点，共跌7.58元。 ②自11.13元高点-7.58元=等幅计算的大波段跌幅满足3.55元。

中纤盘跌至同年9月中旬达大波段跌幅满足附近（最低3.26元），企稳打底 约半年后大幅回升。

操作方式：

02007年11月8日翻黑跌破颈线空单进场，约可空在12.2元附近，突破颈线 也是长黑高点12.9元停损，风险约5\%，颈线站不回，看跌幅满足附近， 空单回补，获利约40\%。 22008年5月26日假突破确立空单进场，约可空在9.7元附近，站上突破颈 线前的整理区高点10.54元反压停损，风险约8\%，反压站不回，看大波段 跌幅满足，空单回补，获利约60\%。

\clearpage

多空转折一手抓

\subsubsection{（二）2007-2008运输类指数周线}

\section{-87.12 146.00}

\section{净森指荐145.62] 146 146.00 5/24跌落颈线下146.18152.89}

\section{M头假突破确立卖出讯号}

135.30

124.60

\subsubsection{7/19跌破颈线M头确立113.90}

\subsection{10320}

\subsection{100.5100.08 104.5}

92.32

\subsection{81.62}

73.04

\subsection{70.92}

60.22

\subsection{49.34}

\subsubsection{→重点提示：M头、假突破、跌幅满足。}

\section{运输类指数周线2008年5月24日，当周跌落前一周突破的颈线之下，假突}

\section{破确立为卖出信号，7月19日当周长黑攒破颈线，M头确立，便能计算跌幅满 足。}

\subsection{观察重点：}

0双头的高点146点至颈线的位置约在100.5点，距离45.5点。 ②跌破颈线的位置约在104.5点-45.5点=等幅计算的跌幅满足59点。

\section{运输类指数跌至跌幅满足附近，便进入盘底后回升。}

\clearpage

技术篇

操作方式（假设此商品可交易）：

02008年5月24日当周跌落线下，假突破确立空单进场，约可空在136点， 停损设在反压区142点附近，也就是站回颈线之上时，风险约4.5\%。 27月19日当周跌破颈线，空单加码，约可空在104点附近，停损点设在当 周长黑高点107.17点，风险约3\%，达跌幅满足附近，空单获利了结，约 43\%（加码空单）～56\%（假突破空单）。

小叮：

0运输类指数周线达第一波段跌幅满足后，若再计算一波当然是极不合理的 位置，因此只看一波。 ②第二波跌幅满足59点-45.5点=13.5点，这对类股指数来说，几乎是不可 能的点，除非是一只炒过头的投机股。

大师语录 “凡是证券交易所里人尽皆知的事，不会令我激动。经常有 人问我，我的信息和想法是从哪来的。我的回答很简单，其实， 我不寻找信息，而是发现信息。”

安德烈·科斯托兰尼（AndréKostolany）

\clearpage

多空转折一手抓

（三）2007-2008苹果日线

196假突破203.1 带量长黑190.4

177.8

\subsubsection{颈线165.3}

\subsection{15063 156 1:58 152.7}

\section{140.0}

\subsection{114.9 102.3 89.60}

价涨量缩12044

\subsubsection{量价背离60603}

0.000

→重点提示：假突破、量价背离。

假突破属于技术面卖出信号，尤其是股价处于高档出现量价配合的假突破， 更是强烈的卖出信号，这是主力高档整理作突破的动作吸引买盘抢进后，开始进 行压低出货，这时便会产生假突破现象。既然是主力出货的行为便属于卖出信 号，通常出现假突破后，大多是非盘即跌，盘跌走空的机率较高，基本面较强 的，则先修正再经过一段整理期后，才有以盘代跌机会扭转弱势型态。

苹果股价2007年一波大多头涨幅惊人，2007年12月24日股价在高档整理后 再度突破创波段新高，不过很明显这次的盘涨突破属价涨量缩量、价背离，初露 败象。

六个交易日后也就是2008年1月4日出现带量长黑，跌回整理区，此时便能 判断属高档假突破，确立为强烈的波段卖出信号，苹果股价随后跌破颈线头部，

10

\clearpage

技术篇

确立大幅回跌至波段起涨点附近，自高点大跌修正幅度高达约43\%。

操作方式：

苹果2008年1月4日带量长黑，假突破确立，空单进场，约可空在189美元， 停损设198美元，也就是站回线上时，风险约5\%，回补点可以192美元及202美 元两个高点形成的双头（M头）计算跌幅满足，约在118美元附近，也是前波起 涨附近的较大支撑区，达跌幅满足，空单获利回补约37\%。

小叮：

苹果股价虽然每次的大幅拉回后都能再创新高，但并不是每次都能那么幸 运。产业总有循环，从技术面判断每个大波段的买卖点，除了能掌握买点外，还 能避开产业自高峰反转的风险。

大师语录 “投机就是投机，不是赌博，这才是正确的投机观念。投机 可以经分析做基础，制订计划，选择项目，这是一整套理性的过 程。这个道理说得简单，却不是人人都明白。没有树立正确投资 观念的人，他的投机可以说就是在碰运气，在赌博，将自己的人 生交给了所谓的运气。成功的投机者会将赌博和投资分得清清楚 楚。

安德烈·科斯托兰尼（AndreKostolany）

\clearpage

多空转折一手抓

\subsubsection{（四）1997日月光日线}

23月1日月光（日运） M9.55高10.03 9.45日期1997/522

\subsubsection{成交21.01}

\subsubsection{20.30假突破20.33}

8/21 9/2 19.09 18.30

\subsection{17.86}

\subsubsection{10/17}

16.65

\subsection{9/916.041}

\subsection{15.41}

\subsection{12.97}

\subsubsection{11.73}

[10.90]10.50 9.266 100480 80384 60288 40192 20096 0.000

\subsection{重点提示：高档头部区假突破的判断。}

\section{通常从量价结构的变化，大多可以判断是否为高档头部区，仔细观察便可抓}

\subsubsection{到转弱的时机。}

\section{如图示，日月光头肩顶的结构，1997年8月21日带量红K突破左肩整理颈}

\section{线，突破整理区的带量红K，表示是主力作价的攻击信号，这根红K低点便不会}

\section{被跌破，一旦被跌破就表示主力拉高的动作是假的，也就是在做高档拉高出货的}

\section{假突破，这是一种骗线动作。因此9月2日低点17.3元跌破红K低点17.86元，便}

\section{能判断这是假突破的骗线动作，骗线就是为了要出货，既然是主力高档出货，当}

\subsubsection{然是随势做空，才容易有波段获利空间。}

48

\clearpage

技术篇

观察重点：

0技术面9月2日判断为假突破确立后，任何反弹皆是空点，高点20.3元为 反压就是空单停损点（高点若太高为了掌控风险，停损可设约5\%～7\% 以内），9月9日再破低点后，8月21日至9月8日便是假突破的区块。 210月17日跌破颈线为技术面最佳加码放空点，通常假突破高档布空利润 较大，但等待期较长，因为主力仍需要时间盘头，将剩余筹码在高档尽量 出干净（如图示，头部完成也确立左肩后还要盘右肩）。跌破颈线表示主 力已出货完毕，此时做空通常下跌的速度最快，跌破颈线后表示颈线以 上，已成为主力出货给散户的压力区就不会再站上，因此跌破颈线空单停 损便设在站回颈线时，颈线站不回则看跌幅满足附近，这便是大赚小赔的 波段操作方式。

操作方式：

9月2日假突破确立后，任何反弹皆空点，若空在假突破前的高点压力区 18.3元附近设高点反压20.3元，风险约11\%太高，可直接设一根停板的7\%，便可 有效控制风险，达P62的“1996-1997日月光日线”图虚线所提跌幅满足附近空 单回补，获利约39\%。

大师语录

\section{“不要忽视警讯——大的亏损很少没有警讯，不要在正在下}

沉的船中等待救生艇。

安德烈·科斯托兰尼（AndréKostolany）

6

\clearpage

多空转折一手抓

\subsubsection{（五）2011正峰新日线}

1538正峰新（日） 0高27.00 24.5日期2010/1119 统交

\section{4/8带量翻黑58.00] 6.10550也是假突破58.00}

\section{假突破确立53.39}

\subsection{6/2 48.79}

47.22

\subsection{44.18}

\subsection{39.58}

\subsection{33.40 33.50 -34.97}

\subsection{30.37}

\subsection{25.76}

\subsubsection{价跌券增主力放空21.16}

16:55

量1585

\subsection{3903.62}

\subsubsection{0.00035814}

\subsubsection{融券空头高档30058}

\subsubsection{强迫回补认输24488}

\section{主力颈线区快速出货18733-14237}

-15000 -15763 -16526 17289

\subsubsection{→重点提示：假突破、主力放空。}

\section{如图示正峰新日线，2011年第二季转亏为盈，主力掌握小型股筹码易掌控的}

\section{投机特性轧空，2011年4月8日轧至新高随后出现带量翻黑，假突破确立的卖出}

\subsection{信号。}

\section{同年6月2日带量长红隔日高点56.5元突破整理区高点56.1元，但3日内便}

\section{跌破长红低点也是属于假突破，股价在主力做最后骗线后随之崩盘。}

50

\clearpage

技术篇

?观察重点一：

散户放空结构 1价涨券增，通常是散户往上放空的结构，容易引起主力轧空。 ②如上图，正峰新日线2011年3月起价涨券增，往上加码放空的融券散户在 主力有计划轧空下，藉由股东会强迫在高档认输回补，切记，放空的技巧 是空弱不空强，会创新高的股票就是强势股，做错就认输停损，不应加码 放空或建立强势股空单，需待有转弱现象出现时才是做空时机（如假突破 确立）。 3正峰新同年6月下旬跌至颈线支撑附近，库存余额急降，显示主力利用颈 线支撑积极作量倒货，长期的库存余额显示主力长期成本低于20元，主 力利用自56.5元高点，跌落至40元以下，吸引1跌深抢进买盘在33元至40 元颈线区做高档压低出货动作（虚线部分），由于主力长期的成本低随便 卖都赚，洞悉筹码结构，便能避开主力在高档区颈线压低出货陷阱。

观察重点二：

主力放空结构 延续上述主力出货后随后一个多月，价跌券增属主力放空结构，9至10月的 反弹最高33.5元，仅碰触上述33至40元颈线区下缘属弱势反弹（支撑跌破变反 压），期间无量，便是主力出货后呈无主力状态，弱势结构可见主力持空信心坚 强，随后股价长期大跌两大波段。

51

\clearpage

多空转折一手抓

\subsubsection{让我们来看看正峰新的周线。}

153正新（） 22高2070倍18猫收19.20 1.0（5133\%）日

\subsubsection{-6075}

\subsubsection{涨幅满足0089 58.00}

\subsection{47,30 52.36}

\subsection{42.70] 6/25 46.71}

41.07

\section{35.42}

3/12 29.78 24.74 21,61] 24.85 24.13 18.49

\subsection{[14.73] 1284}

\subsection{42800 28690 14580 0.000 24968 14655 3997.00}

\subsection{主力高档开始不计成本出货6240.505214.00}

-7267.00 -8293:50

\subsubsection{20139320.00}

\subsection{重点提示：涨幅满足。}

\section{正峰新周线在2011年3月轧空突破前高47.3元，便可计算涨幅满足。}

\subsection{观察重点：}

\subsection{0自14.73元起涨低点至47.3元高点，共涨32.57元}

\section{②自24.85元低点+32.57元=等幅计算的涨幅满足57.42元附近。}

\section{技术面3月12日周线带量长红，突破颈线的颈线支撑38元附近不破，看涨}

\subsubsection{幅满足57.42元附近。}

\section{正峰新达涨幅满足附近后（最高58元），便开始盘头，同年6月25日周线}

52

\clearpage

技术篇

带量翻黑破线，头部确立，就此进入长空，库存余额亦显示主力自涨幅满足高 档，开始长期压低出货。

基本面季报仅当年第二季转亏为盈配合轧空后，又再度陷入长期亏损，负债 比逾50\%，每股净值也跌破10元。投资市场常常是人为的投机，有迹可循，并非 巧合。

操作方式：

0正峰新周线2011年3月12日当周带量突破颈线买单进场，约可买在39元 附近，停损设37元，也就是跌落线下时风险约5\%，达涨幅满足附近，多 单获利了结约47\% ②正峰新日线2011年6月2日长红低点3日后被跌破时空单进场，约可空在 52.8元，停损设高点56.5元，风险约7\% 可以头部方式计算跌幅满足， ·两个高点中间值约57.5元，至颈线约在34.7元距离约22.8元 ·7月13日长黑跌破颈线的位置约在36.5元-22.8元=等幅计算的跌幅 满足13.7元 ③达跌幅满足附近空单获利了结约74\%（正峰新低点随净值跌破票面价值的 10元以下）。

53

\clearpage

多空转折一手抓

\subsubsection{（六）2011晟铭电日线}

3013最花（日） 0 5453.60联租2011/320 EM-

\section{66.62969}

\section{[59.30] 61.20 61.50 6/13跌破长红低点-61.84}

\subsection{[55.20] 56.60] 6/1 57.07}

\subsection{最低59.9 48 7/11跌破前低41.9元52.30}

便可计算波段跌幅满足47.53

42.76

\subsection{37.98}

31.00] 33.21 28.44

跌幅满足

\subsection{16876 11057 5619.18}

\subsubsection{0.00039306}

\section{33445 30467 2587.50 31.50 -2524.50 5080.50 -7743.00}

\subsubsection{重点提示：假突破。}

如图示，晟铭电日线，2011年3月末升段喷出后，股价波段累积大涨逾

200\%，同年4月、5月两个月份库存余额，显示主力拉高后高档震荡出货。

\subsection{观察重点：}

\section{2011年6月1日带量长红突破整理区，随后便拉回跌破带量长红低点，假突}

破确立，由此可知假突破经常是主力骗线的一种手法，利用一段时间的整理，让 市场习惯股价后，再作量转强让市场尤其是散户跟进。如晟铭电当时突破后，累 积涨幅已达约300\%，因此出现假突破后表示主力拉高骗线出货确立，便是风险 最高的时候，也是喜好做空者最佳放空时机，此类股切忌跌深抢反弹，晟铭电最 高约65元（还原权息最高66.5元），大跌逾20元至40余元，只做弱势整理，便

\clearpage

技术篇

又数波大跌（同年12月9日创18.2元最低价）。

2011年基本面季报获利维持在小亏小赢之间，（全年EPS-0.38元）更显股 价炒作至逾50元的投机性，可见当投机过后，都是股价漫漫长夜的开始。

?操作方式：

0做空的技巧就在股价转弱时。如上图示，2011年晟铭电日线6月13日跌破 6月1日带量长红低点59.9元假突破，主力出货，确立败象已露（当日收 58元）。 26月14日及15日仍有两日平盘上布空机会（两日高点皆58.8元），停损设 在突破颈线约61.5元时风险有限，约6\%。 3需注意6月30日除权息融券强迫回补后，7月6日恢复融券将回补的筹码 再空即可。 47月11日跌破反弹低点41.9元，便可计算跌幅满足。 自66.5元高点至41.9元低点，共跌24.6元。 自48元高点-24.6元=等幅跌幅满足23.4元（通常需填权息计算则在23.8 元附近） ③股价8月下旬达跌幅满足附近整理，便是空单波段获利时，约58\%。

大师语录 “我花了整整五年的时间才觉得自己能理智地玩炒股游戏。

杰西·利佛摩尔（JesseL.Livermore）

55

\clearpage

多空转折一手抓

\subsubsection{（七）2010-2011晶电日线}

2448品（日道福） 94.41.18

\subsection{107假突破107.4211171}

\subsection{99.93 29.93带量翻黑破线头部确立100.69}

93.96

\subsection{87.35}

\section{84.18 2011/7/8 80.82}

跌破前低后

\subsection{计算跌幅满足73.89}

70.64

\subsection{63.58 6728}

60.55

\subsection{跌幅满足53.82}

47.09] 47.09

\subsection{48753}

\subsubsection{大量后不涨39377}

\section{19876 10126 0.000}

\subsubsection{重点提示：假突破、跌幅满足。}

晶电2011年3月8日带量突破颈线，三个交易日内便跌破大量低点，随后更 是跌回颈线以下，确立大量不涨的假突破卖出信号，4月19日带量翻黑破线头部 确立，空头成形后股价一路盘跌。

\subsubsection{观察重点：}

07月8日跌破反弹低点75.16元破底后，便可计算跌幅满足。 自107.3元高点至75.16元低点，共跌32.14元。 自88.24元高点-32.14元=等幅计算的跌幅满足56.1元。

56

\clearpage

\subsubsection{技术篇}

28月2日跌破前低65.15元后，也可计算跌幅满足。 自88.24元高点至65.15元低点，共跌23.09元。 自70.64元高点-23.09元=等幅计算的跌幅满足47.55元。

晶电8月达跌幅满足附近企稳筑底反弹（最低47.09元），累积反弹幅度约 七成，1年多后又再度破底。

\subsection{操作方式：}

\section{晶电2011年4月19日带量翻黑破线头部确立空单进场，约可空在95.5元停}

损设99.5元，也就是站回线上时，风险约4.2\%，7月8日破线跌破前低空单加码，

约可空在75元，停损设79元，也就是站回线上时，风险约5.3\%，达跌幅满足， 空单获利了结约25\%～50\% 同样的，让我们来看看晶电的周线。

241品电通71

\section{06T110.54 111.9011843}

107.30

\section{98.24翻黑跌破大量低点102.10}

假突破确立长线卖出信号

92.13

\subsection{假突破结束反弹82.17}

\subsection{[70.84] 66.80 72.37}

\subsection{5435 52.44}

47.09] 42.64 [40.00] -32.68 22.71 146495 117196 87897 58598 29299 0.000

\clearpage

多空转折一手抓

\subsubsection{→重点提示：假突破。}

晶电周线跌破大量低点假突破确立长线卖出信号的位置，便是对应上述日线 同期假突破的位置，也可明显看出晶电2009年的明星产业题材大涨将近四倍后， 处于高档风险区。

\subsubsection{观察重点：}

假突破后约2年的大头也逐渐成形，随后因产业跨入门槛低竞争激烈而进

入获利快速衰退的空头结构，量价结构反应基本面的情况，107.3元高点大跌至 47.09元后的反弹，也在假突破后结束再陷长空格局，高点与低点皆越来越低反 映了产业窘境。

\subsubsection{（八）1993-2011微软月线}

FUSOO微歌（月运福） M5.49高5.83 4.99收5.78量2056748日期1996/1/31 時間爱费成交张肤星量德量

\section{2回染指標59.97 1999/12 60.13}

\subsection{假突破53.59}

\subsection{2007/10}

\subsection{372假突破40.82}

\subsubsection{BS'TE 32.95 34.44}

27.89

\subsection{20.70 21.46] 21.51}

\subsection{14.872009/3 15.13}

\subsubsection{破底翻8.743}

5.870 2.200 3079927 2042126 1037802 0.000 1994 1996/1/319971998199920002001200220032004200520062007 2008 2009 2010 20112012 2013

58

\clearpage

技术篇

→重点提示：假突破、破底翻。

微软月线自1993年8月低点2.2美元至1999年12月高点59.97美元，累计 涨约26倍，6年多的大多头行情在2000年2月跌回突破的颈线以下，确立了长线 卖出信号“假突破”，股价大幅回跌至2000年12月低点20.12美元，才出现反弹 波。

观察重点：

股价长期横盘至2007年10月才做突破颈线动作，但四个月后又跌回突破的 颈线以下再次出现假突破卖出信号，随后便一路盘跌走空至2009年3月最低点 14.87美元，累计跌幅达约60\%股价，超过腰斩。

股价见低点三个月后，同年2009年6月翻回颈线之上破底翻确立，也确立了 长空数年以来的企稳信号，随后3年多来股价维持在颈线支撑21美元附近以上长 期整理。

操作方式：

微软1999年以前，可说是科技类当红炸子鸡，只是再好的产业都逃不过景 气循环，时间长短的问题而已，长线技术面假突破及破底翻是较精确掌握景气是 否反转的判断方式，当出现假突破信号时长线便应见好就收，等待类似破底翻的 企稳信号，才做介入设好停损等待。

59

\clearpage

多空转折一手抓

八、头肩顶跌幅满足计算：空单进场信号

头部 90

左肩右肩

70破线

50

30

内容说明：

如同W底与M头的反向关系，头肩顶是头肩底的反向关系。 如图示我们可以观察到， ①头部高点90至颈线位置70（颈线假设为平行线）距离20。 ②跌破平行的颈线位置70-20=等幅计算的跌幅满足50。 ③若需再计算第二波跌幅满足，

第一波跌幅满足点50-20=等幅计算的跌幅满足30。

60

\clearpage

\section{技术篇1}

操作方式： 跌破颈线时，空单进场，如空在69附近，停损设5\%～7\%以内，也是站 回颈线之上时，反压站不回，看一至两波的跌幅满足附近，空单获利了结，约 27\%\textasciitilde{}56\%

\section{MEMO}

\clearpage

多空转折一手抓

\subsubsection{（一）1996-1997日月光日线}

2311日月光（日） 19

\section{2要文21.43}

\subsection{20.34}

18.47

16.60

\section{14.73}

\subsubsection{10/17 12.85}

带量翻黑破线：

\subsubsection{头肩顶确立LLluau.11.02}

10.90] 9.146

7.276

\subsection{3.535}

\section{60755 40504 20252}

\subsubsection{重点提示：头肩顶、跌幅满足。}

\section{日月光日线1997年10月17日带量翻黑，跌破颈线头肩顶确立，头肩顶的左}

\section{肩、头部、右肩型态相当完整，破线当日，便可计算跌幅满足。}

\subsubsection{观察重点：}

0自头部高点20.3元垂直至颈线位置约在15元附近，距离约5.3元。 ②跌破颈线位置约在16.4元附近-5.3元=等幅计算的跌幅满足11.1元。 日月光同年10月30日达跌幅满足附近（最低10.9元），V型反转打底后回升。

\subsection{操作方式：}

\section{跌破颈线时，空单进场约可空在16元附近，停损设16.8元，也就是站上颈 线反压风险约5\%，反压站不回，看跌幅满足附近，空单回补获利约30\%。}

62

\clearpage

\subsubsection{技术篇}

\subsubsection{（二）2009-2011中釉周线}

1109中种（） 215商2322导2 文

\subsection{30.84头肩顶30.8432.62}

-27.93 26.03

24.98

6/18 22.02 [9861

\subsection{20.25 19.07}

16.16

\subsection{13.20}

跌幅满足10.56

\section{10.25}

7.293

\subsection{4.337}

\section{70928 53196 35464 17732 0.000}

\subsubsection{重点提示：头肩顶、跌幅满足。}

\section{如图示，2011年6月18日中釉周线，当周跌破颈线，约一年两个月的头肩顶}

确立，便可计算跌幅满足。

\subsubsection{观察重点：}

0头肩顶高点30.84元至颈线20.7元，距离10.14元。 ②跌破颈线的位置20.7元（颈线为平行线）-10.14元=等幅计算的跌幅满足 10.56元。 中釉盘跌至同年底达跌幅满足附近（最低10.37元）企稳反弹。

\subsubsection{操作方式：}

跌破颈线时，进场布空单，约可空在20.5元附近，停损设在2周前起跌高点 21.7元，风险约6\%，反压站不回，达跌幅满足附近，空单回补获利约48\%。

63

\clearpage

多空转折一手抓

（三）1995-1999营建类指数周线

2500管建领（运）

\subsubsection{2指595.21595 594.05}

逾二年头部595.98 545,33 549.94

\subsection{涨幅满足}

504.66

8/1 459.39 429.61

\subsection{429上飘旗形413.35}

40758颈线

\subsection{390 368.07}

\subsection{346.57 322.80}

\section{颈线277,13 31129}

4/6 276.76

\subsubsection{约一年底部跌幅满足224 231.48}

216.84

\subsection{185.44}

文量1076400

\subsubsection{→重点提示：合式头肩顶。}

营建类指数周线1996年4月6日当周突破颈线，约一年的底部确立（多单进 场跌破颈线停损），大涨一波拉回整理后再过前高时，便可计算涨幅满足。

\subsection{观察重点一：}

0自216.84低点涨至396.83点，共涨179.99点。 ②自352.8低点+179.99点=等幅计算的涨幅满足532.79点。

\subsubsection{观察重点二：}

营建类指数涨至隔年1997年3月，达波段涨幅满足附近，进入大区间长期

04

\clearpage

技术篇

盘头，隔年1998年8月1日当周跌破颈线，逾两年的头部确立，便可计算跌幅满 足。（头部有两个高点双头，称为复合式头肩顶。） ①自头部两个高点的中间值约595点至颈线位置约390点，距离205点。 ②自跌破颈线的位置约429点-205点=等幅计算的跌幅满足224点。

操作方式（假设此商品可交易）：

跌破颈线时，空单进场，约可空在420点附近，停损设在440点附近，也就 是站上颈线反压，风险约5\%，反压站不回，看跌幅满足附近，空单回补获利约 46\%

1.技术面结构上头部的规模远大于底部的规模，因此容易跌破底部低点。 2.跌破颈线头部完成后的反弹形成上飘旗形（空头中继站），反弹未过颈线， 便是技术面所谓的逃命线。

65

\clearpage

多空转折一手抓

九、假突破（头肩顶）：空单进场信号

\section{跌落颈线下 突破假突破确立卖出信号 90}

\subsubsection{颈线70跌破颈线}

\section{68空单加码信号}

内容说明：

这是一种头肩顶的假突破结构，通常学习技术分析者大多会在跌破颈线确定 头部后，才做放空动作，但是利用假突破的结构，更可进阶判断股价高档转弱而 在相对高档布空。

操作方式：

如图示，高档整理后突破颈线90，随后又回落颈线之下，假突破确立。假 突破是一种主力拉高出货的骗线动作，因此随势做空，停损设在颈线反压区90 附近（或稍高些），待跌破颈线72确定头部后为空单加码信号，空头看反压，

66

\clearpage

技术篇

此时反压自90降至头肩顶的颈线72附近，颈线跌破就不能再站回，因此站回颈 线，加码的空单停损，假突破高档做空的空单停利，反之颈线未站回，看波段跌 幅满足，波段操作小赔大赚。

\section{MEMO}

\clearpage

\subsubsection{多空转折一手抓}

\subsection{（—）2000-2001华邦电日线}

2344琴邦（日） 1:04高3.45收2.25 57063 12518张肤0.0010.985\%期2000/6/13

\subsubsection{特保森指视型德量}

\section{56/14突破后隔日跌落-85.46}

\section{颈线下假突破确立}

78.47

\section{左肩右肩-71.37}

64.70

\subsection{977跳空跌破颈线64.2661.32}

\subsection{头肩顶确立57.15}

\section{56.13]颈线56.5 56.17}

\subsection{49.33上飘旗形50.17}

-43.06

\subsection{32.82 36.67 35.96}

\subsection{28.85}

\subsection{21.91.破底翻-21.75}

\section{188235 141176 94118 47059 0.000}

\subsubsection{→重点提示：假突破、头肩顶、跌幅满足。}

\section{2001年9月7日华邦电日线跳空跌破颈线头肩顶确立，便能计算跌幅满足。}

\subsection{观察重点：}

\section{0头部高点85.46元至颈线约在56.5元，距离约28.96元。}

\section{②跌破颈线的位置56.5元（颈线为平行线）-28.96元=等幅计算的跌幅满足}

\subsubsection{27.54元。}

\section{华邦电跌至跌幅满足附近大幅震荡（最高32.82元，最低21.75元）后，出现}

\subsubsection{破底翻，确立打底反弹。}

68

\clearpage

技术篇

操作方式：

02001年6月14日突破颈线创新高，隔日立即跌落颈线以下，假突破确立头 部空单进场，约可空在82.5元附近，停损设高点85.46元，也是站回线上 时，风险3.6\% 29月7日跌破颈线空单加码，约可空在55元，停损利设在回补缺口，也是 站回颈线时的58.6元，风险约6.5\%，达跌幅满足附近，空单获利了结约 50\%（加码空单）～66\%（假突破空单）。

小叮：

跳空跌破颈线追空的短线风险较高，若无把握，可先建立空单部位，收盘确 定跌破，隔日再做较有把握的加码动作，设好停损点即可。

华那电日线若再计算第二波跌幅满足，则数字为负值，因此不做计算，且8 月出现破底翻的企稳信号，空单只做一波。

大师语录 “谁缺乏耐心，就不要靠近证券市场。

安德烈·科斯托兰尼（AndréKostolany）

69

\clearpage

多空转折一手抓

\subsubsection{（二）2008大统益日线}

1232大益（日运楼）高25.27.5日期200873711 要量2.88

\subsection{假突破头部2. 5/27跌落颈线下32.91}

\subsubsection{31.28假突破确立卖出信号}

\subsubsection{头肩顶30.60 30.87}

\subsection{27.72右肩28.83}

\subsubsection{左肩26.52 2720 26.82}

\subsubsection{距离9.38 8/12 25.23}

247

\subsubsection{24.024.35}

\subsection{颈线23.5 22.73}

\subsection{20.72}

\subsection{19.33}

\subsection{18.368.36 减9.38=等幅计算16.63 的跌幅满足14.97 16.34 14.59}

成交量举森指4030.00

\subsection{2418.00 1612.00 806.00}

\subsubsection{重点提示：假突破、头肩顶、跌幅满足。}

\section{2008年大统益日线5月27日跌落突破的颈线下，假突破确立卖出信号，8月 12日跳空跌破颈线，便可计算跌幅满足。}

\subsection{观察重点：}

\section{0头部高点32.88元至颈线的位置约在23.5元，距离9.38元。 ②跌破颈线的位置约在24.35元-9.38元=等幅计算的跌幅满足14.97元。}

\section{大统益跌至跌幅满足附近（最低14.59元），才出现破底翻企稳。}

70

\clearpage

技术篇

操作方式：

0假突破确立卖出信号空单进场，约可空在29.6元附近，停损若设32.88元 高点，风险超过10\%，太高，可设约7\%，约在31.6元附近。 28月12日跌破颈线头肩顶确立，空单加码约可空在23.9元附近，停损设 25.5元，风险约6.5\%，达波段跌幅满足附近，空单获利了结约37\%（加码 空单）～49\%（假突破空单）。

小叮：

大统益日线属于非常完整的头肩顶型态，通常这类型态大多会有两波段跌幅 满足，但若再计算第二波跌幅满足仅剩5.59元，这严重背离了该股的基本面，再 者，约五个月的头部，跌破颈线后，盘跌的时间也达约两个半月，时间已足够， 所以这是技术面只做一波的主要考量。

大师语录 “如果操作过量，即使对市场判断正确，仍会一败涂地。

乔治·索罗斯（GeorgeSoros）

\clearpage

多空转折一手抓

\subsubsection{（三）2010-2011塑胶类日线}

1300里领（日运得） 30815高308.15297.77收304.94日期2011/3/3

\section{保森指假突破头部322.12322.12 32767}

\section{头肩顶2011/5/4跌落颈线下假突破322.33}

确立卖出信号 305.46] 304,62 309.99

\subsubsection{左肩一右肩297.86}

\subsubsection{距离285.73}

57.12

\subsection{8/2 273.60}

[271.15] 264.03

\section{颈线265 266.35268.5沙一逃命线261.27}

\subsection{251,77}

\subsection{249.14}

237.01

\subsection{减57.12=等幅227.15 224.88}

\section{218.13跌幅满足211.38 221.33 21255}

\section{1262054 946541 631027 315514 0000}

重点提示：头肩顶、假突破、跌幅满足。

塑胶类日线2011年5月4日跌落颈线下，假突破确立卖出信号，8月2日跳空 跌破颈线头肩顶确立，便能计算跌幅满足。

\subsubsection{观察重点：}

0头部高点322.12点至颈线位置约265点，距离57.12点。 ②跌破颈线的位置约268.5点-57.12点=等幅计算的跌幅满足211.38点。

塑胶类指数跌至跌幅满足附近（最低212.55点），才企稳打底反弹。

\clearpage

技术篇

操作方式（假设此商品可交易）：

02011年5月4日跌落颈线下，且在跌破4月21日带量长红突破颈线的低点 （309.54点）时，确定假突破进场做空，约可空在309点附近，停损设高 点322.12点，风险约4\% 28月2日跳空跌破颈线，头肩顶确立空单加码，约可空在266点附近，停损 设前一天，也就是补空点282点，风险约6\%，达跌幅满足附近，空单获利 了结，约18\%（加码空单）～30\%（假突破空单）。

小叮：

①跌破颈线头部完成的时候，经常会有回测颈线的动作，这个动作就叫做 “逃命线”，如塑胶类日线跌破颈线头肩顶确立后，自227.15低点拉高回 测颈线的高点264.03点，这个动作就是技术面所提的逃命线。 ②当类股指数转弱确立时，可以这类族群中的弱势股布空为主。

大师语录 “摊平是一种可能造成严重亏损的业余策略。

威廉·欧尼尔（WilliamO'Neil）

73

\clearpage

多空转折一手抓

十、头肩底涨幅满足计算：多头买进信号

110

92

\section{78}

\section{74 突破}

60

内容说明：

如图示，我们可以观察到， 0低点到颈线的距离=突破颈线后的距离。 ②头肩底低点60至颈线位置78距离=18。 ③突破颈线的位置74+18=等幅计算的涨幅满足92。 ④若需再计算第二波涨幅满足，第一波涨幅满足点92+18=等幅计算的涨幅 满足110。

\clearpage

技术篇

操作方式：

突破颈线时多单进场，如买在75附近停损设5\%～7\%，也是跌破颈线时， 颈线不破，达一至两波涨幅满足时，多单获利了结，约22\%～46\%。

\section{MEMO}

\clearpage

多空转折一手抓

\subsection{（一）1998-2000东元日线}

\subsection{1504東元（日蛋）特查15.27成交商15.44 15.16量收15.34址4806日斯199803/11}

\section{体森指楼26.13}

\subsection{26.16}

\subsection{21.80] 22.42 21.74}

\subsection{20.54}

18.67

\section{16.03颈线17.57 [17.17] 16.83}

\subsection{[5,29]}

\section{14.77] 14.77 14.48 14.21/5突破颈线14.96}

\subsection{11.21}

\subsubsection{左肩2/22}

\subsection{-9.340}

\section{109891 82418 54945 27473 0.000}

\subsubsection{重点提示：头肩底、破底翻、涨幅满足。}

\section{东元日线1999年2月22日，破底翻底部确立，2000年1月5日带量跳空突破}

\subsubsection{颈线，头肩底确立，便能计算涨幅满足。}

\section{观察重点一：}

\subsubsection{0底部9.34元至颈线约在15.3元，距离5.96元。 ②突破颈线的位置约在14.2元+5.96元=等幅计算的涨幅满足20.16元。}

\section{东元达第一波涨幅满足附近，仅做短线震荡，便持续盘涨，以东元长达逾一}

\section{年半以上的底部规模仅涨半个多月，便达首波涨幅满足，因此，可续看第二波涨}

\subsection{幅满足。}

76

\clearpage

技术篇

观察重点二：

第二波涨幅满足计算 第一波涨幅满足20.16元+5.96元=等幅计算的第二波涨幅满足26.12元。 东元达第二波涨幅满足附近（最高26.13元），便大幅拉回。

操作方式：

01999年2月22日破底翻，多单进场，约可买在10.9元附近，停损若设9.34 元低点，风险超过10\%太高，可设在10.2元附近，也就是再度破线确立 时，风险6.5\%。 ②2000年1月5日带量跳空，突破颈线，头肩底确立，多单加码，约可买在 14.6元附近，停损设跌破14.1元，也就是补空跌落线下时，风险约3.4\%， 达两波段涨幅满足，多单获利了结约78\%（加码多单）～140\%（破底翻 买单）。

大师语录 “我更关心的是，怎么控制失利状况。要学会接受亏损。想 赚钱，最重要的就是不能让亏损失控。

马蒂·史华兹（MartySchwartz）

\clearpage

多空转折一手抓

\subsubsection{（二）2008-2009蓝天日线}

2362整天（日限） 25.18 25:72 754 8.52.104\%斯2

\section{2微量42.528/4 44.64}

\subsubsection{假突破4252}

\subsection{39.35}

36.12

\subsection{32.90}

\subsection{29.67}

\subsection{27.21 26.50}

\subsection{23.27}

\subsection{-20.04}

12/19带量突破颈线头肩底确立

16.82

\section{左肩底部右肩13.59}

\section{17622 13216 8810.98 4405.49 0.000}

\subsubsection{重点提示：头肩底、涨幅满足、假突破。}

\section{2008年12月19日蓝天带量突破颈线，头肩底确立，便能计算涨幅满足。}

\subsection{观察重点：}

\section{0底部低点13.59元至颈线约在20.8元，距离约7.21元。 ②突破颈线的位置约在20元+7.21元=等幅计算的涨幅满足27.21元。}

\section{蓝天达涨幅满足附近（最高29.43元）快速急回2日，便持续盘涨。}

\subsubsection{第二波涨幅满足}

\section{③第一波涨幅满足27.21元+7.21元=等幅计算的第二波涨幅满足34.42元}

78

\clearpage

技术篇

蓝天达第二波涨幅满足短线整理后，虽超涨仍进入大幅震荡（最高40.52 元，最低32.68元）

第三波涨幅满足 ④第二波涨幅满足34.42元+7.21元=等幅计算的第三波涨幅满足41.63元

蓝天达第三波涨幅满足附近后（最高42.52元），才出现假突破，进入长期 震荡整理。

操作方式：

2008年12月19日带量突破颈线多单进场，约可买在20.4元附近（突破关前 整理区的近期高点20.34元，确定突破时），停损设当日带量长红低点19.26元， 也是跌落线下时，风险约5.6\%。通常达两波涨幅满足，便获利了结约68\%，若 达第三波涨幅满足后，2009年8月4日出现假突破，确定卖出信号，获利了结约 96\%

大师语录 “我有认错的勇气。当我一发现犯错，马上改正，这对我的 事业十分有帮助。

乔治·索罗斯（GeorgeSoros）

\clearpage

多空转折一手抓

（三）2010扬明光日线

\subsection{假突破183.22}

\subsection{涨幅满足[171,25] 169.43 175.70}

\subsection{168.18}

160.67

\subsection{53.49 153.15}

\subsubsection{.10142.1颈线12/27跌落线下145.76}

13824

右肩130.72

左肩11/15破底翻123.20

头肩底底部116.60) 115.68

7017.98

1754.50

重点提示：破底翻、头肩底、涨幅满足。 扬明光日线2010年11月3日破底，11月15日拉回破底的颈线之上，破底翻 确立，12月3日突破颈线头肩底确立，便可计算涨幅满足。 观察重点：

0颈线142.1元（颈线为平行线）至底部低点115.68元，距离26.42元。 ②突破颈线的位置142.1元+26.42元=等幅计算的涨幅满足168.52元。 扬明光喷出急涨3日，便达涨幅满足附近（最高171.25元）拉回，维持高档震 荡整理，12月22日长红突破再创新高后第三个交易日12月27日，便跌落线下，假 突破确立，也确立扬明光技术面只涨一波，盘头的规模与盘底的规模相当。 操作方式：

2010年11月15日破底翻确立，买进多单，约可买在135元附近，停损设支 撑127元附近，风险约6\%。12月3日突破颈线多单加码，约可买在143元附近， 停损设137元，也是跌落线下时风险约4\%，达涨幅满足附近后的12月27日跌落 线下，假突破确立，无第二波行情时，多单获利了结约17\%（加码多单）～25\% （破底翻多单）。

80

\clearpage

\subsection{技术篇}

\subsection{（四）2011-2012广宇日线}

232学付湿）高2 25.41收25.46

\subsubsection{322假突破32.28}

\section{2TE}

\subsection{30.94}

\section{[TO'OC 29.42 29.42 29.61}

\subsection{26.94}

\subsection{25.2 25.62}

\subsection{24.29}

1/30带量长红22.95

突破颈线头肩底确立

\section{左肩2右肩-21.62}

12/23破底翻底部确立20.28

\section{9925.15 7443.86 4962.57}

\subsubsection{2481.29477.17}

→重点提示：破底翻、头肩底、涨幅满足、假突破。 2012年1月30日广宇日线带量跳空长红突破颈线，便可计算涨幅满足。 观察重点：

\subsection{0底部低点20.28元至颈线约在25.2元，距离约4.92元。 ②突破颈线的位置约在24.5元+4.92元=等幅计算的涨幅满足29.42元。}

\section{广宇达涨幅满足附近（最高32.26元）震荡整理，而高档震荡续创新高后，立}

\section{即回跌出现假突破，且破线后整理的时间过长与底部的规模相当，因此只看一波。 操作方式：}

\section{02011年12月23日破底翻底部确立（突破长黑高点）多单进场，约可买在}

\section{22.6元附近，停损设21.5元，也是再度跌落线下时风险约5\%。}

\section{②2012年1月30日突破颈线头肩底确立多单加码，约可买在24.7元附近（越}

\section{过前高24.63元确定突破时）停损设前一日收盘24.14元，也是跌落线下时}

\section{风险约4\%，达涨幅满足附近至假突破后破线（大约是涨幅满足附近），}

\section{多单获利了结约19\%（加码多单）～30\%（破底翻多单）。}

81

\clearpage

多空转折一手抓

十一、收敛三角形头部跌幅满足计算：空单进场信号

80

65破线三角形需在1/2\textasciitilde{}3/4 处跌破才算有效跌破 60

45

25

内容说明：

三角形头部为三角形底部的反向结构。收敛三角形头部的颈线跌破是一种做 空信号。

观察上图，我们可以发现收敛三角形需在1/2～3/4处跌破才算有效跌破， 整理至越尾端的跌破，越容易产生失败，出现反方向的走势。

82

\clearpage

技术篇

如图示，跌破收敛三角形下缘颈线后，便能计算跌幅满足。

三角形计算方式： ①边长=跌破后的等长。 ②以高点80做计算至60的边长=20。 ③破线处65-20=等幅计算的跌幅满足45。 ④若需再计算第二波跌幅满足：第一波跌幅满足45-20=等幅计算的第二波 跌幅满足25。

操作方式：

跌破收敛三角形下缘，头部确立空单进场约可空在64，停损设5\%～7\%， 也是站回颈线之上时，颈线反压站不回，则达一至两波跌幅满足附近，空单获利 了结29\%\textasciitilde{}60\%

\section{MEMO}

83

\clearpage

多空转折一手抓

\subsection{（一）2002仁宝日线}

2324（日梁）高150 14.76 爱费

\subsection{1746}

\subsection{116.97}

\subsubsection{6.174/24 16.58}

\subsection{15.69]3/22[15.56 [15.65] 15.69}

15.19

\section{[14.34] 14.81}

\subsection{[12.97] 13.05}

\subsection{12:16}

12.13

11.28

\subsection{10.73 10.39}

\subsection{9.508}

\section{54298 40724 27149 13575 0.000}

\subsubsection{→重点提示：收敛三角形头部、跌幅满足、异常量。}

\section{2002年5月2日仁宝日线跌破颈线，收敛三角形头部确立，便可计算跌幅满}

\section{足。注意3月22日及4月24日出现的大量，价未过前高，但却出现突发性最大量}

\section{属异常量。在5月2日破线后这两日的异常量确立，也更确立头部的结构。}

\section{观察重点：}

\section{0三角形边长高点17.29元至低点约在12.7元，距离约4.59元。 ②跌破颈线的位置约在14.9元-4.59元=等幅满足计算的跌幅满足10.31元。}

\subsection{操作方式：}

\section{2002年5月2日带量长黑跌破颈线，空单进场，约可空在14.5元附近，停损}

\section{设在当日长黑高点15.08元，也是站上颈线时，风险约4\%，达跌幅满足附近，空}

\subsubsection{单获利了结约29\%}

84

\clearpage

\subsubsection{技术篇}

\subsubsection{（二）2007-2008台玻日线}

1802台级（日） 25.57收然7 量

\section{原通区32.96 3322}

31.96

\subsubsection{29,12 30.49}

\subsection{6/19跌破颈线 收敛三角形头部确立25.45 24.92}

\subsection{23.7 22.93}

\subsection{21.32 20.46}

20.1

\subsection{8 16.9017.94}

15.42

\subsection{12.90}

[0.38] 10.38 21995 17732 13299 8866.20 4433.10

重点提示：收敛三角形头部、跌幅满足。

\section{2008年6月19日台玻日线跌破颈线，三角形头部确立便能计算跌幅满足。}

\subsection{观察重点：}

0三角形边长高点31.96元至低点约在20.1元，距离约11.86元。 ②跌破颈线的位置约23.7元-11.86元=等幅计算的跌幅满足11.84元。 台玻跌至跌幅满足附近（最低10.38元）企稳回升。

\subsection{操作方式：}

2008年6月19日跌破颈线，收敛三角形头部确立，空单进场，约可空在23.3

元附近，停损设前一日高点24.13元，也是站回颈线时，风险约3.6\%，达跌幅满 足附近，空单获利了结约49\%。

85

\clearpage

多空转折一手抓

\subsubsection{（三）2011扬明光日线}

3504光（日级） 135.72132日期2011/6/13 爱

\section{159.41] 159.62}

\subsection{6/15带量破底收敛三角形确立}

\subsubsection{[27,86] 135.26}

\section{132.8上飘旗形123.08}

109.13

\subsubsection{110.90}

02-001 98.93

\subsection{三角形}

\subsection{77.98 74.57}

6134

\subsection{50.21}

\subsection{5004.90 3753.67 2502.45 1251.22 0.000}

\subsubsection{重点提示：收敛三角形头部、跌幅满足。}

扬明光日线2011年6月15日带量破线后，收敛三角形确立，便能计算跌幅 满足。这波跌势当中的前两波皆无出现底部结束空头的结构，两次的上飘旗形加 上一次弱势整理的三角形，皆是往下破线的空头中继结构，直到第三波才破底翻 打出头肩底回升，因此总共计算了三波跌幅满足。

\subsection{观察重点：}

第一波跌幅满足 0边长的高点159.41元至低点颈线约在132元，距离27.41元 ②跌破颈线的位置约在132.8元-27.41元=等幅计算的跌幅满足105.39元。

86

\clearpage

技术篇

扬明光跌至跌幅满足附近（最低100.7元）反弹。

第二波跌幅满足 ③第一波跌幅满足105.39元-27.41元=等幅计算的第二波跌幅满足77.98元

扬明光跌至跌幅满足附近（最低74元），才又一次反弹，但仅是弱势反 弹，形成三角形的弱势整理结构，往下跌破线后续算第三波跌幅满足。

第三波跌幅满足 4第二波跌幅满足77.98元-27.41元=等幅计算的第三波跌幅满足50.57元

扬明光这次跌至跌幅满足附近（最低50.21元）才出现破底翻，结束长空型 态，打底大幅做倍数以上的弹升。

操作方式：

2011年6月15日带量破线，空单进场，约可空在131元附近，停损设137.5 元，也是站回颈线时风险约5\%，达两至三波跌幅满足时，空单获利了结约40\% （第二波跌幅满足）～61\%（第三波跌幅满足）。（上飘旗形及三角形整理的颈 线跌破，也是放空时机）

大师语录 “永远不要孤注一掷。

乔治·索罗斯（GeorgeSoros）

87

\clearpage

多空转折一手抓

\subsubsection{（四）2009-2010所罗门日线}

\section{29（日） S012高旺M201}

\subsection{16.62}

\section{6.13] 15.83 16.13}

15.47

1/21跌破颈线

\subsection{15.25收敛三角形头部确立14.89}

\section{Z71.614.7 4.7] 1431}

[13.93] 13.75

\subsection{[3.49] 13.34 13.29 13.17}

12.59

破底翻12.01

\subsection{11.43}

\section{10003 7502.51 5001.67 2500.84}

\subsection{重点提示：收敛三角形头部、跌幅满足。}

\section{2010年1月21日所罗门日线跌破颈线，收敛三角形头部确立，便能计算跌幅}

\subsection{满足。}

\subsection{观察重点：}

\section{0三角形边长高点16.61元至低点约在14.7元，距离约1.91元。 ②跌破颈线的位置约15.25元-1.91元=等幅计算的跌幅满足13.34元。}

\subsubsection{③第二波跌幅满足}

\section{第一波跌幅满足13.34元-1.91元=等幅计算的第二波跌幅满足11.43元。}

\section{所罗门刚好达第二波跌幅满足点后，破底翻筑底反弹。}

\subsection{操作方式：}

\section{2010年1月21日跌破颈线，收敛三角形头部确立，空单进场，约可空在14.9}

\section{元，停损设在站回颈线时的15.4元，风险约3.4\%，达两波段跌幅满足附近，空单}

\subsubsection{获利了结约23\%}

88

\clearpage

\subsection{技术篇}

\section{（五）2011华硕日线}

（日）5艺1122第.061959收1570

\section{森指231.67 231.81}

\subsection{223.33}

9/23跌破颈线

\subsubsection{收敛三角形头部确立214.85}

\section{204,25] 200.63 206.36}

\subsection{199.32 197.88}

\subsubsection{186.211184.46] 189.54}

\section{184.06] 189 181.06}

180.30

174.84] 172.58

\section{165.5破底翻164.09}

\section{16228 12171 8114.02 4057.01}

\subsection{重点提示：收敛三角形头部、跌幅满足、破底翻。}

\section{2011年9月23日华硕日线，跳空翻黑跌破颈线，收敛三角形头部确立，便能}

\subsection{计算跌幅满足。}

\subsection{观察重点：}

\section{0三角形边长高点231元至低点约在189元，距离约42元。 ②跌破颈线的位置约在207.5元-42元=等幅计算的跌幅满足165.5元。}

\section{华硕跌至跌幅满足附近（最低155.61元），出现破底翻的企稳信号。}

\section{操作方式：}

\section{2011年9月23日跌破颈线，收敛三角形头部确立，空单进场，约可空在}

\section{204.5元附近（跌破前低204.57元破线确立），停损设211.5元，也是回补跳空缺}

\section{口站回颈线时，风险约3.4\%，达跌幅满足或破底翻时，空单获利了结约19\%，且}

\subsubsection{破底翻应翻空为多。}

89

\clearpage

多空转折一手抓

十二、收敛三角形底部涨幅满足计算：多头买进信号

80

65

55

\subsection{50突破三角形需在1/2\textasciitilde{}3/4}

处突破才算有效突破

40

内容说明：

前面提过，收敛三角形需在1/2～3/4处突破才算有效突破，整理至越尾端 的突破，越容易产生失败，出现反方向的走势。

如图示，突破收敛三角形上缘颈线后，便能计算涨幅满足。

90

\clearpage

技术篇

三角形计算方式：

①边长=突破后的等长。 ②三角形边长以低点40做计算至55边长=15。 ③突破处50+15=等幅计算的涨幅满足65。 若需再计算第二波涨幅满足：第一波涨幅满足65+15=等幅计算的第二波 涨幅满足80。

操作方式：

突破收敛三角形上缘颈线，底部确立多单进场，停损设5\%～7\%以内，也 是跌落线下时，颈线支撑不破，则看一至两波涨幅满足附近，多单获利了结，约 30\%\textasciitilde{}60\%

\section{MEMO}

91

\clearpage

多空转折一手抓

\subsection{（一）1990-1994荣化周线}

1704乘北（通蛋）

\subsection{14.62}

\subsubsection{[3.31 13.36}

涨幅满足13.03 13.06

12.09

\subsection{10.81}

9.65]

\subsection{9.533}

\subsection{8.278}

7.001

5.723

12/4带量突破颈线4.446

\subsection{3.169}

\section{106418 79813 53209 26604}

\subsection{→重点提示：收敛三角形、涨幅满足。}

\section{荣化周线1993年12月4日带量突破颈线，收敛三角形底部确立，便可计算}

\subsection{涨幅满足。}

\subsection{观察重点：}

\subsubsection{1三角形的边长=突破颈线后的等长。}

\section{②低点3.17元至三角形垂直边长高点约在10.7元，距离7.53元}

\section{③突破颈线的位置约在5.5元+7.53元=等幅计算的涨幅满足13.03元。}

\section{荣化盘涨至1994年8月20日当周，进入涨幅满足附近，便开始大区间盘头}

\subsection{后大幅回跌。}

92

\clearpage

技术篇

操作方式：

带量突破颈线时多单进场，约可买到5.6元附近，跌破颈线设停损，可设破 5.3元后，风险不到6\%，多头看支撑，支撑不破，看满足，达涨幅满足附近获利 约120\%，大赚小赔。

小叮：

0三角形除了收敛三角外，还有正三角形与倒三角形一，通常正三角形 为头部型态，倒三角形为底部型态。 2荣化周线虽有正三角形的倾向，但12月4日的带量突破是往上突破，属量 先价行的结构，量价结构配合，多头就能够顺利进行。

大师语录

“坚持自己的想法，别人的想法只是别人的，你所能用的只

有自己的节奏，所以要不随波逐流。这样你才能在如流的人海中 找到一块真正属于自己的领土。

安德烈·科斯托兰尼（AndréKostolany）

93

\clearpage

多空转折一手抓

\subsubsection{（二）2008-2009联电日线}

20电（日）限8.03 51.73

量12.46

\subsection{12.46 1248}

\section{11.70}

10.92

\subsection{10.16}

9.714

\subsection{9.379}

边长距离约2.59 下降楔形8.601

\subsection{8.42]}

\subsection{-下飘旗形7.837}

\section{7.26}

\subsection{[6.57 6.75] 3/9长红突破颈线7.059}

\subsubsection{收敛三角形底部确立6281}

\subsection{5.503}

\subsubsection{护保指楼5.35 5.505.50}

\section{360760 270570 180380 90190 0.000}

\subsubsection{重点提示：收敛三角形底部、涨幅满足、下飘旗形。}

\section{2009年3月9日联电日线长红突破颈线，收敛三角形确立，便可计算涨幅满}

\subsubsection{足。}

\section{观察重点：}

\section{0三角形边长高点7.94元至低点约在5.35元，距离约2.59元。 ②突破颈线的位置约7.26元+2.59元=等幅计算的涨幅满足9.85元。}

\section{联电达第一波涨幅满足附近（最高10.42元），便拉回整理，完成下降楔}

\subsubsection{形越过前高后，便可再计算第二波涨幅满足。}

6

\clearpage

技术篇

3第二波涨幅满足计算 第一波涨幅满足9.85元+2.59元=等幅计算的第二波涨幅满足12.44元。

联电达第二波涨幅满足附近（最高12.46元），便盘头回跌。

操作方式：

2009年3月9日长红突破颈线，收敛三角形底部确立，买单进场，约可买在 7.4元附近，停损设当日长红低点6.94元，也是跌落线下时，风险约6\%，达两波 段涨幅满足附近获利了结约68\%。

小叮：

3月9日突破的当时，或许没有把握是否为有效突破（有时候会留上影线收 盘没过），此时除了看五分钟线或小时线是否带量突破，也可以前高7.39元为依 据，突破便确立立刻买进，所以在操作方式会买在7.4元附近，停损也设好风险 便已锁定。

大师语录

“不顾市场情况，每天以感情冲动进进出出，是华尔街很多

炒手亏钱的主要原因。他们试图像做其他工作一样，每天都能拿 一笔钱回家。世界上没有比亏钱更好的老师。当你学习怎么做才 不会亏钱时，你开始学习怎么赚钱了。

杰西·利佛摩尔（JesseL.Livermore）

95

\clearpage

多空转折一手抓

\subsection{（三）2009中航周线}

2612中航（蛋） 931 16100

\subsubsection{77,.44 77.44}

\subsection{59.73 -72.05}

66.57

\subsection{61.37}

\subsection{[59.31]}

\subsection{52.92 55.61}

50.22

\subsection{43.89 41.9 44.74}

2/7突破颈线三角形收敛底部确立39.26

33.78

\subsection{28.30 2830}

\section{26363 19772 13181 6590.70 0.000}

重点提示：收敛三角形底部、涨幅满足。

\section{中航周线2009年2月7日当周突破颈线，收敛三角形底部确立，便能计算涨 幅满足。}

\subsubsection{观察重点：}

第一波涨幅满足 0三角形边长高点43.31元至低点约在26.8元，距离约16.51元 ②突破颈线的位置约在41.9元+16.51元=等幅计算的涨幅满足58.41元。

中航涨至涨幅满足附近（最高58.89元）周线出现长上影线，仅做短线大 幅震荡，便再度盘涨，无反转迹象，续看第二波涨幅满足。

96

\clearpage

技术篇

③第二波涨幅满足 第一波涨幅满足58.41元+16.51元=等幅计算的第二波涨幅满足74.92元。 中航达第二波涨幅满足附近（最高77.44元），进入长期盘头后反转。

操作方式：

中航2009年2月7日当周突破颈线，收敛三角形底部确立，买单进场约可买 在42元附近，停损设当周带量长红低点40.88元，也是跌落线下时风险约3\%，达 两波段涨幅满足，多单获利了结约78\%

小叮：

中航周线的收敛三角形倾向倒三角形，倒三角形大多属于底部的结构，因为 上缘为平行，突破后至少近期的上档无压，多头容易进行，尤其是跌深后的底部 型态颈线是平行的，突破后大多有相当大的涨幅。

大师语录 “买股票时，需要想像力；卖股票时，需要理智。

安德烈·科斯托兰尼（AndréKostolany）

97

\clearpage

多空转折一手抓

\subsubsection{（四）2011-2012新兴日线}

（印）509 23.34高23.4低2307收23.29 305

\subsubsection{2森指翼受真成交}

\subsection{27.61 27.29 27.62}

\subsection{26.98}

\subsubsection{加3.3=27.1 26.99}

26.36

边长距离约3.3 26.13

\subsection{25.90]破线确立头部25.72}

\subsection{25.09}

25.05

24.47

\subsection{1/31突破颈线}

收敛三角形底部确立23.84

23.21

\subsection{22.44}

\subsection{22.1 21.94}

\subsection{6409.30 4806.98 3204.65 1602.33 0.000}

\subsubsection{重点提示：收敛三角形底部、涨幅满足。}

\section{新兴日线2012年1月31日突破颈线，收敛三角形底部确立，便能计算涨幅 满足。}

\subsection{观察重点：}

\section{0三角形边长25.4元高点减低点约22.1元，距离约3.3元。 ②突破颈线的位置约在23.8元+3.3元=等幅计算的涨幅满足27.1元。}

\section{新兴涨至涨幅满足附近（最高27.61元）进入整理，由于整理时间过长未过 前高且破线形成头部，因此仅能计算一波涨幅满足。}

98

\clearpage

技术篇

操作方式：

2012年1月31日突破颈线，收敛三角形确立，多单进场，约可买在23.9元附 近，停损设23.4元，也就是跌落线下时，风险约2\%，达涨幅满足附近，多单获 利了结约13\%

小叮：

新兴日线收敛三角形上线颈线偏斜，突破后仍会有压，尤其边长高点，也就 是大量高点25.4元为实质压，加上先前的盘跌层层压力，属于非大跌后的底部， 相对涨幅便有限，也就是经过大跌后的收敛三角形，才容易有相对大的涨幅。

大师语录 “长期来看，证券市场无法脱离经济。投机者必须仔细观察 分析本国的经济形势，当然还必须观察分析世界的经济形势。但 要注意，具决定性的不是过去的发展，而是未来的发展。

安德烈·科斯托兰尼（AndréKostolany）

99

\clearpage

多空转折一手抓

附注：时间波

内容说明：

时间波的计算与看法，我不记得曾在任何技术分析上看到相关的理论，这是 我长期研究观察型态以来，认为对行情判断有相当大的助益，尤其是当时间到达 对称时经常会出现变盘，时机非常巧合。主要原因是当结构转好或坏后，大多需 要与转好或坏前，也就是左边（前面）整理区大约相同的整理时间来消化卖压， 这时便会产生左右对称的结果。

计算方法：

时间波的计算要在相对高度的整理区。 ①若是转空型态（如假突破），对称的时间内通常是非盘即跌。在对称的时 间内随时都会跌的状况下，拖过对称的时间过久，才有扭转空头的机会。 ②若是翻多型态（如破底翻），对称的时间内通常是非盘即涨。在对称的时 间内随时都会涨，拖过对称的时间过久，才有再度转弱的机会。

熟悉后便能事先布局，再配合技术分析，在转折点适时加码。

举例说明：

（一）2013华邦电日线

2013年6月7日华邦电跌落线下，假突破确立，此时便能计算出突破前的左 边整理区为26天，就表示假突破后右边对称的26天里，属于非盘即跌的风险期， 也就是大约26天或以内破线下跌的机率相当高，除非拖过26天后才有以盘代跌 的机会（用时间换取空间），结果华邦电在26天内的第二十天翻黑破线完成头

100

\clearpage

\subsection{技术篇}

\section{肩顶，直到达头肩顶计算的等幅跌幅满足6.56元附近才企稳（最低6.57元），因}

\subsubsection{此在对称的风险期间做空为主，设好停损即可。}

2344禁邦（日道础） M6.11 6:34 5.05收.15 3233里4583源10门.653\%）日期2013/3711

\subsubsection{更贝10.10}

\section{假突破10.10 2013/6/7跌落线下10.11}

\section{假突破确立}

\subsection{26天9.39 26天9.600}

\subsection{8.95第二十天翻黑破线9.171}

\subsubsection{头肩顶完成}

8.592

\section{8.10颈线8.25] 8.3 18.3231]8.36 [7.93] 8.084}

\subsection{7.26 7.575}

\subsection{7.075}

\section{6.34跌幅满足6.56 6.567}

6.058

-5.550

\section{森指标.55成交量124646}

\section{74788 -49858 24929 0.000}

\subsection{MEMO}

\subsubsection{101}

\clearpage

多空转折一手抓

\subsubsection{（二）2011-2012大盘日线}

0000大临指数（日）随7101.20毫7109.347055 回森指8170.728170.31 8170.72

\subsubsection{3/29带量破线7997.88}

\subsection{7825.00}

\subsection{对称第二十天变盘7652.14}

\subsection{7479.287529.08}

\subsubsection{第15天7349,85}

\subsubsection{19天突破颈线7303.48}

\subsection{19天第二十天翻红起涨}

713062

\subsection{12/21}

\subsection{-6784.90}

\subsubsection{破底翻6609.11}

1693726

大盘日线2011年12月21日长红站回颈线，破底翻确立，也确立终结当时欧

\section{债风暴的空头结构翻空为多，此时便可计算出左边整理区共19天，表示右边也}

应有对称约19天非盘即涨的多头安全期，也就是说破底翻后的低点6609点不应 该会再破的，破了就是判断错误多单停损的时候，结果台股在安全期内的第十五

\section{天突破颈线确立底部，更在对称整理的19天后，也就是第二十天，带量翻红，}

准时变盘，持续滚量上攻大涨。

操作上破底翻时多单进场，停损设前波整理区低点6744点，风险仅约3\%， 突破颈线底部确立后，按照头肩底方式计算两波段涨幅满足，或3月29日带量破 线转弱确立，波段获利了结即可。

102

\clearpage

\subsubsection{技术篇}

\section{这波底部完成后，上攻的多头波在3月29日带量翻黑破线转弱确立后，4月}

\section{3日刚好就是左右两边对称的时间波，第二十天准时一日不差地翻黑跌破颈线，}

\section{完成M头型态，向下变盘，操作上3月29日及4月3日翻黑破线皆为放空时机，}

\section{站回线上停损，风险皆在2\%以内，获利点可按照M头等幅计算两波段跌幅满}

\subsubsection{足，或7422点反弹后再破底时，计算波段跌幅满足即可。}

\subsubsection{（三）2013法国日线}

\section{P004法巴（日通） 8415万9西4200.3141665收4193.17}

\subsection{彩指特货0.00更0.00成交0.000.000绝量0}

\section{11/17带量上影线假突破4336.61 17天4357.11}

\section{17天4309.151\_4309.92 4311.81 4309.20}

\subsection{4228.00}

\subsection{颈线42254211.82 4222.31 4214.22}

415971

\subsection{4118.42}

\subsection{4105.67跌幅满足4094 4071.34}

\subsubsection{405125] 4023.44}

3975.53

\subsection{3927.63}

交量1991.00 1592.80 1194.60 796.40 398.20 0.000

\section{法国股市日线2013年11月17日带量上影线假突破，左边的整理区共17天，}

\section{表示有对称17天非盘即跌的风险，右边也整理17天后隔日变盘，开盘先跌破6}

\section{日低点小颈线，再破颈线头肩顶确立，达头肩顶型态计算的跌幅满足4094点附}

\subsubsection{近后，才企稳拉高。}

\subsubsection{103}

\clearpage

多空转折一手抓

\subsection{（四）2011大盘周线}

0000大微（） 7074.39高7139.046548.57收7120.51日期2012/177

\subsubsection{2回第森指9220.699220.69}

\subsubsection{9099.759089.472011/6/10空袭警报9220.69}

长空起始 8931.60

8642.52

\section{19周19周}

\section{-835343}

\section{8064.34}

8070.52

\subsection{7770.36}

\subsection{7481.27}

\subsection{7162.78}

\subsection{6903.10}

6877.12

\subsubsection{6609.11] 6609.11}

\subsection{84342519}

67474016

\section{2011年大盘周线M头型态，蓝色虚线为技术分析的M头颈线，但是在计算}

\section{时间波是以相对高度的整理区作对称来计算（深蓝色实线），左边的头整理时间}

\section{共19周，右边也整理19周后隔周变盘，注意这左边19周整理后无法立即计算右}

\section{边的对称19周，而是右边配合2011年6月10日判断转空后，在整理接近19周型}

\subsubsection{态完整后才开始计算，提供转折变盘的有效参考。}

104

\clearpage

\subsubsection{技术篇}

\subsubsection{（五）2011-2012耀华日线}

2367（甘蛋楼） M12.28 12.22 12.04

\section{2集森指格34卖14.62}

\subsection{14.34}

[14.01]

\subsection{13.88}

13.50] [13.45] 13.40

\subsection{12.98 12.94}

\section{20天20天12.47}

\subsection{12.00}

\subsubsection{11.57第二十一天带量}

\subsubsection{突破颈线往上变盘11.54}

11.06

10.59] 10.59

破底翻 10.12 25069 20055 15234 10028 5013.80 0.000

\section{耀华日线2011年12月下旬破底翻企稳确立，此时便可计算出左边的整理期}

\section{共20天，表示右边也会有对称20天非盘即涨的多头安全期，结果第二十一天带}

\section{量突破颈线往上变盘（包括突破三个交易日前疑似带量不涨的高点化解危机）， 三个交易日后脱离底部区完成头肩底，操作上第二十一天带量突破颈线变盘时多}

\section{单进场，停损设短期整理区低点11.29元风险不到5\%，底部完成时，再计算头肩}

底的等幅涨幅满足，做为获利点参考即可。

\subsection{大师语录}

\section{“一个人必须相信自己才能在这行生存。我从不接受别人的}

\subsubsection{点子或内幕消息。}

杰西·利佛摩尔（JesseL.Livermore）

105

\clearpage

\section{2}

\section{进阶篇}

109

\subsection{76.00]}

跌幅

森指標 成交量 价涨量增量价背离

\clearpage

多空转折一手抓

\subsection{一、大盘}

接下来看这10几年来的台股各时期量价结构与型态，只要熟悉技术篇型态 结构的方法，应该都不难掌握较大趋势的转折机会。观察大盘，随大盘趋势，操 作个股才会有更高的成功率。了解型态、掌握多空趋势才能随势操作，而作期指 的大多是短线操作，这是不对的。指数的操作也是必须波段操作，一年可以做的 波段没有几次，较大的波段甚至可能少至一两次，问题是做得到的波段趋势利润 极为可观，绝对不是短线操作者可以想像的，投资人千万不要尝试短线，那会成 为习惯，必败无疑。做短线的，要及时回头戒掉恶习，才有成功的机会。底下几 乎是连接每个年份的大盘多空趋势，只要胆大心细，一旦出现难得的转折，无论 是多头转折或空头转折，积极介入，严设停损，小赔大赚，一年只要抓到一两次 的较大波段转折，便是成功的操作。

\subsubsection{（一）1997-1998台股日线}

0000大指他（日乘他） 971高9030.558969.56收9030.55 回体指10256.10] 10264 9784.54

\subsubsection{破线头部确立9/2 9378.529337.61}

\subsection{4/14 破线头部确立8841.41 8532.25}

8116577/30 8361.86 破线头部确立796223

\section{假突破}

\subsection{7488.43}

7375,14 7218.39 7418.73 7040.54] 7073.22 6939.18 爆量后6467.62 价涨量缩跌幅满足

量价背离908869

108

\clearpage

进阶篇

·重点提示：量价背离、跌幅满足、假突破。

大盘1997年7月17日爆2969亿元天量，技术面大量后通常都还会有高点， 随后价涨量缩量价背离。同年9月2日长黑损破颈线头部，确立为长线卖出信号 （红色虚圈部分，后续的红色虚圈都是反弹后破线头部确立），结束了1995年8 月15日4474低点盘底后，长达约2年逾5500点的大多头行情。

观察重点：

进入长空后的反弹波至9378高点盘头至1998年4月14日破颈线头部确立， 直到跌破前低7040点后，便可计算跌幅满足。 0自10256高点跌至7040低点，共跌3216点。 ②自9378高点-3216点=等幅计算的跌幅满足6162点。

大盘达跌幅满足附近最低6219点企稳反弹，只是约三个月约1200点的缓步 反弹，在长线的结构上属于相对弱势反弹，反弹波在假突破后结束，随后跌破颈 线再跌一波才见底。

?操作方式：

01997年9月2日长黑损破颈线，头部确立空单进场，约可空在9340点附 近，突破9650点也就是站回颈线时停损，风险约3.3\%（本波可计算头部 的等幅跌幅满足，共连跌三波） 21998年4月14日破线头部确立空单进场，约可空在8735点附近，突破 8980点也就是站回颈线时停损，风险约2.8\%。 ③同年1998年7月30日破线头部确立，空单进场，约可空在7780点附近， 突破7950点，也就是站回颈线时停损，风险约2.1\%，达跌幅满足附近， 空单回补获利约20\%～34\%。

109

\clearpage

\subsubsection{多空转折一手抓}

\subsection{（二）1997-1999台股周线}

高7659707.22

\section{10256.10] 10256}

\subsubsection{9722.07}

9188.04

\section{B532 8644.95}

\subsection{8110.92}

\subsection{7488.43 7576.89}

7375.14

\subsection{2073.22 7033.81}

\subsection{跌幅满足656313}

5965.74

\section{跌幅满足}

\subsection{5422.66}

森指理

\subsubsection{>重点提示：跌幅满足。}

\section{台股1997年8月高点10256跌至7040点，反弹至9378点，然后回跌破7040}

\subsection{点，便能计算跌幅满足。}

\section{观察重点一：}

\subsection{绿色虚线部分}

\subsubsection{0自10256高点跌至7040点，共跌3216点。 ②自9378高点-3216点=6162点。}

\section{1998年9月台股跌至跌幅满足附近（最低6219点）后进入约两个半月的反}

\section{弹，反弹虽达1269点（最高11月下旬7488点），但就10256高点累积跌4037点}

110

\clearpage

进阶篇

的结构来看仍属相对弱势反弹。

随后同年12月中旬破线，隔一个月1999年1月上旬跌破前低6219点便可计 算跌幅满足，这次的跌幅满足有两个高点可供计算。

观察重点二：

紫色虚线部分 0自9378高点至7073，共跌2305点。 ②自8116高点-2305点=5811点。

绿色虚线部份 ③另自8116高点至6219点，共跌1897点。 4自7488高点-1897点=5591点。

台股1999年2月上旬达跌幅满足附近（最低5422点），仅分别停留1-2周便 做破底翻结束长空翻多回升。

操作方式：

参考前一篇“1997-1998台股日线”，本篇周线做为辅助操作。

大师语录 “多样化只是一种避险的无知。

威廉·欧尼尔（WilliamO'Neil）

\clearpage

多空转折一手抓

\subsection{（三）1998台股日线}

\subsubsection{0000大路（日）8116.57网：52454高6534.15166412.72收5476.80 497031451（4.76）日期99}

\subsection{8116.57}

\section{破线头部确立7703.76] 7818.37}

\section{7488.43假突破7520.17}

\subsection{7221.97}

\subsection{6923.77}

6620.516554.81

\subsection{632231}

6219.89破底翻

\subsection{5725.91}

5422.66

\subsubsection{重点提示：跌幅满足、破底翻}

\section{台股1998年12月上旬假突破后再跌一波，技术面假突破确立通常都会跌破}

\subsubsection{前低（6219点），跌破前低后便能计算跌幅满足。}

\section{观察重点：}

\subsubsection{0自8116高点跌至6219低点共跌1897点。 ②自7488高点-1897点=等幅满足计算的跌幅满足5591点。}

\section{大盘跌至跌幅满足附近（最低1999年2月5日5422点）企稳，短线整理后连}

\section{续二个跳空长红，在第二根长红（1999年2月22日）站回5988颈线附近破底翻}

\section{确立，技术面破底翻确立后通常都会突破前高，盘势随后盘涨突破前高7488点。}

112

\clearpage

进阶篇

·操作方式：

破线头部确立时空单进场约可空在7750点附近，停损设在7830点附近，也 就是站回颈线反压，也是翻黑高点，风险仅约1\%，反压站不回看跌幅满足附近， 空单获利约30\%

小叮：

0本篇连接前两篇“1997-1998台股日线”，弱势反弹后的假突破接续操 作。 ②同样也是以前一篇“1997-1999台股周线”做辅助参考操作。 ③如同许多例子一样，假突破翻空后，最终以破底翻企稳买进信号结束空 头。

大师语录 “很多交易员碰上商品（或股票）仓位下跌时，就采取对冲 避险来保护自己，也就是说抛空某些商品（或股票）以弥补亏损。 这实在是大错特错。”

威廉·江恩（W.D.Gann）

113

\clearpage

多空转折一手抓

\subsubsection{（四）1999台股日线}

0000大业（日运） 0090高0028.397937777979.40量16240739室1215617.362175）

\section{8678 8710.71 8710.71}

\subsubsection{7/13 8347.42}

\subsubsection{8058破线翻黑}

\subsection{A7924头部确立7984.14}

\subsection{7705,63}

\section{7488.43假突破677长红7614.69}

7249突破颈线

\subsubsection{7209.51] 7304.44 7251.41}

\subsection{6888.12}

6577.926550突破颈线底部确立6771.7 6518.68 95A79 破底翻6155.39

\section{559507 5422.66 298815}

重点提示：破底翻、涨幅满足。

台股日线破底翻后于同年1999年3月15日长红突破颈线，底部确立，技术 面当日长红低点6440点不破，看底部计算的涨幅满足附近。

\subsection{观察重点一：}

0颈线约在6550点-低点5422点=1128点。 ②突破颈线的位置约在6533点+1128点=7661点。

台股4月下旬达涨幅满足附近便进入约两个月的整理期，同年6月7日突破 整理区上缘颈线再创新高后，便可计算第二波涨幅满足。

114

\clearpage

进阶篇

第一波涨幅满足7661点+1128点=等幅计算的第二波涨幅满足8789点，本 波涨至8058短线整理再创新高后，也可计算涨幅满足，都在大约的位置。

观察重点二：

①自低点7304点至短波高点8058点，共涨754点。 ②短波高点后的整理区低点7924点+754点=8678点。

台股达涨幅满足附近（最高8710点）进入高档震荡盘头后，7月13日带量长 黑破线将近一个月的头部确立，5422点以来的波段多头行情也告一段落，大幅回 跌。

操作方式：

1999年3月15日长红突破颈线底部确立买进多单，约可买在6580点附近， 跌破长红低点6440点，也就是跌落线下时，停损风险约2.2\%，达涨幅满足7661 点附近多单减码，再达第二波涨幅满足附近，尤其是7月13日破线翻黑头部确立 时，多单获利了结约16\%\textasciitilde{}31\%

大师语录 “我更关心的是，怎么控制失利状况。要学会接受亏损。想 赚钱，最重要的就是不能让亏损失控。

马蒂·史华兹（MartySchwartz）

115

\clearpage

多空转折一手抓

（五）1999-2001台股周线

\subsubsection{10394}

涨幅满足 9747.39

\subsubsection{9209.482000/7/1 9101:20}

8710.71跌破颈线头部确立

8455.00

\subsection{8281.68}

7808.80

7261311999/12/25 7151.65

\section{边长5771.75突破三角形整理区6505.45}

约2180 6198-22

\subsection{5256.B7 跌幅满足 带量破底翻4555.91}

→重点提示：涨跌幅满足、三角形整理、双头、破底翻。

涨幅满足、三角形整理： 台股1999年7月达8710高点多头后告一段落，大幅拉回并做长时间震荡整 理，整理时间达约半年，后于年底12月25日当周突破三角形整理区上缘颈线， 2周后突破8710点前高便能计算涨幅满足，有两种计算方式的满足点都落在大致 的位置上。

\subsubsection{观察重点一：}

绿色虚线1部分为大波段涨幅满足。 0自5422低点涨至8710高点，共涨3288点。 ②自6771低点+3288点=10059点。

116

\clearpage

进阶篇

观察重点二：

深蓝色部分为三角形等幅满足计算方式，突破颈线后的涨幅=三角形边长的 等幅。 0三角形边长高点8710点的直线约2180点。 2突破三角形上缘颈线的位置约在8020点附近+2180点=10200点（颈线反 压突破变支撑，不破看涨幅满足附近，也就是买进停损风险约2.5\%，获利 约22\%）

双头、跌幅满足、破底翻：

台股2000年2月达涨幅满足附近（最高10393点）进入盘头，同年7月1日 当周跌破颈线头部确立，颈线跌破后，便可计算跌幅满足。

?观察重点三：

010393点与10328点双头两个高点，连成直线的中间位置约在10360点附 近。 210360点至颈线的位置约在8385点附近幅度1975点。 3跌破颈线的位置约在8550点附近-1975点=6575点。

而这波空头直接进行两波跌幅，可见空头未见企稳迹象前抢反弹风险颇高， 第二波跌幅满足计算，便是第一波的跌幅满足6575点再减1975点=4600点，要 注意台股当年底达跌幅满足附近（最低4555点），才出现带量破底翻做较大幅 度反弹。而7月1日跌破颈线做空的停损风险约6\%，风险稍高，但波段获利更 大，约44\%

117

\clearpage

多空转折一手抓

（六）2001台股日线

\section{2/16假突破6198.22}

\subsection{5890.34 -558247}

291

颈线5269.37

\subsection{假突破4653.62}

\subsection{4340.53}

跌幅满足

\subsection{4032.65}

10/11 3761.30 跌幅满足跳空长红 8910底部确认3411.68

重点提示：跌幅满足。

台股日线2000年12月28日达前一篇“1999-2001台股周线”计算跌幅满足 附近，最低4555点弹至2001年2月16日6198点，出现假突破，结束1643点的反 弹行情，盘跌的结构至同年6月下旬，跌破短期反弹低点4902点，便能计算跌幅 满足。

\subsubsection{观察重点一：}

绿色虚线1部分 0自6198高点跌至4902低点，共跌1296点。 ②自5318高点-1296点=4022点。 台股跌至同年7月24日达跌幅满足附近（最低4008点）反弹。

118

\clearpage

进阶篇

弹至4715高点后，隔日跌回线下，出现假突破，结束反弹，再跌一波。同 年9月13日跌破反弹低点4008点，便能计算跌幅满足。

观察重点二：

绿色虚线部分 0自5318高点跌至4008点，共跌1310点。 ②自4715高点-1310点=3405点。 台股跌至同年9月下旬跌幅满足附近（最低9月26日3411点）打出短底，于 10月11日跳空长红突破底部颈线，结束6198点以来的空头型态，也结束10393 点长期盘跌的长空进入回升多头结构。

操作方式：

02001年2月16日突破后隔日立即跌回线下，假突破确立空单进场，约可空 在5940点附近，突破前高6198点停损，风险约4.3\%，达首波跌幅满足附 近（绿色虚线0部分），空单获利了结约32\%。 ②同年8月20日跌落线下假突破确立，反弹波结束空单进场，约可空在4540 点附近，突破前高4715点停损，风险约3.8\%，达跌幅满足附近（绿色虚 线②部分），空单回补获利约21\%。

大师语录

“我有认错的勇气。当我一发现犯错，马上改正，这对我的

事业十分有帮助。

乔治·索罗斯（GeorgeSoros）

119

\clearpage

多空转折一手抓

（七）2002台股日线

0000大（日）

\subsection{4假突破6490.69}

\subsection{5/2长黑破线-6145.38}

\subsection{5466.27 5120.96 4781.41}

4412.42价涨量缩4441.85

价量背离 4096.54 3756.99 10/11跳空长红3613.11

\section{3411.68底部确立3411.68}

610768

重点提示：量价背离、涨幅满足、假突破。

台股3411点短底完成回升，达4722点压回再做突破时，便可计算涨幅满足。

\subsubsection{观察重点：}

0自3411低点涨至4722高点，共涨1311点。 ②自4376低点+1311点=5687点。 台股2001年底达涨幅满足附近后震荡两个多月才进入缓步盘涨。

技术面2002年3月再创新高，缓步盘涨的过程中，很明显为价涨量缩的量价 背离现象，直到2002年4月下旬出现假突破，多头露出败象，尤其在5月2日长 黑损破颈线，确立了假突破的长线卖出信号，同时也确立头部，进入大幅回跌，

120

\clearpage

进阶篇

至接近起涨点附近（最低3845点）。

操作方式：

2001年10月11日跳空突破颈线，买进约可买在3750点附近，跌破颈线3600 点附近停损，风险约4\%，达涨幅满足5687点附近，获利约51\%。

小叮：

①达涨幅满足后未出现确定的方向前，以观望为主。 ②当量价背离的结构出现后，再来注意转折的出现。 ③假突破便是转折信号，操作方式请参考下一篇“2002台股周线”。

大师语录

“心理造就90\%行情。认为影响股价（特别是短期走势）起

伏的原因，主要并非受经济发展的影响，而是受制于投资者对消 息的反应及心理因素。好消息未必一定令股价上升，大众对消息 的反应才是影响股价的主要因素。”

安德烈·科斯托兰尼（AndréKostolany）

121

\clearpage

多空转折一手抓

\subsection{（八）2002-2004台股周线}

0000大推指（） 4 (2.1905]

\subsection{7135.00}

涨幅满足

\subsection{3/19 6723.62}

\section{假突破：涨幅满足枪击事件}

100L/7074 6182.2 6312.24

5893.89

\subsection{5482.52}

\subsection{5450.72 5329.12}

6/21 5255.065071.14 带量突破颈线底部确立

\section{465279}

\section{D破底翻5/31}

\subsection{4241.41}

\section{跌幅满足7 4044.收敛三角形底部形态3830.03}

\subsection{1210292 9664272}

\section{0.000}

重点提示：跌幅满足、破底翻、收敛三角形底部、涨幅满足。

\section{台股周线2002年5月上旬假突破翻空，技术面4808点反弹至5460点，回跌}

再破前低4808点后，便可计算跌幅满足。

\subsection{观察重点一：}

\subsubsection{0自6484高点减4808低点，共跌1676点。 ②自5460高点-1676点=等幅计算的跌幅满足3784点。}

\section{打底的结构约半年，至2003年5月底翻红，站回颈线，出现破底翻多头买进}

\section{信号，3周后带量突破颈线，收敛三角形底部型态确立，便可计算涨幅满足。}

122

\clearpage

进阶篇

观察重点二：

0三角形边长约1415点（3845点至边长高点，约在5260点附近的距离）。 ②突破颈线的位置约在4810点附近+1415点=6225点。 台股盘涨至同年11月涨幅满足附近（最高6182点），进入整理约两个半月， 才再度突破整理区颈线，续涨一波。

观察重点三：

技术面突破前高6484点，便可计算大波段涨幅满足。 ①自3411低点涨至6484高点，共涨3073点。 ②自4044点+3073点=等幅计算的涨幅满足7117点。 台股4044低点起涨的多头波至2004年3月上旬，达涨幅满足附近后（最高 7135点），震荡3周，才在当时的“3·19”枪击事件利空后结束多头翻空。

操作方式：

02002年5月4日当周跌回颈线以下假突破确立放空，约可空在5980点 附近，突破当周高点6258点停损风险约5\%，达跌幅满足回补，获利约 36\% 22003年5月31日当周破底翻，约可买在4470点附近，跌破破底前的前低 4240点，停损风险约5\%，3周后6月21日当周带量突破颈线加码买进，约 可买在4930点附近，跌破颈线也是前一周低点4715点，停损风险约4.5\%， 达收敛三角形涨幅满足6224点附近，多单获利约26\%～39\%。 32004年1月10日当周带量突破颈线多单买回，约可买在6150点附近，跌 破颈线约6020点停损风险约2\%，达大波段涨幅满足附近，多单获利约 15\%

123

\clearpage

多空转折一手抓

\subsubsection{（九）2004台股日线}

\subsubsection{7135.00319枪击利空高点7135.00}

\subsection{一2004/4/22带量破线6792.93}

\subsubsection{破线不带量5/5 6450.85}

\subsection{613555}

带量破底翻6131.97

5766.70

\subsection{08/19}

\subsubsection{541883}

\subsection{5076.75}

4734.68

4392.60

\subsection{4044.73}

\section{→重点提示：无预警利空的主力心态、破线有无带量、破底翻。}

主力作价偶尔也会遇到突发性利空落跑不及，主力来不及跑，大多会有一波 拉高解套波，因此利空短线大跌后，大多还是会拉高到利空起跌点附近，做主力 解套的逃命波，主力跑了才真正会大幅回跌。

\section{台湾地区股市业内盘特性，业内具滚量作价能力，一旦介入行情波动相当}

大。如图示，2004年3月上旬达大波段涨幅满足附近震荡，高档连续作价出量主 力便有调节获利机会，但一波多头行情需在高档震荡一段期间，主力筹码才有较 充裕的出货时间，这次震荡时间太短，遇到当年“3·19”枪击突发性利空，短 线跳空大跌主力无法完全顺利出货，因此3月下旬至4月下旬约一个月的盘涨拉 高至利空起跌高点附近震荡。

124

\clearpage

进阶篇

观察重点：

拉高后于当年4月22日出现一根带量长黑破线翻空信号，技术面当日高点 6912点已成反压，是翻空操作的空头停损点，结构上3月和4月也形成约两个月 的高档震荡期，足以让4044点盘涨以来的多头主力有充裕的时间出货，4月22日 带量长黑破线，便是主力出货完毕的退出翻空信号，主力跑了台股自然大幅下 跌，破利空低点。

短线急跌至大颈线区附近仅支撑2日，在5月5日又见长黑破长达约八个 月的大颈线支撑（可见主力落跑后切勿在颈线支撑做所谓的逢低承接），问题是 接下来的这次的长黑破线并未带量，因此虽然破线也完成的反弹逃命线型态，短 中期盘弱持续空头走势，但随后的跌幅并不太深，未出现长空走势而是中期修 正，盘弱的走势至同年8月19日带量长红，突破约一个月的整理区颈线，随后带 量站回颈线，破底翻企稳确立，结束了中期拉回的修正波。接下来振荡整理的时 间超过一年。

操作方式：

2004年4月22日带量翻黑破线做空，约可空在6720点附近，突破长黑高点 6912点停损风险约3\%，8月19日量突破下降趋势线，带量长红突破颈线，底部 确立，空单回补约在5550点附近，获利约17\%

大师语录 “多数的交易员因为缺乏纪律而遭致失败而非因知识不足而 失败。

安德烈·科斯托兰尼（AndréKostolany）

125

\clearpage

多空转折一手抓

\subsubsection{（十）2005台股日线}

0000大编指做（日星）

\subsubsection{291819假突破648334}

\subsection{6380.39}

\subsubsection{6267.52涨幅满足6279.16}

\subsection{6186.126171.87 6177.93}

6088.12

\section{6074.99}

\subsubsection{上升趋势线5976.40 5973.76}

\section{5769.58}

\section{S734.87跌幅满足}

5668.36

\subsection{5618.90] 5585.41}

5565.41

\section{033}

\section{0}

\subsubsection{→重点提示：破底翻、涨幅满足、假突破、跌幅满足。}

\section{2005年4月18日台股日线跌破长期整理的颈线，直到同年5月5日带量长红}

破底翻企稳确立，6月初突破整理区，便可计算涨幅满足。

\subsection{观察重点一：}

\subsubsection{0自5565低点涨至6001高点，共涨436点。 ②自5879低点+436点=等幅计算的涨幅满足6315点。}

\section{台股涨至同年6月中旬达涨幅满足附近（最高6401点），便进入高档震荡。}

126

\clearpage

进阶篇

同年8月上旬出现假突破卖出信号，随后8月15日跌破上升趋势线进入 盘跌趋势，同年9月下旬跌破反弹低点，便可计算跌幅满足。

?观察重点二：

0自6481高点跌至5976低点，共跌505点。 ②自6186高点-505点=等幅计算的跌幅满足5681点。

只是这波曾短暂拉回破底的颈线之上，但未过反弹高点6186点，后续再破 底后就型态而言，仍续看跌幅满足。台股跌至同年10月20日最低5618点，然后 止跌回升。

操作方法：

02005年5月5日破底翻买进约可买在5930点附近，停损设在当日长红低点 5851点，也是颈线附近，风险约1.4\%，达涨幅满足6315点附近获利约 6.5\% ②同年8月8日假突破确立做空，约可空在6380点附近，停损设在前高6481 点，风险约1.6\%，达跌幅满足附近5681点附近，获利约11\%

大师语录 “要学会迅速而利落地砍断亏损。不要期待自己永远不犯错， 而是在错误发生时，尽快砍掉亏损。

伯纳德·巴鲁克（BernardBaruch）

127

\clearpage

多空转折一手抓

\subsubsection{（十一）2005-2006台股日线}

0000大额指（日蛋）号720375座722935713621收7136.21量17114 32.726-8.457\%5瓶2005/4/2

\section{森指楼涨幅满足24 7476.077476.07}

7476.07

\subsection{7064.91}

\subsection{6747.436726.73趋势线跌破上升6859.34}

4/3突破颈线-6653.76

\subsection{6436.14破底翻6444.70}

\section{跌破：6373.63 6344.73}

6171.87上升趋势线6239.12

\section{6033.54}

10/31 5827.96

\section{068195破底翻5618.90}

5618.90

重点提示：破底翻、涨幅满足。

2005年台股日线达跌幅满足后，同年10月31日出现破底翻企稳信号，随后 一路回升，直到隔年2006年1月18日长黑跌破上升趋势线，多头才告一段落， 整理约两个月。

同年3月底再度出现破底翻买进信号，随后突破整理区颈线，第二波涨势确 立，盘涨越过前高6797点后，便可计算涨幅满足。

\subsubsection{观察重点：}

0自5618低点涨至6797高点，共涨1179点。 ②自6344低点+1179点=等幅计算的涨幅满足7523点。 技术面4月3日突破整理区颈线后颈线压力变支撑，支撑不破看涨幅满足。

128

\clearpage

进阶篇

台股同年5月上旬达涨幅满足附近（最高7476点），5月10日带量回跌，隔 日跌破上升趋势线，结束当年多头大幅回跌。

操作方式：

02005年10月31日破底翻买进，约可买在5750点附近，跌破5618点，停损 风险约2.2\%，2006年1月18日跌破上升趋势线卖出，约可卖在6550点附 近，获利约14\%。 2同年3月30日破底翻买回，约可买在6530点附近，跌破破底翻的前低 6436点停损风险约1.4\%，达涨幅满足附近，获利约14\%。

小叮：

这次达涨幅满足的反转相当快速，是属于V型反转，若出脱不及可以三日低 点为参考，当三日低点跌破时，也同时跌破上升趋势线，涨高后直接V型反转在 台股指数属于较为少见的现象。

大师语录 “关键在于判断市场处于哪个阶段。

安德烈·科斯托兰尼（AndréKostolany）

129

\clearpage

多空转折一手抓

\subsubsection{（十二）2006台股日线}

\subsection{0000大期（）体指6401.55 7476.07}

\section{假突破17476.07 7478.33}

\subsubsection{j5/11 734270}

-7209.32 跌破

\subsubsection{上升趋势线7073.69}

\subsection{6967.38}

6940.31

\subsection{679111}

\section{6671.30}

6537.93

\subsection{跌幅满足6404.55}

\subsection{5290.53 6268.92}

6/96268.92268929

\section{3876878 0.000}

\subsubsection{→重点提示：假突破、跌幅满足}

2006年5月上旬台股日线达涨幅满足附近，后拉回跌破上升趋势线，同时在 型态上也属假突破确立的卖出信号（颈线突破压力变支撑，多头看支撑，跌破支 撑便属转弱），盘跌反弹后同年6月5日长黑跌破反弹低点，便能计算跌幅满足。

\subsection{观察重点：}

0自7476高点跌至6791低点，共跌685点。 ②自6976高点-685点=等幅计算的跌幅满足6291点。 台股同年6月9日达跌幅满足附近（最低6268点）止跌进入打底期。

130

\clearpage

进阶篇

操作方式：

2006年5月11日跌破上升趋势线，也跌破5月4日带量突破整理区的红K低 点7278点，空单进场约可空在7270点，停损设5月11日破线黑K高点7361点风 险约1.5\%，达跌幅满足附近，空单获利了结约13\%

小叮：

02006年5月11日跌破上升趋势线及带量红K低点，同时也是上篇“2005一 2006台股日线”提到的三日低点被跌破。 26月5日跌破反弹低点6791点时，也是空单加码点，同样站回线上停损， 达跌幅满足附近，空单获利了结。

大师语录 “大多数菜鸟散户在亏损还算少的时候死抱不放。他们原本 可以在受伤不重的情况下退出，却受制于情绪和人性，一厢情愿 地一再延误，造成亏损越来越大，终致损失惨重。，

威廉·欧尼尔（WilliamO'Neil）

131

\clearpage

多空转折一手抓

\subsection{（十三）2006-2007台股周线}

0000大指数（据） 51773高6582.216436.14收6490.58量41018520至250648962.9800.961\%）期2006/3/11 9807.91

\section{9807.91}

涨幅满足 9345.08

888225

8411.57

\subsection{涨幅满足： 7987.61] 7948.74}

\section{10u476.07突破颈线}

涨幅满足-7485.91

7306.07

\subsection{7015.24}

突破颈线底部确立6552.41

\subsection{4.73 6268.926232.49 6089.58}

\subsection{5618.90}

成文量第森指程5618.90

106093

000

\subsubsection{重点提示：涨幅满足。}

\section{台股周线2006年6月达跌幅满足后，同年9月下旬突破整理区上缘颈线，完}

\section{成约三个月的底部，颈线6730点附近压力变支撑，支撑不破，看一路盘涨过前}

高后的涨幅满足附近。

\section{图示为5618低点计算三个波段高点的涨幅满足。绿色虚线1部分同日线请}

\subsubsection{参考“2005-2006台股日线”}

\subsubsection{观察重点一：}

\subsubsection{红色虚线部分}

\subsubsection{0自5618低点涨至7476高点，共涨1858点。}

132

\clearpage

进阶篇

②自6232低点+1858点=等幅计算的涨幅满足8090点。 台股盘涨至隔年2007年1月中旬达涨幅满足附近（最高7999点），然后进 入约三个月的整理期。

整理期间的日线3月初破底，3月中旬出现破底翻买进信号（周线为下影 线），图示周线4月7日当周突破颈线，越过7999点前高，便可计算涨幅满 足，同样颈线突破压力变支撑，支撑不破，看涨幅满足附近。

观察重点二：

绿色虚线2部分 0自5618低点涨至7999高点，共涨2381点。 ②自7306低点+2381点=等幅计算的涨幅满足9687点。 台股盘涨至同年7月下旬，达大波段涨幅满足附近（最高9807点）便爆量反 转，随后大区间震荡作头进入长空型态。

操作方法：

02006年9月23日当周突破颈线，买进约可买在6870点附近，跌破颈线 6730点附近，停损风险约2\%，达涨幅满足附近，获利约13\%。 22007年3月20日破底翻买进，约可买在7750点附近，跌破破底翻前的前 低7637点，也是破线确立，停损风险约1.5\%，达绿色虚线2部分的大波 段涨幅满足附近，获利约25\%。

133

\clearpage

多空转折一手抓

\subsubsection{（十四）2006-2007台股日线}

0000大指数（甘爱）

\section{复合式头肩顶7999.42/头7990.55双右肩8069.61}

\subsection{795020跌破颈线}

7899.19

\subsubsection{左肩7716.85 7731.61}

7637.54]破底翻7564.03

7393.61

\subsection{7306.07}

\subsection{7058.45}

688803

6720.45

\subsection{6550.03}

净森指楼6550.03

重点提示：复合式头肩顶、跌幅满足、破底翻。

图示便是承前文“2006-2007台股周线”，所提“整理期间的日线3月初破

\section{底，3月中旬出现破底翻买进信号（周线为下影线）”，2007年3月1日带量跌}

破颈线，确立了约三个月的头部完成，便可计算跌幅满足。

\section{图示为复合式头肩顶，除了有左肩与右肩外，还有与M头一样的双头，测}

\subsubsection{量的方式与M头相同。}

\subsubsection{观察重点：}

\section{①双头高点连接与颈线的直线距离=跌破颈线位置后的同等距离。 ②双头高点连接的中间值约7996点，与颈线7599点距离约397点。}

134

\clearpage

进阶篇

③跌破颈线（3月1日）的位置约在7730点-397点=等幅计算的跌幅满足 7333点。 ④技术面颈线压力未站回前看跌幅满足。

台股仅约四个交易日便达跌幅满足附近（最低7306点），而这次是以V型 反转直接站回颈线，出现破底翻买信号，随后盘涨再创新高，进入前面周线所提 的绿色虚线?部分的大波段涨势。

·操作方式：

2007年3月1日带量跌破颈线做空约可空在7670点附近，站回颈线约7760点 附近停损，风险约1.2\%，达跌幅满足7333点附近，获利约4.4\%。

小叮：

①这次计算跌幅满足获利幅度有限，可以不作操作。 ②若做空以相对规模来看，欲看第二波跌幅满足者，在破底翻确立后，须做 停损动作，前一篇“2006-2007台股周线”为下影线未破线，也是重要的 判断工具。

大师语录 “股票投资是一门艺术，历史和哲学在投资决策时显然比统 计学和数学更有用。

安德烈·科斯托兰尼（AndréKostolany）

135

\clearpage

多空转折一手抓

\subsection{（十五）2007台股日线0}

0000大指数（日强） 38.49高8388.49336.53收346.30 137023453 224372380：285

保指9807.91

爆量假突破9807.91

-9501.75 9219 9195.59

8884.24

8578.08

\subsection{8271.91}

\subsection{7990.55] 7950.20}

\section{-7960.56}

\subsection{7654.40}

破底翻

\subsection{7306.07 734824}

\subsection{7036.89}

\subsection{→重点提示：假突破、跌幅满足。}

\section{台股日线2007年7月下旬到达“2006-2007台股周线”所计算的大波段涨幅}

\section{满足附近后（绿色虚线部分），7月26日出现爆出几乎与历史天量相同的历史次}

\section{天量，重点这是爆量假突破的结构，极为明显的长线卖出信号。}

\subsection{观察重点：}

假突破跌深反弹后再破前低，便能计算跌幅满足。 ①自9807高点跌至8727低点，共跌1080点。 ②自9219高点-1080点=等幅计算的跌幅满足8139点。

\section{台股同年8月17日达跌幅满足附近（最低7987点）仅1日便大幅弹升回高点}

\subsubsection{136}

\clearpage

进阶篇

做第二个头，整个逾1500点的高档大区间头部震荡期约半年。

所谓做空赚得快，2007年5月21日突破8100点，越过前高，花了两个多月 的时间才到达高点，但9807点爆量反转至跌破8000点，共跌1807点的时间仅半 个多月，这便是空头迷人的地方，掌握时机获利的速度颇为惊人。

·操作方式：

2007年7月27日首次受次贷风暴消息面影响，跳空大跌才确立了技术面7月 26日的爆量假突破，因为当天大幅跳空的关系，空单介入的风险相对较高，约可 空在9250点附近，停损设在9566点也就是前一日收盘补空的位置（利空跳空的 缺口通常不会被回补），风险约3.4\%，达跌幅满足附近空单回补获利约12\%。

小叮：

此处为消息面影响急回，虽有空头波段利润，但头部结构不完整并非长空时 机，操作一波空单后，需等待较完整的头部结构出现。

大师语录

“职业投资者的工作，95\%是在浪费时间，他们在阅读图表及

营业报告，却忘记思考，但对投资者来说，这才是最重要的。

安德烈·科斯托兰尼（AndréKostolany）

137

\clearpage

多空转折一手抓

\subsection{（十六）2007台股日线②}

0000大指街

\subsection{9997.74}

\section{978375量价背离9863.19}

爆量假突破假突破

9650.75

9441.85

\section{9275.72 923295}

9020.51

\section{8811.61}

8602.70

839026

-818136

811028

\subsection{7968.92}

\section{2586987}

\subsubsection{重点提示：涨幅满足、量价背离、假突破。}

\section{承前文8月17日达跌幅满足附近后，以V型反转的方式回升，7月26日的爆}

\section{量，通常在正常的情况下大量后大多还会有高点，当次却出现假突破的卖出信}

\section{号，主要是7月27日突发性的次贷风暴令市场措手不及，因此只要注意量价结构 的变化，仍可在假突破后全身而退，甚至反手做空。}

当然高档爆量，表示主力已开始出货做获利了结的动作，但大多仍需在高档震

\section{荡盘头，才能完全顺利出货完毕，而这次的爆量因突发性利空，主力无法在高档完}

全顺利出货，如前面提到的“3·19”枪击事件，因此才有达跌幅满足后的V型反 转再度拉高，且突破爆量高点9807点（最高9859点），还是完成了大量后还有高 点的结构，所谓大量后几乎都还有高点，就是主力高档震荡出货手法产生的现象。

138

\clearpage

进阶篇

图示达跌幅满足弹升一波后整理，同年9月26跳空长红，越过整理区高点便 能计算涨幅满足（整理区上缘9078点附近压力突破变支撑，支撑不破看满足点）

观察重点一：

0自7987低点涨至9078高点，共涨1091点。 ②自8845低点+1091点=等幅计算的涨幅满足9936点。

台股同年10月底涨至涨幅满足附近（最高9859点）再度出现假突破反转。

操作方法：

这波涨高至10月高档震荡了约一个月，已足以让主力顺利完全出货，10月 29日跳空长红突破整理区随后再创新高，已明显量价背离，随后便是假突破卖出 信号，头部完成后，曲终人散，以空单操作为原则。以这次为例，假突破后整理 区上缘9670点附近为反压，空头看反压（反压是做空的停损点），反压未突破 前，一路顺势做空直到趋势改变为止，这次量价背离假突破，产生了第二个高点 头部，才是长空的开始。

量价背离

通常涨高后的量价背离，大多是主力已出货不再积极滚量作价，且相对高 档市场认同度不高，因此价创新高量却未过高，反而出现量缩趋势是量价背离 警讯，很容易引发市场筹码的反手杀出，形成追高者套牢的头部区。

139

\clearpage

多空转折一手抓

\subsubsection{（十七）2007-2008台股日线}

\subsubsection{0000大姐（日）森指9859.6520674第385608206:74838160量1270677057992282.230.9915}

\subsection{9859.65 9859.65}

心量价背离假突破

\section{缺口566056 9545.49}

\subsection{假突破9231.32}

\section{涨幅满足9049.23}

8911.84

\subsection{98598}

\subsection{8546.20颈线8597.67}

\section{7964.02 7900.80 破底翻7649.86 双底}

\section{跌幅满足7384.61] 7335.70}

LLOISAN 701621

2179139

\section{000}

\section{重点提示：跌幅满足、破底翻、W底、涨幅满足、假突破。}

\section{2007年10月底量价背离假突破，进入一波空头修正，同年12月13日跌破反 弹低点，便能计算跌幅满足。}

\subsection{观察重点一：}

\subsubsection{0自9859高点跌至8207点，共跌1652点。 ②自8804高点-1652点=等幅计算的跌幅满足7152点。}

台股跌至隔年2008年1月下旬达跌幅满足附近（最低7384点，通常都会超 跌），这波有未达满足229点的误差（约2.3\%），不过仍是共跌2475点可接受的 误差，筑出短底后，同年2月中旬破底翻企稳确立。

140

\clearpage

进阶篇

观察重点二：

同年2月25日突破颈线双底确立，便可计算W底型态的涨幅满足。 1双底连线中间值至颈线的距离=突破颈线后的距离。 2双底连线中间值约7510点至颈线约8540点距离约1030点。 32月25日突破时的颈线位置约在8140点+1030点=等幅计算的涨幅满足 9170点。 台股同年4月17日最高9194点达涨幅满足附近后便进入盘头，盘头期至同 年5月中下旬，出现再创新高后的假突破卖出信号，结构翻空后的空头走势，随 著当年次贷金融风暴的失控，共崩跌5354点，相当惨烈。

操作方式：

02007年10月底量价背离后，11月12日跳空跌破颈线，假突破确立，进场 做空，约可空在9300点附近，停损设9598点，也是前一日收盘的缺口（通 常空头第一个跳空缺口大多不会回补），也刚好是站回颈线，风险约5\%， 达跌幅满足附近或破底翻的位置（约7900点），空单获利了结约15\%。 ②破底翻的位置翻空为多，约可买在7900点附近（2008年2月18日），停 损设破底翻前长红低点7707点风险约2.7\%，达涨幅满足附近多单获利了 结约16\% 3这波反弹波达涨幅满足附近进入盘头后，高点出现在同年5月20日9359 点，隔日立即出现假突破为反弹波结束确立空单进场，布空点约可空在 8950点附近，停损设4月14日大量高点反压9194点风险约2.7\%，达下一 页日线计算的跌幅满足4352点附近，或下文“2007-2009台股周线”周线 计算的跌幅满足4225点附近空单获利了结都有50\%以上。

141

\clearpage

多空转折一手抓

\subsection{（十八）2008台股日线}

0000大指（日）高161.480.89 00.537

\subsection{2森指9194.86 930995]}

\section{假突破9319.98}

8718.35

\subsection{8106.69}

\subsection{7376.69 7525.11}

\subsection{6708.46] 6809.96 6341.90}

5740.27

\section{5148.67}

\subsection{跌幅满足4557.06}

\subsection{3955.43}

1922462

4851073

\subsubsection{→重点提示：跌幅满足。}

\section{承“2007-2008台股日线”一文，台股2008年5月下旬假突破进入长空后，}

\section{同年9月1日跌破反弹整理区颈线，隔日跌破前低6708点，便可计算跌幅满足。}

\subsubsection{观察重点一：}

\subsubsection{0自9309高点跌至6708低点，共跌2601点。 ②自7376高点-2601点=等幅计算的跌幅满足4775点。}

\section{这次在达大波段跌幅满足附近时，也是次贷风暴爆发最烈时，全球股市恐慌}

\section{性杀盘带动下台股亦进入超跌，不过达大波段跌幅满足后，即使超跌也大多是相}

\section{对低档区（周线图会更精确的计算出大波段跌幅满足），当时台股甚至实施跌幅}

142

\clearpage

进阶篇

减半，不过在利空下技术面仍可较精确地计算跌幅满足的波段低点。

从7376点进入第二大波跌势，9月18日出现5530点的反弹低点，随后10月 6日再破反弹低点后，便可计算跌幅满足。

观察重点二：

0自7376高点跌至5530低点，共跌1846点。 ②自反弹高点6198点-1846点=等幅计算的跌幅满足4352点。

台股10月28日达跌幅满足附近（最低4110点），便进入约四个多月的打底 期，最低虽曾见到3955低点，但随后破底翻底部确立，进入一波倍数涨幅的多 头趋势。9309点起跌的这波长空趋势中多头伤亡惨烈，说明了随势而为的重要 性，空头成形需随势空单波段操作，除非站上反压（空头看反压）停损利，或达 跌幅满足附近，才作波段空单的获利或减码，直到技术面出现底部企稳信号才翻 空为多，空头千万不要有捡便宜的心态去接掉下来的刀子，尤其是投机抢跌深反 弹，切记随势而为。

?操作方式：

继“2007-2008台股日线”5月20日出现假突破的翻空空点后，9月1日跌破 颈线为技术面又一次空单进场点，约可空在6900点附近，停损设7050点反压， 也是站回颈线时，风险约2.2\%，达跌幅满足附近，空单获利了结约40\%。

143

\clearpage

多空转折一手抓

\subsubsection{（十九）2008-2009台股日线}

0000大款（日）

2森指[7084.83]

\subsection{6950 6739.07}

\section{6198.48涨幅满足6393.31}

\subsection{4/30 6100.08 6047.56}

\section{5695.94}

\section{边长1185下飘旗形}

\section{5095.98}

\subsection{1颈线4580 5004.42}

4652.81

\subsection{4307.05}

\subsection{614918 3955.43}

\section{3955.43}

1957273

重点提示：收敛三角形涨幅满足、下飘旗形。

2009年3月5日大盘日线突破颈线，收敛三角形确立，便可计算涨幅满足。

\subsubsection{观察重点：}

\section{0三角形边长高点5095点至低点约在3910点，距离1185点。 ②突破颈线的位置约在4580点+1185点=等幅计算的涨幅满足5765点。}

\section{大盘达涨幅满足附近，完成第一波涨势后进入震荡（最高6071点，最低}

\section{5571点）。4月30日跳空突破颈线，完成下飘旗形（多头中继站），第二波 涨势确立。}

\clearpage

进阶篇

③第二波涨幅满足计算： 第一波涨幅满足5765点+1185点=等幅计算的第二波涨幅满足6950点。

大盘达第二波涨幅满足后大幅拉回。

操作方式：

02009年3月5日突破颈线多单进场，约可买在4620点附近，停损设4500点 附近，也就是跌落线下时，风险约2.5\%，达两波段涨幅满足附近，多单获 利了结约50\% ②操作上若在跌幅满足附近布局多单，则底部确立时便是积极加码时，稳健 的波段操作者则在确定底部时快速积极进场，停损设在跌破颈线时（技术 面带量突破颈线后拉回便不能再跌破颈线），如此便可在有限的风险下获 取大波段利润，当年大盘在突破颈线后就此喷出大涨，且多头维持了约两 年。

大师语录

“三种迅速致富的可能：一透过带来财富的婚姻，二透过幸

运的商业点子，三透过投机。

安德烈·科斯托兰尼（AndréKostolany）

145

\clearpage

多空转折一手抓

\subsection{（二十）2007-2009台股周线}

\subsubsection{弹森指精0000大验量教（证）S96586 57}

\subsection{9870.71}

9207.31

8554.97

\subsection{7891.58}

\subsubsection{[7376] 7185.52}

\subsection{7239.24}

6586.90

\subsubsection{5923.506133.58}

3/14 5271.16

\subsubsection{4817.44带量突破颈线底部确立}

\subsubsection{跌幅满足4618.83}

\section{4.1g三角形收敛型态3955.43}

\section{8522491 8371956 1301}

\subsection{→重点提示：跌幅满足、三角形底部收敛型态。}

\section{周线除了同日线9309高点计算的大波段跌幅满足外，当年最高9859点也可}

\section{以同样的方式计算跌幅满足，但无论哪种方式计算出来的位置都属相对低档底部}

区。

\subsection{观察重点一：}

\subsubsection{红色虚线部分}

\subsubsection{0自9859高点跌至6708低点，共跌3151点。 ②自7376高点-3151点=等幅计算的跌幅满足4225点。}

\section{台股空头在次贷风暴的助威下2008年11月达跌幅满足附近，企稳盘底。}

146

\clearpage

进阶篇

在9309点起跌的波段跌势中，亦可计算出跌深反弹低点附近位置。

观察重点二：

绿色虚线部分 0自9859高点跌至7384低点，共跌2475点。 ②自9309高点-2475点=等幅满足计算的跌幅满足6834点。

台股2008年11月中旬跌至跌幅满足附近（最低6708点），才短线企稳，进 行约一个月的反弹至7376点，结构上属弱势反弹，因此长空的结构当中，以不 抢反弹逢高空为原则，而这次反弹后的最佳布空点，“9月1日跌破反弹整理区 颈线”便是空单再次进场时机，停损设在站回颈线时，颈线反压未突破站回则波 段空单一路抱至大波段跌幅满足附近，这便是小赔大赚的波段操作方式。 在大盘大跌后，除非有把握的个股逢低承接布局多单（但需注意资金控 管），否则在技术面底部确立后，才做积极介入波段多单的动作。

·操作方式：

图示周线，当台股达大波段跌幅满足附近后，经过约四个月的打底期至2009 年3月14日当周，带量长红突破颈线底部确立，为波段买进时机，结构上明显为 三角形收敛型态，亦可视为3955点及4164点两个低点的双底型态，也就是所谓 的W底，研究并深入了解技术面的量价结构，相对容易切入最佳时机的空点或 买点，弥足基本面的盲点。

147

\clearpage

多空转折一手抓

\subsection{（二十一）2009-2010台股日线}

0000大编指微日 2回指8395.39

\subsubsection{涨幅满足8190.01 840301}

\section{7875.48破线7946.00}

涨幅满足

\section{2 7496.61}

\subsection{破线后高档破底翻7039.61}

\section{相对位置套牢区70090}

\subsubsection{上升趋势线6590.22}

6140.83

\subsubsection{5683.835546.73}

5234.44

4785.05

\subsection{2451164 195727}

\section{.000}

\subsubsection{重点提示：涨幅满足、破线、高档破底翻结构。}

\section{2009年开始的大波段多头中，首波7084高点拉回至6100低点后进行的第二}

\subsubsection{大波中的小波段，亦可计算涨幅满足。}

\section{越过前高至7185点拉回后同年9月7日再创新高，便能计算涨幅满足，且}

\subsubsection{6229低点确立与前低6100点，形成一条上升趋势线。}

\subsubsection{观察重点一：}

\subsubsection{绿色虚线·部分}

\subsubsection{0自6100低点涨至7185点，共涨1085点。 ②自6629低点+1085点=等幅计算的涨幅满足7714点。}

148

\clearpage

进阶篇

台股同年10月中旬涨至涨幅满足附近（最高7811点）便做明显拉回，但仍 维持在上升趋势线支撑之上。同年11月17日再过前高7811点，便能计算涨幅满 足。

观察重点二：

绿色虚线2部分 0自6629低点涨至7811高点，共涨1182点。 ②自7218低点+1182点=等幅计算的涨幅满足8400点。

台股隔年2010年1月达涨幅满足附近（最高8395点），盘头大幅拉回。

这次1月21日破线短头确立，及1月26日跌破上升趋势线，首度出现结束波 段多头的危机，此时波段多单通常先卖再说，避开进入空头的风险。指数在破线 后约半个月跌约千点，直到3月8日站回线上出现高档的破底翻，才扭转空头颓 势。不过在角度较高的地方做破底翻，虽然也是企稳信号，但相较低档破底翻的 风险高出很多，因此在多单的重新介入除了慎选个股外（以先转强的类股为主， 原则上量价不表态不介入），资金控管少量操作，尽量降低容错率。

在时间对称的结构下当1月26日破线转弱后，表示往前2009年9月上旬至12 月中旬约三个半月，已成为相对位置的套牢区（紫色框区块），包括头部区共约 四个半月的套牢区，理论上扭转结构也需要对称以上的时间，以盘代跌整理消化 卖压。2010年3月8日拉回此套牢区，高档破底翻企稳确立后，此套牢区共震荡 整理了约七个月，才再度盘涨进入第二大波多头的延续。

149

\clearpage

多空转折一手抓

·操作方法：

本篇可做上一篇多头操作的延续参考，主要的重点在于1月21日破线短头确 立空单进场，约可空在8175点附近，突破前高8395点停损，风险约2.7\%。

1月26日跌破上升趋势线空单加码，约可空在7820点附近，突破8015点也 就是站回线上时，停损利风险约2.5\%，但在3月8日高档破底翻扭转空头趋势 后，约在7760点附近空单需退出观望。

小叮： ①本书以成功的操作型态为主，但实际的操作上严格执行停损才是成功的重 要关键。 ②失败的双底：

突破颈线 买进信号

跌落颈线下 停损

150

\clearpage

\section{2}

进阶篇

③失败的收敛三角形底部

突破颈线 买进信号

跌落颈线下 停损

④失败的头肩顶

站回颈线停损

跌破颈线 放空信号

151

\clearpage

多空转折一手抓

\subsubsection{（二十二）2009-2011台股周线}

0000大后款（调

\section{体指带量跌破颈线9230.55}

\section{涨幅满足M头确立}

\subsection{863895}

\subsection{8170.72)}

\subsection{7787 8057.21}

\subsubsection{9/18突破颈线7475.46}

7084

\subsection{703240量放过速6883.86}

突破整理区颈线

\subsection{6302.12}

\section{5720-3R. 5621.78 5128.77 带量突破颈线底部确立 4547.03}

\section{三角形收敛型态3955.43}

3955.43

10593377

\section{8522491}

\subsubsection{→重点提示：涨幅满足、M头、量放过速。}

2009年3月14日带量突破颈线，底部确立为技术面最佳加码时机，波段大涨

\section{至同年6月间拉回整理，同年9月上旬突破整理区颈线越过前高，便可计算第二}

大波涨幅满足。9月18日当周突破颈线为买进信号，跌破颈线停损。

\subsubsection{观察重点：}

0自3955低点涨至7084点，共涨3129点。 ②自6100低点+3129点=等幅计算的涨幅满足9229点。

\section{台股盘涨至2011年2月中旬达涨幅满足附近（最高9220点），便进入长期}

\subsection{盘头。}

152

\clearpage

进阶篇

盘头期进行到同年8月6日，周线带量长黑损破颈线，M型头部确立，M头 的时间规模达约十一个月，盘跌的过程至隔年2012年2月4日，周线带量长红， 突破颈线，出现底部型态，但底部的时间规模五个多月，远低于M头规模，且 放量过速，如当时部落格预测主力求急不求长久，看短不看长，属于M头长空 型态的反弹逃命波，台股随后盘头跌回起涨点。

?操作方式：

2009年3月14日当周大量突破底部确立多单进场，约可买在4780点附近停 损设带量起涨低点4688点风险约2\%，达涨幅满足附近，获利了结约88\%。就型 态来看，若在2010初的震荡被中途洗出多单，也有约58\%的波段多单获利。

小叮：

“2008-2009台股日线”为同时期同样的三角形收敛底部操作，日线操作较 即时，因此型态上有转折机会时，可观察较短K线，甚至是小时线或五分钟线的 量价结构，以求更即时积极的介入时机。

大师语录

“使我犯错的是我没有足够的毅力按计划做，即只有在先满

足入场条件时才入场。对每天都要买卖的人来说，他不可能有足 够的理由和知识，使他每天的买卖都是理性的。”

杰西·利佛摩尔（JesseL.Livermore）

153

\clearpage

多空转折一手抓

\subsubsection{（二十三）2011-2012台股周线}

0000大照拍数（） 12584商71396证99897

\section{9220.699180}

\subsection{M头9225.58}

893215

\subsubsection{2011/8/6 8643.60}

\section{8170.72 4/7 8355.05}

\subsection{8245跌破颈线-8061.62}

\subsubsection{[8070.52]头部确立}

\subsection{7773.07}

\subsection{7484.53}

量放过速7191.09

\section{7032.40] 2/4 6922736902.55}

跌幅满足6877.12 6857.35 6609.11] 6609.11

\subsubsection{9114亿91827012}

7332455

\section{0.000}

\subsubsection{重点提示：M头、跌幅满足、量放过速}

\subsection{观察重点一：}

\section{2011年空头的结构8月跌破颈线后可以M头的型态计算跌幅满足}

0颈线约在8245点附近至双头连线的位置，约在9180点距离约935点。 ②跌破颈线的位置约在8490点-935点=等幅计算的跌幅满足7555点。

台股跌至第一波跌幅满足附近，大幅震荡后进入第二波跌势。

154

\clearpage

进阶篇

观察重点二：

第二波跌幅满足计算为第一波跌幅满足点-935点=6620点。 台股达第二波跌幅满足附近（最低6609点），配合日线出现破底翻后，便 应翻空为多。

台股破底翻进行反弹波多头至2012年2月4日，周线带量长红，突破颈线， 型态上突破颈线后，明显6877点及6609点双底成形也就是W底，不过这次却是 量放过速，属于主力求急不求长久、看短不看长的结构。

大盘周线长红，突破颈线型态上约五个半月的底部成型，但相对上档8200 点以上约10个半月的大头部区，规模差异仍大。因此未达相对时间规模的整理 之前，应以反弹视之。

操作方法：

0量价结构上量放过速表示主力求急，而非大多头逐步放量求长久结构，表 示业内作价看短不看长，但是量虽放过速也表示业内成功作量吸引散户进 场（融资增加）。通常都还会有一段类股轮动的盘头或盘涨期，以利其边 拉边跑的动作。直到同年4月7日当周跌破颈线，完成短期头部结束反弹 波，指数大幅回跌至起涨区，结果如当时部落格预测的“量放进速，主力 求急而非求长久” ②部落格网址：ts888.blogspot.tw，参考2012年2月7日PO文。

155

\clearpage

多空转折一手抓

\section{MEMO}

156

\clearpage

进阶篇

二、个股 先了解大盘的趋势，再介入个股操作，将会大幅提高成功率。因此随势选股 非常重要。当大盘脱离底部进入大波段起涨时，尤其是第一波段，几乎是鸡犬升 天、涨多涨少的问题而已。从经验上看，起涨第一波，跌深的低价转机股通常涨 幅最为惊人。大盘若有第二波以上时，选股难度便随之上升，技术面量价结构产 生的型态，便提供了相当大的帮助。

反之，进入空头的结构，状况大多是覆巢之下无完卵的。一般头部的完成， 多半是在市场最无戒心的时候。这方面，技术面量价结构提供相当重要的判断， 能率先市场抢得先机卖出股票，甚至反手放空。无论是大盘或个股的多空型态， 前半部技术篇介绍的12种方法都足以提供判断。当然学无止境，不过切记，无 论型态如何变化都还是建立在量价结构的基础上。另外，个股的多方趋势或空方 趋势操作，至少要建立三至五档组合，切忌单压个股，有较大把握的趋势成形 时，按个人的风险承受能力，先投入五成以上资金并严设停损。趋势若如预期， 个股扩大获利时，再加码轮动股即可，届时即使是多单或空单满档，都是坐在波 段趋势的轿子上。

大师语录 “相对于赚钱，我比较注意赔钱。别光想赚钱，要专注于保 护你的战果。”

保罗·都铎·琼斯（PaulTudorJones）

157

\clearpage

多空转折一手抓

\subsubsection{（一）1997-1998华泰日线}

2329禁鑫（旧绿餐） 3

\subsubsection{50.29头肩顶50.29}

\subsection{47.12}

\subsubsection{涨幅满足43.94}

\subsection{M头}

\section{颈线40.6 40.72}

37.38] 37.54

\subsubsection{上飘旗形36.15 4/28}

\subsection{34.37}

10/17 34.8

\section{颈线32 31.14}

颈线292922

1/20 27.97 跌幅满足

\section{21.57跌幅满足212 21.70 21.85 22.02W底21.57}

\section{45501 34019 22963 11481 0.000}

→重点提示：头肩顶、跌幅满足、上飘旗形、W底、涨幅满足M头。

\section{华泰日线1997年9月1日跌破颈线，三尊头型态的头肩顶确立便可计算跌幅 满足。}

\subsection{观察重点一：}

0头部高点50.29元至颈线的位置约40.6元，距离9.69元。 ②跌破颈线的位置40.6元（颈线为平行线）-9.69元=等幅计算的跌幅满足 30.91元。

华泰跌至跌幅满足附近（最低30.81元）反弹，同年10月17日长黑破线 上飘旗形确立，隔日再破30.81元前低，便可再计算跌幅满足。

158

\clearpage

进阶篇

观察重点二：

第一波跌幅满足30.91元-9.69元=等幅计算的跌幅满足21.22元。

华泰达跌幅满足附近（最低21.7元）后企稳筑底，打底至隔年1998年1 月20日突破颈线，W底确立，便可计算涨幅满足。

观察重点三：

0双底低点中间值约21.85元至颈线位置约29元，距离7.15元。 ②突破颈线的位置29元（颈线为平行线）+7.15元=36.15元。

华泰约三个月的底部突破约一周，便达涨幅满足附近，就规模而言，可 直接计算第二波涨幅满足。

观察重点四：

第一波涨幅满足36.15元+7.15元=等幅计算的第二波涨幅满足43.3元。

华泰达第二波涨幅满足附近后进入大区间震荡盘头（2月中旬最高42.59 元，4月上旬最高45.76元），同年4月28日跌破颈线，M头确立，便可计算 跌幅满足。

159

\clearpage

多空转折一手抓

观察重点五：

0双头高点中间值约44元，至颈线位置约33.2元，距离10.8元。 ②跌破颈线的位置约34.8元-10.8元=等幅计算的跌幅满足24元。

华泰跌至同年6月达跌幅满足附近企稳。

操作方式：

01997年9月1日跌破头肩顶颈线时，空单进场，约可空在40元附近，停损 设在42元附近，也就是站上颈线反压，风险约5\%，反压站不回，看跌幅 满足附近，空单回补获利，第一波跌幅满足附近获利约22\%，第二波跌幅 满足获利约47\% ②1998年1月20日突破W底颈线，多单进场，约可买在29.5元，停损设在 28元附近，也就是跌破颈线支撑，也是当日长红低点风险约5\%。支撑不 破，看涨幅满足附近，多单获利了结，第一波涨幅满足附近获利约22\%， 第二波涨幅满足附近获利约42\% ③1998年4月28日跌破M头颈线时，空单进场，约可空在34元附近，停损 设36元附近，也就是站上颈线，反压风险约6\%，反压站不回，看跌幅满 足附近，获利约30\%。

160

\clearpage

\subsection{进阶篇2}

小叮：

有许多型态，如华泰的头肩顶型态，除了计算头肩顶型态的跌幅满足外，也 可用波段跌幅满足做计算，大多会落在大约的位置上，只是如华泰有相当完整的 头肩顶型态，便会以头肩顶为主计算跌幅满足。

实线=头肩顶两波段涨幅满足 虚线=波段跌幅满足

大师语录 “大钱不存在于股票的日常小波动，大钱只存在大势之内。 因此你需要判定大势的走向。

杰西·利佛摩尔（JesseL.Livermore）

161

\clearpage

多空转折一手抓

\subsection{（二）1997-1998联电日线}

2303電（日） 17.76高18.10 17.57 森指槽40.1940.19

\subsection{假突破—8/43g35头肩顶40.24}

37.18

\section{涨幅满足34.72 34.18}

\section{-31.17}

\subsection{28.17 25.67}

\subsection{25.11}

\subsection{20.92 20,92破底翻2211}

跌幅满足 19.11

\subsection{16.63 16.1016.07}

1997/5/29带量突破15.94 15.52

.05 13.26 13.05

\subsection{244995}

\section{146997 97998 48999 0.000}

\section{→重点提示：涨幅满足、假突破、头肩顶、跌幅满足、破底翻。}

\section{联电日线1997年首波涨至6月下旬拉回，7月初再创新高时便可计算涨幅满}

足。

\subsection{观察重点一：}

\subsubsection{0自13.05元低点涨至25.67元高点，共涨12.62元。 ②自22.1元低点+12.62元=等幅计算的波段涨幅满足34.72元。}

\section{联电涨至7月中旬，达涨幅满足附近，拉回后进行大区间盘头。8月4日}

\section{创40.19元历史新高，隔日便拉回突破的颈线以下，技术面假突破确立长线卖}

\section{出信号。9月2日跳空长黑跌破颈线，头肩顶确立，便可计算跌幅满足。（注}

162

\clearpage

进阶篇

意：破线隔日急弹未突破颈线为技术面的逃命线。）

观察重点二：

0自头部高点40.19元至颈线位置约31.3元距离8.89元。 ②自跌破颈线位置约34元-8.89元=等幅计算的跌幅满足25.11元。

联电9月15日达跌幅满足附近，短弹几日后回跌再破前低，便可计算第 二波跌幅满足。自第一波跌幅满足25.11元一8.89元=等幅计算的第二波跌幅 满足16.22元。

观察重点三：

联电盘跌至同年10月30日，达跌幅满足附近（最低15.94元），才企稳筑 底，盘底至隔年2月底才以破底翻确立底部，大幅反弹。

操作方式：

01997年5月29日带量突破颈线，多单进场，约可买在16元附进，停损设 15元也就是跌落线下时，风险约6.2\%，达涨幅满足获利约117\%。 ②跌破颈线时，空单进场，约可空在33元附近，停损设在35元附近也就是 站上颈线，反压风险约6\%。反压站不回，看跌幅满足，空单回补获利， 第一波跌幅满足获利约24\%，第二波跌幅满足获利约52\%。

163

\clearpage

多空转折一手抓

\subsubsection{（三）1998-2000友讯周线}

2332（道） 17-7高176016.00收17.07量25514

\subsubsection{②回体森指量-44.24}

\subsection{44.24}

\section{涨幅满足43.86头肩顶}

40.44

8590

\subsection{3272}

28.93

25.20

\section{23.6颈线2293 26.324.46 28.5 25.06}

1/8 2.60]

\subsection{21.4破底翻21.20}

17.41

\subsection{13.55}

\section{头肩底跌幅满足}

\subsubsection{9.6898.091}

\section{152601 114814 77027 39240 0.000}

\section{→重点提示：头肩底、涨幅满足、头肩顶、跌幅满足、破底翻}

\section{如上图，2000年1月8日友讯周线当周带量长红，突破颈线，约一年半的大}

\subsubsection{头肩底确立，便可计算涨幅满足。}

\subsection{观察重点一：}

0底部低点12.37元至颈线位置约23.6元，距离11.23元。 ②突破颈线的位置约21.4元+11.23元=等幅计算的涨幅满足32.63元。

\section{友讯同年3月中旬达涨幅满足附近后，周线一周大幅拉回逾15\%，隔周}

再大幅回升约25\%，充分反映当时网络泡沫投机的气氛，技术面越过前高便 可再计算第二波涨幅满足。

164

\clearpage

进阶篇

观察重点二：

第一波涨幅满足32.63元+11.23元=等幅计算的第二波涨幅满足43.86元。

友讯4月8日当周达第二波涨幅满足附近（最高44.24元），V型反转， 股价急转直下，反映当时友讯搭上网络泡沫的末班车后，随着泡沫戳破，股 价怎么上去怎么下来的投机走势。同年7月1日当周，跌破颈线，头肩顶确 立，便可计算跌幅满足。

观察重点三：

0自头部高点44.24元至颈线位置约在26.3元，距离17.94元。 2自跌破颈线位置约在28.5元-17.94元=等幅计算的跌幅满足10.56元。

友讯长期盘跌至隔年2001年1月初，达跌幅满足附近（最低9.69元）才企稳 破底翻，大幅反弹。

操作方式：

①突破颈线时买进，约可买在21.5元，停损设20.5元，也就是跌落线下时， 风险约5\%，达涨幅满足获利了结，约51\%～104\% ②跌破颈线时空单进场约可空在28.3元附近，停损设30元附近，也就是站 上颈线，反压风险约6\%，反压站不回，看跌幅满足附近，空单回补获利 约60\%

165

\clearpage

多空转折一手抓

\subsubsection{（四）1998-2001茂矽周线}

2342动（）期575.70两584.90低541收962.05 强肤

\subsubsection{森指854,50854.50 848}

\subsection{复合式头肩顶 涨幅满足800 70.19}

\subsection{626.55 9/9}

\section{535 551.44}

\section{475.02}

[420,23]

\subsection{375.62带量破底翻398.60}

\section{296.81 323.49}

9/4带量突破颈线247.07

\section{171.81收敛三角形底部187.35跌幅满足193 170.54 170.64}

\section{710743 482941 246026 0.000}

→重点提示：收敛三角形底部、涨幅满足、复合式头肩顶、跌幅满足、破底

\subsubsection{翻。}

茂矽周线（图示为茂矽几次减资后的还原权值图）1999年9月4日带量突破 颈线，逾一年半的底部确立便可计算涨幅满足。

\subsubsection{观察重点一：}

01999年8月7日波段最低点187.35元=当时价位20.5元。 22000年4月8日波段最高点854.5元=当时价位93.5元。 32001年1月6日波段最低点170.64元=当时价位16.5元。

4三角形边长的高点420元至低点约在165元，距离255元

③突破颈线的位置约在290元+255元=等幅计算的涨幅满足545元。

166

\clearpage

进阶篇

茂矽盘涨至隔年2000年2月，达涨幅满足附近，震荡整理，整理后再度 突破整理区高点，便能计算第二波涨幅满足。

观察重点二：

第一波涨幅满足545元+255元=等幅满足计算的第二波涨幅满足800元。

茂矽同年4月初达第二波涨幅满足附近，便进入大区间震荡盘头，先盘 出M头后，再盘出复合式头肩顶。同年9月9日，长黑跌破颈线，复合式头 肩顶确立，便可计算跌幅满足。

观察重点三：

0头部高点中间均值约848元至颈线位置约在506元，距离342元。 ②跌破颈线的位置约在535元-342元=等幅计算的跌幅满足193元。

茂矽跌至同年底达跌幅满足附近，仅震荡2周便带量翻红，破底翻企稳确 立

操作方式：

01999年9月4日当周突破收敛三角形底部型态颈线时，多单进场，约可买 在295元附近，停损设267元，也就是跌破颈线支撑，也是当周大量长红 低点，风险约9\%。支撑不破，看涨幅满足附近，多单获利了结，第一波 涨幅满足附近获利约84\%，第二波涨幅满足附近获利约170\%。 22000年9月9日当周，跌破复合式头肩顶颈线时，空单进场，约可空在 530元附近，停损设在563元，也就是站回颈线反压，也是当周长黑高点， 风险约6\%；反压站不回，看跌幅满足附近，空单回补，获利约63\%。

167

\clearpage

多空转折一手抓

\subsubsection{（五）1998-2003台积电周线}

2330台费增（运）打25高品收托1 6\#746 26（4.42\%）日期2000/3/1

\subsubsection{2体指成交51.42旅量锦量}

\section{51.42}

2000/7/21长黑跌破 颈线收敛三角形头部确立46.93

42,11 42.35

\section{42.5逃命线2002/6/8长黑跌破}

\subsection{39.5颈线头肩顶头部确立37.62}

36

\subsection{34.9 33.29}

\subsection{25.79突破28.71}

\subsection{破线24.22}

22.45

\subsection{18.83 119.15 19.64}

\subsection{15.07}

收敛三角形底部

\section{破底1998/11/21带量10.50}

10.5010.50翻红站回颈线破底翻 文量518514 414811 315097 207406 103703 0.000

\section{→重点提示：破底翻、涨幅满足、收敛三角形头部、头肩顶、跌幅满足、收}

\subsubsection{敛三角形底部。}

图示为台积电1998年至2003年周线k线型态（填权息），1998年8月29日 当周破底，同年11月21日当周，出现带量长红破底翻，为底部企稳确立的波段 买进信号。隔年1999年6月26日当周，见首波高点30.5元，拉回最低23.4元，当 股价回升，至同年9月4日再过前高30.5元，便能计算波段涨幅满足。

\subsection{观察重点一：}

\subsubsection{0自最低点10.5元至30.5元高点，共涨20元 ②自23.4元低点+20元=等幅计算的涨幅满足43.4元。}

168

\clearpage

进阶篇

台积电盘涨至隔年2000年1月下旬，达涨幅满足附近后，虽然持续超涨 （最高51元），但型态上仍是以大波段涨幅满足为中心，做高档大区间震 荡，长达约半年。同年7月21日当周，长黑损破颈线，收敛三角形头部确 立，破颈线的位置约在42.5元附近，这是对应1988年11月21日当周，站回 颈线破底翻约在12.8元附近以来的多头结束确立，也就是技术面波段操作会 买在破底翻底部确立的13元附近。卖在跌破颈线头部确立后的42元附近，时 间约一年八个月，波段获利约220\%

观察重点二：

①当收敛三角形颈线跌破头部确立后，便可计算跌幅满足。 计算方式：三角形边长=跌破颈线后的等长。 ②三角形边长高点约51元-颈线位置约39.5元=11.5元。 ③跌破颈线的位置约42.5元-11.5元=等幅计算的跌幅满足31元。

台积电跌至首波跌幅满足附近后，以满足点为中心，做大区间震荡约两 个月（最低24.56元，最高36.91元）后再度破底。

观察重点三：

第二波跌幅满足计算： 第一波跌幅足点31元再减11.5元=等幅计算的跌幅满足19.5元。

台积电2001年9月跌破长期收敛整理型态下缘颈线3周后，达第二波跌 幅满足附近（最低18.8元），然后直接V形反转回升，先站回颈线，破底翻 后持续大涨至最高42.11元才反转。

169

\clearpage

多空转折一手抓

观察重点四：

2002年6月8日，长黑跌破颈线头，肩顶型态的头部确立，便可计算跌幅满 足： 0头部高点42.1元至颈线位置约在34.9元，距离7.2元。 ②跌破颈线的位置约在36元-7.2元=等幅计算的跌幅满足28.8元。

台积电首波跌至同年7月6日当周最低29.37元，反弹至最高34.14元，完 成逃命线后再度破底。

观察重点五：

第二波跌幅满足计算： 自第一波跌幅满足点28.8元-7.2元=等幅计算的跌幅满足21.6元。

台积电达第二波跌幅满足后，以满足点为中心，作区间打底达，约八个月 （最低16.67元最高25.79元）。2003年5月31日当周长红突破整理区上缘颈线， 完成收敛三角形底部，进入一波多头行情。

操作方式：

01998年11月21日当周带量长红破底翻时，多单进场，约可买在13.3元， 停损设在12.5元，也是当周大量长红低点支撑，风险约6\%，支撑不破， 看涨幅满足附近，多单获利约226\%。 22002年7月21日当周长黑损破颈线时，空单进场，约可空在42元附近， 停损设在44元附近，也就是站上颈线反压，也是长黑高点，风险约5\%。 反压站不回，看跌幅满足附近，空单获利了结，第一波跌幅满足附近获利 约26\%，第二波跌幅满足获利约53\%

170

\clearpage

进阶篇

32002年6月8日当周长黑损破颈线时，空单进场，约可空在34.5元附近， 停损设在36元附近，也就是站上颈线反压，风险约4.3\%。反压站不回， 看跌幅满足附近，空单获利了结，第一波跌幅满足获利约15\%，第二波跌 幅满足附近获利约36\%

小叮： 由“1998-2003台积电周线”的例子来看，约五年的时间里，只有两次明显 突破的波段买点（2001年达两波段跌幅满足后的破底翻或可加算一次），及两次 跌破颈线完成头部的波段卖点（空点）。操作波段要有耐心，机会是等出来的。

\section{MEMO}

跌破颈线后反弹未站回颈线称之为逃命线（粉色箭头部分）

171

\clearpage

多空转折一手抓

（六）1999-2002金宝周线

2312金（银） 22.10 23.03借71.25收22.54 667500.160.709\%）日期2000/3/11

\subsubsection{保森指槽26.73成交报量硬量}

\subsubsection{26.73] 26.73}

\subsubsection{江假突破}

\subsection{头肩顶24.44}

22.10

\subsection{19.59 19.77}

7/29跌破颈线 16.96] 17.48

\section{涨幅满足02}

\subsection{14.9716.7 15.15}

\subsubsection{14.19}

12.82

\subsection{10.53}

10.36]

9.34破底翻12/8 8.198

\section{跌幅满足5.92带量突破颈线5.867}

\section{285572 214179 142786 71393 0.000}

\section{→重点提示：假突破、头肩顶、跌幅满足、头部与底部的规模、收敛三角涨}

\subsection{幅满足。}

\section{金宝2000年2月突破上缘颈线又跌回破线，假突破确立长线卖出信号，股价}

\section{盘跌至同年7月29日当周，跌破颈线，头肩顶确立，便能计算跌幅满足。}

\subsubsection{观察重点一：}

\section{0自头部高点26.73元垂直至颈线位置，约在15.7元，距离11.03元。 ②自跌破颈线位置约在16.7元-11.03元=等幅计算的跌幅满足5.67元。}

金宝盘跌至隔年2001年初，达跌幅满足附近（最低6.08元）后破底翻急弹， 这次破底翻的底部规模两个多月，远小于头肩顶十个多月的头部规模，通常短线

172

\clearpage

进阶篇

反弹视之（但因跌深弹幅也达倍数），3周急弹后进入长期盘跌，至同年10月27 日当周，带量长红，突破颈线，底部确立，此时对应年初低点底部规模已达十个 多月，较易有大波段多头走势，12月8日带量长红，突破大收敛三角形颈线，便 可计算涨幅满足。

?观察重点二：

0自6.08元低点垂直至上缘颈线位置，约在12.8元，距离6.72元。 ②自突破颈线的位置约在9.8元+6.72元=等幅计算的涨幅满足16.52元。 金宝波段大涨至隔年2002年4月，达涨幅满足附近（最高16.96元）才盘头 翻空。

操作方式：

①空单操作 2000年2月21日假突破确立后，空单进场，约可空在22.9元附近，停损设在 站上颈线24.4元时，风险约6.5\%，同年7月29日当周，跌破颈线确立为空单加码 点，约可空在16元附近，站回颈线16.7元之上，停损利点风险约5\%（假突破高 档空单停利，破线加码空单停损），颈线反压站不回，看跌幅满足附近，空单回 补获利，获利约50\%以上。

②多单操作 2001年10月27日当周带量突破颈线后，多单进场，约可买在7.5元附近，风 险设约7\%，也就是跌破颈线时，同年12月8日当周再度带量突破大颈线后，为 多单加码信号，约可买在9.7元附近，风险设约7\%，停损也就是跌破颈线时，颈 线支撑不破，看涨幅满足附近，获利约70\%～120\%。

173

\clearpage

多空转折一手抓

\subsubsection{（七）2002-2004友达周线}

2409发建（更量） 1435商14.76倍13.04收14.75

\subsubsection{2回体森指播49.2149.21}

49.28

\subsubsection{涨幅满足44.77 44.89}

\subsubsection{6/12 40.49}

36

\subsection{36.17}

\subsection{31.78W}

27:38 跌幅满足

\subsection{24.74 23.06}

22.19

\subsection{15.40] 14.71] 1428}

\section{89'01 5/31破底翻9.880}

\section{1061018 798290}

\section{0.000}

\subsubsection{重点提示：破底翻、涨幅满足、头肩顶、跌幅满足。}

\section{友达周线2002年5月31日出现带量破底翻买进信号，股价就此数倍波段大}

\section{涨，第一大波涨至30.88元拉回修正，2004年2月21长红突破前高，便能计算涨}

\subsection{幅满足。}

\subsection{观察重点一：}

\subsubsection{0自9.88元低点涨至30.88元高点，共涨21元。 ②自23.77元低点+21元=等幅计算的涨幅满足44.77元。}

\section{友达同年4月中旬达涨幅满足附近（最高49.21元），不到3周便急速反}

\section{转盘头，同年6月12日跌破颈线头肩顶确立，便可计算跌幅满足。}

174

\clearpage

进阶篇

观察重点二：

0自头部高点49.21元至颈线位置，约在35.5元，距离13.71元。 ②自跌破颈线位置约36元-13.71元=等幅计算的跌幅满足22.29元。

友达盘跌至同年11月6日达跌幅满足附近（最低22.19元），企稳打底后回 升。

操作方法：

0多单操作： 友达日线2003年5月26日破底翻确立，隔日买进约在11.6元附近，停损设 破底前低点10.68元跌破后（破底前的区间才是主力成本区），风险约8\%。当周 周末5月31日收盘长红周线破底翻，达周线涨幅满足44.77元附近的波段利润约 286\%，利润是风险的35倍。唯有一次日线11月17日破线，短线转弱，可能卖出 的位置在26.5元附近，波段利润也高达约128\%

②空单操作： 友达周线2004年6月12日当周跌破36元颈线，便是空单进场布空时。破线 后放空的位置约在35元附近，36元颈线跌破支撑变反压，空头看反压因此站回 颈线反压之上停损，可设约5\%停损。反压站不回看头肩顶波段跌幅满足22.29元 附近空单回补，波段利润约36\%

175

\clearpage

多空转折一手抓

\subsection{（八）2003-2008光磊周线}

2340光磊（） 1429高14.96低13.1收14.37 4751呈7745肤7231\%） 爱贸贵惠

\subsection{43.42}

\subsubsection{涨幅满足3922}

37.74

\subsection{31.11}

\section{28.11}

\subsection{22.18}

\subsection{17.97 17.13颈线16.2 15.92 15.5 17.91}

\section{4/15 15.14 [16.33 13.63}

\subsection{[.82] [10.98] 9.354}

5.078

\section{保森指308003}

\section{186234 124156 62078 0.000}

\subsubsection{→重点提示：底部、涨幅满足、头肩顶、跌幅满足。}

\section{光磊周线2006年4月15日当周突破颈线，带量长红，三年四个多月的底部}

\section{确立，便可计算涨幅满足，只是突破颈线后的回测颈线时间为少见的约半年之}

\subsection{久。}

\subsection{观察重点一：}

\section{0自底部低点5.08元至颈线位置，约在16.2元，距离11.12元。 ②突破颈线的位置约在15.5元+11.12元=等幅计算的涨幅满足26.62元。}

\section{光磊2007年初，达涨幅满足附近整理达三个多月，同年3月4日当周，突破}

\subsubsection{旗形整理上缘，便可计算第二波涨幅满足。}

176

\clearpage

进阶篇

第一波涨幅满足26.62元+11.12元=等幅计算的第二波涨幅满足37.74元。

光磊同年7月底，达第二波涨幅满足附近后进入盘头，11月3日当周长黑跌 破颈线，头肩顶确立，便可计算跌幅满足。

观察重点二：

0自头部高点43.42元至颈线位置约在29.1元，距离14.32元。 ②自破线位置约在31元-14.32元=等幅计算的跌幅满足16.68元

光磊盘跌至隔年2月中旬，达跌幅满足附近企稳反弹。

操作方式：

02006年4月15日当周带量突破颈线买进，约可买在16元附近，跌破颈线 约14.8元停损，风险约8\%，达两波段涨幅满足26.62元及37.74元附近， 获利65\%\textasciitilde{}135\%。 ②2007年11月3日当周长黑跌破颈线空单进场，约可空在30.5元，停损设 32元也就是站回线上时，风险约5\%，达跌幅满足空单回补获利约45\%。

大师语录 “别想买到最低、卖到最高。

伯纳德·巴鲁克（BernardBaruch）

177

\clearpage

\subsubsection{多空转折一手抓}

\subsection{（九）2003-2013宏基日线}

3353宝基（思强）高收46.44 6318第2239380.551.389\%日期2006/3/1

\subsection{森指楼要成交量热量94.3}

\subsection{92.16M头-94.35}

-85.08

\subsection{颈线区75.81}

\subsection{58.04带量压低出货66.53}

\subsection{51.08 57.26}

\subsection{48.21}

\subsubsection{高档无量}

\subsection{[37.52 38.94}

\subsection{31.77 29.67}

\subsection{20.39}

\subsection{主力进货区11.12}

\section{265544}

0.000

\subsection{价跌资增118413}

\section{散户一路套7851339900}

\subsection{262567 171297}

\subsection{主力出贷80027}

\subsection{主力进货-11242}

\section{2005 2010 2011 2013 -106315}

\subsection{重点提示：高档压低出货}

\section{自从2008年高油价产生的全球金融风暴崩盘至当年底见低点后，庞大的印}

\section{钞救市，以及不断的量化宽松政策造成泡沫式的资金多头行情。台股亦至当年低}

\section{点弹升，个股涨幅甚至以数倍计，因此许多个股高涨后，常见主力在时间与空间}

\subsection{的有利条件下，高档做压低出货的方式（尤其是外资主力）。}

\section{观察重点：}

\section{宏碁自2008年11月21日低点31.77元（填权息）波段大涨，涨幅高达约}

\subsection{200\%，后修正再大区间盘头。}

\section{2011年初库存余额显示主力在高档积极落跑压低出货，融资散户则以为短线}

\subsubsection{178}

\clearpage

进阶篇

急跌逢低买进，股价一路盘跌惨遭套牢，可见基本面里，苹果颠覆市场的IPad对 平板电脑及NB市场的冲击相当大。

如图示，2003年上半年宏碁在14元附近区间震荡打底，随后一路盘涨，拉 高进货，至2005年5月7日带量长红，突破颈线，发动攻势，主力进货区的平 均成本粗估约在26元附近或以下。筹码集中在法人手中。于2009年至2010年拉 高至65元附近至90余元，高档震荡无量。2010年底自高档92.16元压回至颈线 附近，很明显出现价跌资增、主力库存减的高档压低出货结构（咖啡色虚圈部 分），股价在颈线之上震荡盘弱约一个月后，只做短暂3天反弹便跳空崩跌，可 见高档“颈线支撑区是主力法人压低出货的最大陷阱”。

?操作方法：

误判颈线支撑低接的散户在跌破支撑后，噩梦才刚开始，后续虚圈部分宏碁 爆量下跌，融资急速增加且主力库存更快速出货，属带量压低出货，股价一路盘 跌下来连环套，越来越多无知的散户套牢。主力约26元以下的原始筹码随便卖 随便赚。其实因误判而在颈线支撑买进并不可怕，可怕的是跌破支撑不知严格执 行停损，因为停损后的损失有限，不知停损则如宏基跌破支撑一年的大M头成 形股价腰斩再腰斩，操作上宁可做错也不要留错，做错还能小赔大赚，留错则翻 身困难。

更多详细内容请参考2011年3月23日部落格P0文。

179

\clearpage

多空转折一手抓

\subsubsection{（十）2004-2012台达电周线}

2308台建章（运）高47.70售44.50收5747

要成交195 126.73

\subsection{126.73假突破126.73}

\subsubsection{涨幅满足114.74}

2/12 98.81 102.54

\subsection{2010/7/10}

\section{87.87破底翻90.34}

78.34

\subsubsection{跌66.14}

\subsection{满足58.4553.94}

-41.95 29.75 25.75

\subsection{17.55}

成文量128939 103348 77757 52166 26575 0.000

\subsubsection{重点提示：涨幅满足、假突破、跌幅满足、破底翻。}

\section{台达电周线2010年7月10日当周长红，突破前高95.33元后，便可计算大波}

\subsubsection{段涨幅满足。}

\subsection{观察重点一：}

0自17.55元低点涨至95.33元高点，共涨77.78元。

②自40.29元低点+77.78元=等幅满足计算的涨幅满足118.07元。

\section{台达电同年底达涨幅满足附近进入盘头（最高126.73元），至隔年2011}

\section{年2月12日当周带量长黑，确立了假突破翻空的长线卖出信号，结束两年多}

\section{头行情，进入长空修正。同年8月7日当周跌破反弹低点87.87元，便能计算}

180

\clearpage

进阶篇

跌幅满足。

·观察重点二：

0自126.73元高点跌至87.87元，共跌38.86元。 ②自98.81元高点-38.86元=等幅计算的跌幅满足59.95元。

台达电盘跌至同年11月底达跌幅满足附近，进入打底（最低58.45元），年 底12月24日当周长红破底翻，底部确立，结束约一年的空头大跌修，正进入另 一波大多头行情。

操作方式：

2011年2月12日当周带量长黑，跌破颈线，假突破确立，放空，约可空在 118元附近，停损设125元站回颈线时，风险约6\%，达跌幅满足59.95元附近，空 单回补，获利约49\%。

小叮：

这篇的型态分析，主要是确定高档假突破卖点的做空时机。

大师语录 “犯错并没有什么好羞耻的，只有知错不改才是耻辱。

乔治·索罗斯（GeorgeSoros）

181

\clearpage

多空转折一手抓

\subsubsection{（十一）2006-2008神达日线}

2315神迪（甘蛋）高30.81

\section{2指涨幅满足40.37 39.41 46.6748.14}

37.56带量假突破 -43.67 40.59

\subsection{31.38 37.58}

34.65

\subsection{31.64}

破底翻25.15 28.56

\subsubsection{跌幅满足！ 25.63}

22.62

\subsection{19.54}

\subsection{交58126 38964 19801 0.000 96075 76690}

主力库存减3328644493 -55701 68908 78582

\subsubsection{重点提示：破底翻、涨幅满足、假突破、跌幅满足。}

\section{如上图，神达日线2006年6月26日破底翻筑底后，盘涨进入一波长多走势，}

\section{隔年2007年5月28日突破拉回整理后的前高便可计算大波段涨幅满足。}

\subsubsection{观察重点一：}

0自16.53元低点涨至31.75元，共涨15.22元 ②自25.15元低点+15.22元=等幅计算的涨幅满足40.37元

\section{神达同年10月1日最高39.41元达涨幅满足附近，隔日带量跌回线下假}

突破确立长线卖出信号，结束多头行情进入价跌资增、主力库存减的长空结 构，同年12月13日跌破反弹低点便可计算跌幅满足。

182

\clearpage

进阶篇

观察重点二：

0自39.41元高点跌至27.61元低点，共跌11.8元。 ②自31.38元高点-11.8元=等幅计算的跌幅满足19.58元。

神达盘跌至隔年2008年1月下旬达跌幅满足附近，才筑短底反弹。

操作方式：

02006年6月26日破底翻底部确立，买进，约可买在19.5元附近，跌破18 元停损，风险约8\%，达涨幅满足40.37元附近，尤其在10月2日带量假突 破后，多单获利了结约88\% ②当日随之翻多为空，约可空在36.8元附近，停损设在19元附近站回颈线 时，风险约6\%，达跌幅满足附近，空单回补，获利约46\%。

大师语录 “投资法则只有两条：第一条是绝对不要赔钱，第二条是绝 对不要忘了第一条。”

沃伦·巴菲特（WarrenBuffett）

183

\clearpage

多空转折一手抓

\subsubsection{（十二）2006-2009友讯周线}

2332） 442 43.21收43.74 11680 0.53(-1.205\%） B福2007/47 量

\section{回回森指69.2268.4569.22 72.B9}

\subsection{M头 11/3 6288}

\subsection{56.42}

\section{50.08}

43.62

\subsubsection{VEZE6/21破线37.17}

3/28带量长红

29.38突破颈线30.82

跌幅满足-28双底确立 -24.37

\subsection{17,95] 17.91}

\section{13.55L}

11.46

高档带量出货81879 65628 48752 33127 16876 0.000

重点提示：M头、三波段跌幅满足、W双底型态。

友讯周线2007年11月3日当周爆量长黑，跌破颈线，M头确立。而破线后 连续两周爆量长黑，极为明显为高档压低出货结构，后续随之出现三波段的极弱 势长空盘跌型态。

\subsection{观察重点一：}

第一波跌幅满足，周线跌破颈线后，便能计算跌幅满足。 0M头两个高点的中间值约68.45元，至颈线位置约54元，距离14.45元。 ②跌破颈线的位置约56.9元-14.45元=等幅计算的跌幅满足42.45元。

\section{友讯同年12月22日当周达跌幅满足附近，反弹约2周。}

184

\clearpage

进阶篇

·观察重点二：

第一波跌幅满足42.45元-14.45元=等幅计算的第二波跌幅满足28元。

友讯隔年2008年2月23当周达第二波跌幅满足附近企稳，虽仅弱势反 弹，但整理期长达约四个月。直到同年6月21日当周破线结束整理，转弱确 立，隔周跌破前低，便能计算第三波跌幅满足。

观察重点三：

第二波跌幅满足28元-14.45元=等幅计算的第三波跌幅满足13.55元。

友讯盘跌至同年11月下旬达第三波跌幅满足附近（最低11.46元）才企稳， 整理盘底至隔年2009年3月28日，带量长红，突破颈线，约五个月的W双底确 立，结束了长达约一年半的长空型态，正式翻为多头。

操作方式：

02007年11月3日带量跌破M头颈线，空单进场约可空在56元附近，停损 设58元附近站回颈线时，风险约4\%，达第二波跌幅满足附近，空单获利 了结约45\% ②2008年6月21日当周破线，空单再度进场，约可空在32元附近，突破34.2 元站回颈线时停损，风险约7\%，达第三波跌幅满足附近，空单获利了结 约57\%

185

\clearpage

多空转折一手抓

\subsection{（十三）2007-2008鸿准周线}

2354酒津（通蛋） 72.20 73.28 7.04收银乐

\subsection{207,25 212.5 218.50] 218.50}

196.15

\subsection{12/15长黑跌破颈线173.41}

\section{M头确立}

151:06

\section{159.5 132.98}

134.79上飘旗形128.32

\subsection{105.58}

\subsection{68.87}

\subsection{6049}

\subsection{-37.75}

\subsection{15.01}

\section{95110 71333 47555 23778 0000}

\subsubsection{→重点提示：M头、跌幅满足}

\section{2007年12月15日当周鸿准周线带量跌破颈线，M头确立，便可计算跌幅满}

足。

\subsection{观察重点：}

\section{0双头高点连接的中间均值约在212.5元，至颈线约在150元，距离约62.5}

\subsubsection{元。 ②跌破颈线的位置约在159.5元-62.5元=等幅计算的跌幅满足97元。}

\section{鸿准跌至跌幅满足附近（最低86.08元）反弹，形成上飘旗形破线后，便可}

\subsubsection{计算第二波跌幅满足。}

186

\clearpage

进阶篇

1第一波跌幅满足97元-62.5元=等幅计算的第二波跌幅满足34.5元。 ②鸿准达第二波跌幅满足附近（最低34.59元）便筑底回升。

操作方式：

2007年12月15日当周带量跌破颈线，M头确立，空单进场，约可空在152.5 元（跌破前低152.84元时确定破线），停损设162元附近站回颈线时，风险约 6.2\%，达两波段跌幅满足附近，空单获利了结约70\%（第一波跌幅满足则有约 36\%)。

小叮：

上飘旗形完成的两波段跌势，也可以波段跌幅满足做计算。 0上飘旗形反弹高点132.98元-73.94元低点=59.04元。 ②反弹高点97.66元-59.04元=等幅计算的跌幅满足38.62元，也是落在大约 的底部位置下缘处（底部区高低点34.59～54.9元）。

大师语录 “投资者分两类：固执的和犹豫的。胜利者是固执的人。

安德烈·科斯托兰尼（AndréKostolany）

187

\clearpage

多空转折一手抓

\subsection{（十四）2007-2011台泥周线}

1101合泥（德）

\subsection{回45.4 46.29}

\section{44.07假突破44.86 46.29}

\section{M头42.43}

涨幅满足

\subsubsection{2008/6/14跌破颈线38.58}

34.67

\subsection{31}

30.82

\subsubsection{23.05}

\subsection{17.96] 跌幅满足15.35}

10.79

\section{134647 90887 45443 0.000}

\subsubsection{→重点提示：M头、假突破、跌幅满足、涨幅满足。}

\section{台泥周线2007年9月下旬拉高盘头，至2008年6月14日当周跌破颈线，M}

\section{头确立，便可计算跌幅满足（第二个头部已明显出现假突破的高档卖出信号）。}

\subsection{观察重点一：}

M头两个高点连线的中间位置约在45.4元附近，垂直相对的颈线中间位置约 在33.2元附近，45.4元-33.2元=12.2元。 跌破颈线的位置约在38.2元-12.2元=等幅计算的跌幅满足26元。

\section{台泥跌至首波跌幅满足附近（最低26.6元），短线快速反弹，弹后回跌再}

\section{破前低，便可计算第二波跌幅满足，头部的规模够大且跌势快又急，通常表示}

188

\clearpage

进阶篇

都还会有第二波跌势（除非整理的时间够长，以时间修正争取空间的修正）。

观察重点二：

第一波跌幅满足26元-12.2元=等幅满足计算的第二波跌幅满足13.8元。

台泥跌至跌幅满足附近初期，虽有超跌（最低11.44元）但立即V型反 转，随后盘涨共进行两年多的两大波多头走势，2011年4月2日当周带量突 破前高，便可计算大波段涨幅满足。

观察重点三：

0自11.44元低点至32.19元高点，共涨20.75元。 ②自20.6元低点+20.75元=等幅计算的大波段涨幅满足41.35元。

台泥同年中涨至涨幅满足附近盘头而下大幅回跌。

操作方式：

02008年6月14日当周跌破颈线，进场做空，约可空在37.5元附近，站回颈 线，停损可设在39元，风险约4\%，达两波段跌幅满足，空单获利了结约 63\% ②若以基本面来看，台泥为绩优龙头股之一，当年的金融风暴造成股价崩 跌，而在当年以前至少连续3年以上EPS都逾1.5元以上（2005年至2007 年1.97元、2.26元、2.55元逐年成长，2008年则受金融风暴影响为1.75 元）。当台泥达跌幅满足13.8元附近时，还原权息约在17元附近（最低 14.15元），且当年净值仍高达逾20元，达跌幅满足附近，便具现股买进 长期投资价值，达大波段涨幅满足，多单获利了结约200\%。

189

\clearpage

多空转折一手抓

\subsubsection{（十五）2008-2009彩晶日线}

616品（日） 森指程9.03

\subsection{5/8 9.030}

4/10 8.384 7.90涨幅满足 7.727 7.081 6.424

\subsection{5.767}

5.80量缩后

3/24带量突破支带量翻黑5.121 颈线多头加码时机 破底翻4.464

\subsection{高档}

\section{跌幅满足！不稳定量3.150}

266449

\section{→重点提示：跌幅满足、收敛底型、涨幅满足、不稳定量盘头。}

所谓随势而为，便是随大盘趋势作多空的选股顺势操作，如前文“2007一 2009台股周线”提到，2009年3月14日周线带量长红，突破颈线，三角形收敛 底部型态确立，而以日线来看，则是在当周的3月11日为确定突破（前一周3月 5日已微幅带量突破）。此时便是随势积极介入个股多单的时机，通常跌深后的 低价转机股大多弹幅可观，可做为多单布局标的。

\subsection{观察重点一：}

彩晶当时介入多单的主要因素有， 1达跌幅满足盘底约四个月。 ②虽然低价但每天成交量两万至三万张，流动性足，非常好进出（当时仍可

190

\clearpage

进阶篇

融资信用交易）。 ③大盘3月11日突破颈线的同时，彩晶也带量破底翻，确立为企稳翻多信 号。 4彩晶当时减资后净值回到票面价值附近，股价远低于净值，预估至少有一 季以上的安全期。

技术面随势破底翻后，当天收盘价4.25元，距破底翻后的支撑4元附近 符合停损约5\%～7\%的低风险，即使破前低3.75元才停损也是相对风险有 限，所以此股成为当时笔者积极布局多单时占资金最大比重的个股。

接下来3月24日带量突破颈线（4.8元附近）为多头加码时机，技术面突 破颈线后两个交易日前（3月20日）的低点4.51元已成为提高的支撑。跌破 支撑为加码多单的停损点，也是破底翻买进的多单停利点（颈线突破再拉回 跌破属假突破），随后股价喷出。接下来是等待卖出时机，在股价突破颈线 时，便能计算涨幅满足。

观察重点二

收敛底部的边长=突破颈线后的等长。 0边长低点3.15元至颈线约4.8元，距离1.65元。 ②突破颈线的位置约4.8元+1.65元=等幅计算的涨幅满足6.45元。

彩晶喷出一路不破3日低点属极强势，且转机股大多是一波到底的投机 行情。因此直接计算第二波。

191

\clearpage

多空转折一手抓

观察重点三：

第一波涨幅满足6.45元再加边长距离1.65元=等幅计算的第二波涨幅满足8.1 元。

操作方式：

03月11日带量破底翻，多单进场，约可买在4.25元，跌破4元，停损风险 约6\%，3月13日带更大的9万多张量长红，底型时的量最多5万多张，量 先价行，随大盘翻多趋势，随势积极买进多单，约可买在4.4元，停损设 出量前的短线整理区低点4.07元，风险约在7.5\%。3月24日带量突破颈 线，多单做最后加码，约可买在4.8元，停损设4.5元，也就是跌落线下 时，达两波段涨幅满足多单获利了结约68\%～92\% ②型态上来看，涨幅满足8.1元附近也是前波起跌高点压力区，为多单获利 调节时机，4月10日爆量也符合爆量后大多还有高点。4月15日也是在涨 幅满足附近，才见喷出以来的首破三日低点。高档震荡至5月期间三个带 量都是开盘拉高走低，为高档不稳定量易做头的结构，6月8日为量缩后 带量不涨且翻黑，隔日破线，头部确立为波段卖出信号，波段剩余多单需 积极获利了结退出。 3这是预估安全期内操作的转机股波段行情，转机股一旦有结束迹象，便应 见好就收，无需留恋，因为随着时间拉长，不确定性增加，风险也随之增 加。彩晶大区间震荡至隔年初2010年1月4日高点8.7元后，随着面板景 气陷入长空，盘跌至2011年12月19日最低1.2元。 4商品一年大多有一两次转折的大波段机会财，无论多头转折或空头转折， 要胆大心细，一旦转折出现，设好停损，积极介入操作，便能以有限的风 险获取难得一次的大波段利润。如彩晶3月11日的破底翻4.25元至4月中 旬涨至涨幅满足8.1元附近（最高8.63元，波段最高9.03元），仅约一个

192

\clearpage

进阶篇

月获利约90\%至逾倍数，突破颈线的加码点更仅约半个月便达涨幅满足。 股票不要经常进出短线操作，勤做功课才有实力掌握波段时机，以逸待 劳。

\section{MEMO}

193

\clearpage

多空转折一手抓 AU7

\subsubsection{（十六）2008-2011台化周线}

1326台化（）随40：13离40.89低3829收39.36量31504103511.150）期2005/6713

\section{伊森指機安票净跌、宝量：热量101,32101-32 106.14}

\subsection{涨幅满足97.4 92.99 81.17-84.51 79.35 76.18}

\subsection{-67.70}

\subsection{59.22}

\subsection{54.09] -50.89}

\subsection{44.79颈线43.5 45.8643.251}

\subsubsection{8/15突破颈线42.41}

33.93

\subsection{25.45头肩底25.45}

\subsection{量价背离100073124853}

\section{50513 25733 0.000}

\subsubsection{重点提示：头肩底、涨幅满足、头肩顶。}

\section{台化周线2009年8月15日长红突破颈线，约一年的头肩底确立，便可计算}

\subsubsection{涨幅满足。}

\subsection{观察重点一：}

\subsubsection{第一波涨幅满足}

\section{①底部低点25.45元，至颈线约在43.5元，距离18.05元。 ②突破颈线的位置约43.25元+18.05元=等幅计算的涨幅满足61.3元。}

台化盘涨至隔年2010年4月达涨幅满足附近才做较大回档整理，约四个 月的整理过后再过前高时，便可计算第二波涨幅满足。

194

\clearpage

进阶篇

观察重点二：

第二波涨幅满足。 第一波涨幅满足61.3元再加18.05元=等幅计算的第二波涨幅满足79.35元。

台化同年11月达第二波涨幅满足附近后，强势整理约一个月便持续盘 涨，盘涨的结构未出现卖出信号，续做第三波涨幅满足计算。

观察重点三：

第三波涨幅满足。 第二波涨幅满足79.35元+18.05元=等幅计算的第三波涨幅满足97.4元。

台化盘涨至隔年4月达涨幅满足附近才进入盘头，同年8月6日当周爆量 长黑，损破约八个月的颈线，完成头肩顶（也可计算跌幅满足），才结束为 期约两年的大多头走势。

操作方式

2009年8月15日当周突破颈线买单进场，约可买在44.9元附近，设停损在跌 破42元后（跌破颈线确立），风险约6\%，支撑不破看满足，至少有第一波获利 约35\%，通常会在第二波涨幅满足附近，获利了结约75\%，第三波则约110\%。

195

\clearpage

多空转折一手抓 AAA

\subsection{（十七）2008-2012元太日线}

8069元太（日建） M35.00 139.46 456h 163651501.7后

\subsection{保染指特要真：成文： 74,06]}

\subsection{假突破74.06}

\subsection{610/25 66.85}

翻黑破线头部确立59.63 11/16 52.41 45.19

\subsection{34.35 39.73}

\subsection{-30.75}

30.65

\section{-带量长红突破颈线颈线支1 23.53}

\subsubsection{撑压低18.7516.31}

主力进货区出货9.096 73711 -49411 25110

\subsubsection{000090519}

\section{33086 3899.00}

\subsubsection{第一波出货36831-14344}

\subsubsection{第二波出货59318}

-81805 105229

\section{→重点提示：主力颈线支撑区压低出货陷阱、波段跌幅满足、头肩顶。}

\section{元太2008年11月最低9.1元，筑底盘涨至隔年2009年5月5日带量长红，突}

\section{破18元附近颈线，接着连续长红滚量至25元附近都属主力进货区，主力平均成}

\section{本粗估低于20元以下。股价转机题材大涨至2010年1月2日带量长黑，假突破确}

\section{立，库存余额自此一路下降，显示主力长期炒高，开始自高档持续性压低出货}

\subsubsection{（65元附近压低至35元附近），此为第一波出货。}

\section{2011年当年度高获利的转机题材发酵，再拉一波（2010年由亏转盈EPS3.85}

\section{元，2011年大幅成长至EPS6.05元），达涨幅满足后，盘头至10月25日出现带}

\section{量假突破（股价约60元附近）。11月16日翻黑破线，头部确立，库存余额自10}

\section{月25日假突破后便开始一路下降，重点就在盘跌至2009年8月附近及2010年7月}

196

\clearpage

进阶篇

附近及2011年3月附近，三个低点连接形成长达约两年四个月的大颈线，支撑约 在35元附近，主力就在此处震荡出货达约四个月之久。

主力的初期成本在20元以下，第一波高档压低出货，获利数倍，成本大降。 第二波拉高后再度出货，虽然股价自最高66元附近跌至35元区间，看似跌幅已 深，但对主力而言随便卖随便赚（包括公司派也是主力的一种），尤其是法人长 期操控的股票，大涨后最喜好用这种手法出货。不知情的散户以为跌深后支撑区 逢低承接，便掉入其压低出货陷阱。元太震荡后破底长期盘跌，此类跌深至颈线 区，融资持续增加，库存余额持续减少，通常就是法人颈线支撑区压低出货的最 大陷阱区。元太隔年2012年获利急转直下，由高获利转为亏损，春江冷暖鸭先 知，了解结构便能避开风险。

观察重点：

0元太技术面的最佳空点在翻黑破线头部确立时，也是主力第二波出货的起 始点，下图便是头部确立进入下跌波区域的截图。 ②当股价自43.03元反弹后，再跌破反弹低点43.03元，便可计算跌幅满足。 ·自65.31元高点跌至43.03元低点，共跌22.28元。 ·反弹高点52.3元-22.28元=30.02元。

操作方式：

02011年11月16日破线头部确立进场，通常跌破整理半个多月整理区的低 点55.94元，确立翻空，约可空在55.9元附近，停损设59元也就是站回线 上时，风险约5.5\%。 ②元太达跌幅满足附近（最低30.65元）才出现破底翻，企稳作中期反弹， 达跌幅满足附近或破底翻时空单获利了结约37\%以上。

197

\clearpage

多空转折一手抓

1069元太（日）

\subsubsection{2森指要66.68}

\subsection{65.31 -66.68}

\subsection{62.72}

2011/11/16 -58.69 54.67 50.71

\subsection{47.76 46.69}

43.06

\subsection{38.71}

34.68

\subsection{590破底翻30.65}

\section{53993}

\section{32561 21845 11128 0.000}

同样的图形也可以头肩顶作计算，空点相同。（如下图）

\section{2011年11月16日跌破颈线头肩顶确立，便可计算跌幅满足。}

\subsubsection{第一波跌幅满足}

\section{0头部高点66.68元至垂直颈线位置约在52.6元，距离14.08元。 ②跌破颈线的位置约在57.5元-14.08元=等幅计算的跌幅满足43.42元。}

元太达第一波跌幅满足附近（最低43.03元），便开始短线急弹，作逃命反弹。

\subsubsection{第二波跌幅满足 第一波跌幅满足43.42元-14.08元=等幅计算的第二波跌幅满足29.34元。}

\section{元太两波跌幅满足与上图计算的波段跌幅满足，大约相同的位置。}

198

\clearpage

\section{2}

\subsubsection{进阶篇}

99元本（日） 58 [0.3145）日期2011/1/

\subsubsection{货66.68}

\section{头部66.68右肩3 66.68}

\subsubsection{头肩顶}

\subsection{6272}

\subsubsection{左肩2011/11/16}

\subsection{58.69}

54.67

\subsection{52.6 50.78}

\subsection{47.76] 46.69}

\subsection{42.66}

\subsection{34.68}

\subsection{33.11 30.65破底翻30.65}

\section{43277 32561 21845 11128 0.000}

\section{大师语录}

\section{“只有仔细分析失败，才能从中获利；只有惨败才会让人回}

\subsection{到现实。”}

\subsubsection{安德烈·科斯托兰尼（AndréKostolany）}

199

\clearpage

多空转折一手抓 AA

\subsection{（十八）2009裕日车日线}

2227格日（日） 30630低56.80收67 2.39

\subsubsection{[114.00]114.00}

\subsection{104.60}

-95.03 89.75 [82.60] 81.30 85.47

\section{涨幅满足76.70 77.50 76.07}

69.00

\subsection{66.50}

\section{.80 60.90 LAUS24 56.50] 56.93}

\subsubsection{跌幅满足}

\subsubsection{突破前高47.53}

45.40

37.97

\section{28.40] 28.40}

\section{成文量保染指模40 3432.20}

\subsection{2069.80 1388.60 707.40 0.000}

\subsubsection{→重点提示：涨跌幅满足。}

\section{裕日车日线2009年6月26日突破前高63元后，便可计算涨幅满足。}

\section{观察重点一：}

0自28.4元起涨低点至63元，共涨34.6元。 ②自45.4元低点+34.6元=等幅计算的涨幅满足，约在80元附近。 股价最高达82.6元附近，便做约半年的拉回修正，幅度约31\%。

\section{下图裕日车技术面接续上图达跌幅满足附近后，2010年3月3日带量长红，破}

\section{底翻企稳确立。同年4月7日长红突破颈线，隔日4月8日再度长红且带量突破新高}

\section{（前高82.6元），此时便可计算涨幅满足。这次计算的是更大波段的涨幅满足。}

200

\clearpage

\section{2}

\subsubsection{进阶篇}

\subsubsection{[14.00]114.00}

114.00

\subsection{涨幅满足}

\subsection{95.03}

\subsection{85.47}

76.07

\subsubsection{带量长红66.50}

\section{63.00突破颈线}

\subsubsection{带量翻红： -56.93}

\section{5650破底翻}

47.53

\subsection{37.97 28.40 3432.20 2751.00 2069.80 1388.60 707.40 0000}

\subsection{·观察重点二：}

\subsubsection{0自28.4元起涨低点至82.6元，共涨54.2元。 ②自56.5元低点+54.2元=等幅满足计算的涨幅满足约在110.7元附近。}

\subsection{操作方式：}

\section{0操作上，在突破颈线后，不破颈线支撑74元上看波段涨幅满足110.7元附}

\subsubsection{近。}

\section{②裕日车股价达涨幅满足附近后（最高114元）便盘头至同年4月30日长黑}

跌破头部颈线，转弱确立，大幅回档。

201

\clearpage

多空转折一手抓

\subsubsection{（十九）2009-2012宏达电周线}

349奖速（） 1001.54日期2011/4/5

假突破1069.601115.09

\subsection{颈线2011/6/18 971.45}

\subsubsection{带量翻黑跌破颈线及873.29}

\section{电阻式触控电容式触控}

上升趋势线空头确立 775.14

604,29 679.38 581.22

上升趋势线483.07

\section{1/17带量破底翻325.00] 307.50 387.31}

\subsection{285.81] 362.40}

\section{241.03)破底翻V型反转191.00}

91.00

\subsubsection{价跌资增53989}

\section{10333}

\subsection{库存减少-27204-7397.00}

筹码越跌越乱-47010 66816

\section{-86623}

\subsubsection{重点提示：破底翻、假突破、价跌资增。}

宏达电周线图2010年初以前，股价停滞不前，主要原因是当时使用电阻式

\section{触控的操控性低于竞争对手的电容式触控。自从产品改用电容式触控后，股价在}

\section{技术面也出现破底翻的买进信号，所以基本面的改变会反应在技术面上。宏达电 自此一路大涨至千元以上，飙涨约400\%。}

\subsection{观察重点：}

\section{2011年6月开始自高档积极压低出货，6月18日当周带量长黑破线，假突破 确立，果然如当时部落格PO文的预期翻空股价一路崩跌，型态上出现了V型反}

\section{转。通常V型反转大多出现在投机股上，贵为当时股王的宏达电也如投机股般V}

\section{型反转，而且大波段起涨的低点207.08元也被跌破，这种情况通常都是公司长期}

202

\clearpage

\section{2}

进阶篇

衰退的征兆，最好的情况唯有长期整理，以时间换取空间。筹码面明显价跌资增 且库存持续大减，表示筹码越跌越乱，有了威盛的前车之鉴（与威盛同一董事 长，威盛股价最高629元，长期盘跌至最低4.89元，这是不计算减资的价位，如 下图威盛月线），果然不值得信任。

\subsubsection{?操作方式：}

宏达电周线2010年4月3日当周破底翻，确立底部，多单进场，约可买在 276元附近，停损设257元也就是跌落线下时，风险约7\%。2011年6月18日当周 带量翻黑，跌破颈线及上升趋势线，假突破翻空确立，多单获利了结，约可卖在 900元，波段获利约226\%，同时翻多为空，约可空在900元，停损设在940元附 近也就是站回线上时，风险约4.5\%。2012年11月17日当周带量站回线上，破底 翻企稳确立，空单获利回补，约可补在248元，波段获利约72\%。

23期城盛（月） 46.10高4.0041.90收43.70量2778904.008.388\%日期2002/11/30 烫费成交苹量量 集森指福670.38 630.72 -560.04 491.08 422.11 353.15 282.46 -213.50

\subsection{144.54}

\subsection{38.15 28.1531.50 75.58}

\subsection{14.70 4.890}

成交量1256046 832813 423233 0.000 2014

203

\clearpage

多空转折一手抓

（二十）2010苹果日线

FUS简果AFL（日通） 2森振槽294.73294.73 294.93

\subsection{279.011假突破283.26}

2010/6/29

\subsection{271.58}

\subsection{2010/9/1-248.43}

\subsubsection{239.60]破底翻23675}

\section{225.28}

\subsubsection{-213.60 时间波对称 左右各两个月201.93 [199.25] 190.25187.03}

重点提示：假突破、破底翻。

图示苹果2010年在一波多头走势后于当年6月29日再次出现带量翻黑，假 突破确立的卖出信号，也就是非盘即跌。

当卖出信号出现后就会有一段危险期，而危险期可以用对称的时间波来做计 算，苹果型态上假突破确立后高点往左边的整理时间约两个月（图示紫色虚线部 分），此时便可以同样高点计算往右约两个月的对称时间为危险期。

\subsubsection{观察重点一：}

同年8月24苹果在接近危险期的末端跌破整理区下缘颈线，但在对称的危险 期过后，同年9月1日翻回整理区出现破底翻的企稳确立买进信号，量价结构上

204

\clearpage

进阶篇

这次的破底不带量，但翻回整理区的量较大，表示筹码经过整理后已稳定换手， 扭转弱势结构。苹果股价在破底翻企稳买进信号出现后，随之盘涨，突破整理区 上缘，再涨一波。

观察重点二：

因此，当卖出信号出现后大多是非盘即跌，做空间或时间的修正，盘跌是空 间的修正，以盘代跌则是时间的修正。有强劲的基本面，股价才能在涨多后经过 一段整理期消化不安定的筹码，再攻一波。以盘代跌这段时间就是市场对公司是 否持续成长的观望等待期。苹果2008年的假突破后走空大跌约43\%做空间修正 （详见P60～P61），这次2010年的假突破后，拉回修正约15\%，才以破底翻扭 转盘弱结构作时间修正。

操作方式：

操作上出现明显卖出信号时，作必要性的退出才能避开风险（对长期操作者 而言则是减码的依据），直到企稳信号出现时才做进场布局动作。苹果2010年 6月29日假突破卖信后，同年9月1日破底翻前第六个交易日，股价破线，为技 术面的放空点，约可空在239美元附近，停损设252也就是站回线上时，风险约 5.5\%。9月1日带量翻红站回线上破底翻企稳确立，刚好也是时间波两个月危险 期的末端，表示技术面以时间代替空间的修正的扭转空头翻多信号，空单需果决 执行停损进而翻空为多，约可买在252美元附近，停损设低点附近也就是再度跌 落线下时风险约6\%，随后苹果完成时间修正后进入另一波盘涨。

205

\clearpage

多空转折一手抓

\subsubsection{（二十一）2012苹果日线}

\section{US25乘APP（甘要）574高57B.4857038收572.16 10491800室}

\subsubsection{回美森指交量-}

\subsubsection{705.07假突破705.07}

2012/9/25

\subsection{假突破654.00 669.11}

\section{涨幅满足2012/4/13 619.87 632.53}

\subsection{595.94}

\subsection{565.61 570.00 559.98}

\section{526.2以盘代跌拖过危险期}

-523.40

\subsubsection{时间对称[505.75] 486.81}

486.63

450.85

\subsection{12/20 377.68}

\section{交量77 ②价量背离5378970}

\section{①价涨量增③价跌量增4303176}

075794

重点提示：涨幅满足、假突破、量价背离、价跌量增（时间换取空间）。

\section{如图示，2012年第一季苹果一波多头达涨幅满足附近即635美元附近（最高 644美元），涨势告一段落。}

\subsection{观察重点：}

涨幅满足计算 ①第一波低点377.68美元涨至526.2美元=涨148.52美元。 ②第二波低点486.63美元+148.52美元=635.15美元。

\section{苹果价涨量增（数字①）至3月中旬出大量，随后便出现近1个月盘涨的价}

涨量缩量价背离（数字②）现象，股价盘涨至涨幅满足附近（最高644美元），

206

\clearpage

进阶篇

4月13日翻黑，假突破确立，股价进入拉回修正约19\%。 至6月下旬，时间波对称后，股价维持在颈线附近，量缩整理，以盘代跌， 拖过危险期，也就是完成了时间的修正，经过整理后再涨一波。苹果接续股价再 涨一波后，同年9月25日翻黑出现假突破卖出信号，而这次的卖出信号后价跌量 增为筹码渐乱的空头走势（数字③），股价随后大跌。

·操作方式：

苹果2011年12月20日长红突破颈线，为买进信号，多单进场，约可买在 396美元，停损设387美元也就是跌落线下时，风险约2.5\%，达波段涨幅满足附 近或2012年4月13日翻黑假突破卖出信号，多单获利了结约55\%。

小叮：

苹果2012年9月25日这次假突破后，为价跌量增的空头结构，空单进场约 可空在674美元，停损设695美元风险约3.1\%，再以波段跌幅满足计算，满足点 在395美元，空单获利约41\%。苹果自高点705美元大跌至2013年4月19日见波 段低点385美元，才筑底反弹，跌幅达45\%。科技股的股价动力来自于创新，一 旦创新能力不足便会反应在技术面上，反之亦然，微软、思科等皆有前例可循， 这说明股票市场要小心注意风险。

大师语录 投机家的条件：资金、耐心、坚强的神经、独立思考、经验 丰富与敏锐眼光。

安德烈·科斯托兰尼（AndréKostolany）

207

\clearpage

多空转折一手抓

\section{附录：}

\subsection{一、认识投机股结构2012普格日线}

注意量价结构大多可以避开投机股的风险，我操作股票以来，难以避免也接 触过投机股，但都因为坚持严设停损而能避开风险。许多投机的结构大多有迹可 循，认识投机的结构后当量价出现问题时毅然退出，便能避开该股票随后发生的 问题，甚至是退市。

（日）张6/06 7 31.35 7/20 31.35

\subsection{28.10 28.36}

25.29

\subsection{22.22}

17,65] 19.15

\subsection{091 16.0814.58}

\section{13.00] 1301}

轧空利用股东会强制回补12/26 9.943

\subsubsection{6.874}

\section{联霖指槽4.41] 3.730}

\section{成交量融资断头7024.00}

470840

\subsubsection{融券异常6065.29}

\subsubsection{主力库存559.04248.04}

下降-870.04 -1181.04 1505.00

普格日线2012年3月中旬开始大涨，融券也开始增加，随后主力利用台股4

月陆续开始的密集股东会在4月上旬强拉喷出轧空，迫使空头在高档因股东会强 制回补，这已是台湾地区股市几乎每年都会上演的戏码，也突显空头在台湾地区

的不公平性，通常轧完空后会跌回起涨点的股票大多属投机性质。

208

\clearpage

此股最大的问题在同年7月20日带量突破整理区高点后，融券快速增加但价 却不涨的异常现象，融券余额已超过日均量，具备轧空条件，但价不涨反跌便是 异常，属主力放空现象。7月30日股价长黑破线，连续超过5根跌停，大跌爆量 后融券快速获利回补，极为明显的内线交易。（价涨券增通常主力轧散户空，价 跌券增通常是主力做空）

大跌爆量后出现了短线低点17.35元，短线反弹又未过爆量高点20.35元，当 回跌再破17.35元低点的时候，便是多头最后的停损机会（通常下跌大量后的低 点不能再破），也就是持有该股在7月30日跌停之后的连续无量跳空跌停，苦无 卖出机会的，应该在此时毅然退出。

普格10月下旬发布重大信息公司疑遭商业集团诈骗，估计最大影响金额4.6 亿元，以公司5.79亿元股本计算，最高将损失每股纯益7.9元，股价崩跌连续约 20根跌停，11月20日股价最低3.73元爆量打开跌停，融资尤其是券商此时才得 以见量执行断头，融资多头损失惨烈。

跌深反弹后又出现了相当大的问题，12月26日爆量破线头部确立，当日高 点6.95元成为反弹结束后的反压，重点是当时连续两日的库存余额异常大减，通 常都是公司派或大股东积极出脱持股，随后库存余额呈下降趋势，“春江冷暖鸭 先知”，可见此股长期疑虑在股价崩跌后仍不减反增。盘跌后2013年初，虽然 随着大盘跌深反弹，类股概念反弹也未过6.95元反压，同年3月4日利空再度跳 空跌停又一波崩跌，股价可怜必有其可恨之处。

209

\clearpage

多空转折一手抓

\subsection{二、压低出货陷阱-1997-2001思科日线}

\section{P（强） S 1584}

\section{保森指推82.2082.00 9558}

\subsubsection{2000/3/27假突破8215}

70,00] 73.55

64.95

56.49

\subsection{倒货区 3929}

起涨最大支撑30.84

\subsection{也是法人最后出货陷阱2223}

\section{1997/4/22 13.63}

\subsection{[11.04]}

\section{进货区5.031}

\section{225304}

\section{0.000}

\subsubsection{重点提示：假突破、颈线支撑区压低出货。}

\section{全球领先的网络设备厂商思科2000年网络泡沫破灭前当红，股价自1997年}

\section{4月22日低点5.03美元，涨至2000年3月27日82美元（分割前的股价为146.75}

\section{美元），累计涨幅约15倍，突破高档整理区创新高后约2周，便又跌回颈线下极}

\subsubsection{为明显的假突破长线卖出信号。}

\section{股价长期盘跌腰斩，至起涨区企稳反弹，整理约两个月。通常起涨区附近视}

\section{为一波主力作价成本起点，因此为相当大的支撑区，不过一旦法人长线看空，便}

\subsubsection{是最大的压低出货陷阱。}

210

\clearpage

\section{附录}

观察重点：

0注意量价结构上，红色虚框的前后两个大量区位置，1997年4月前后期间 的大量在股价大涨后便视为进货量，当时股价约在6美元上下，2000年高 点长期盘跌至年底跌回起涨大支撑区的大量，在股价续破底大跌后便视为 出货量。 ②当时股价即使腰斩，也仍约在33美元至44美元区间，法人主力长期的筹 码仍是随便卖随便赚，因此除了产业的变化外，需注意量价结构来配合判 断，以免以为捡了便宜，却陷入更大的风险之中。思科长空至2001年出 现13.19美元及11.04美元两个波段低点。

\section{MEMO}

211

\clearpage

多空转折一手抓

\subsubsection{三、崛起的大国中国大陆股市}

P016上海综合指（月通福）期978.94高6005.134857.04收4861.11量8308659-1093.6618.356\%）日期2007/11/30 特· 0.00安费0.00成交0.00万0.00量量0量0 森指楼6124.04 6279.80 OTJO.Z0 5557.65 4991.26 4424.88

\subsection{3858.49}

\subsection{2011/12/31 3277.94}

\subsection{2711.55}

2145.17

\subsubsection{1664.92原始上升趋势线1578.78}

\subsection{2012/12/31带量破底翻998.23}

成交量4211722

\section{14191675 0.000}

\section{2012 2013 2014}

\subsubsection{（一）上证月线}

上证指数月线2011年12月31日当月跌破原始上升趋势线，这对于一个国家 的指数来说是相当罕见的现象，也因此进入长期的低迷。2012年12月31日当月 翻红带量破底翻为技术面长线企稳信号，股价也进入一年多来的打底，大量区间 1949点至2269点为月线的长线支撑区。

重点在一个新兴国家初期的投机高点大多会在10年内再度突破，主要是

因为先走投机，再走基本面高成长的循环，如台湾地区加权指数1990年2月达 12682点投机行情高点后崩跌至2485点，七年多后1997年8月达10256高点（这 是未填权的指数），上证指数2007年10月达6124点投机高点，拉回至今已进入 第七年，配合月线的长线结构来看，仅缺关前整理的量缩动作，第二季以前完成

212

\clearpage

\subsubsection{附录}

\section{底部的机率颇高，当然这一切都必须在技术面不破底部支撑的前提之下，月线的}

\section{颈线在2330点附近突破，尤其是前高2444点突破为底部确立信号，两年多的底}

\subsubsection{部便可确立。}

016上合鱼）12275系19561215

\subsubsection{指.00成之0.00 0.00登05124.04}

\subsection{6133.92}

-5561.09 4988.26

4415.44

3852.49

3210.53

\subsection{2444-80 2270 2706.83}

2004/9\textasciitilde{}2005/12 -2143.88

\subsection{1949.461849.65 1984.821571.06}

\subsubsection{2006/1/7突破颈线完成大底后盘涨至喷出9/14 998.23}

交量

\subsection{（二）上证周线}

\section{如图示，上证周线2004年9月至2005年12月约一年多的底部，在2006年1}

\section{月7日当周突破出现技术面的买进信号，随后盘涨至喷出进入数倍的投机行情，}

\section{注意量价结构上大量后的过高都是一波多头行情，反之涨多爆量后的盘头都是一}

\section{波盘跌空头行情，最近一次的爆量是2013年9月14日当周的高点2270点，表示}

\section{未来若支撑不破而突破2270点就是一波多头的开始，技术面在稳固的双底支撑}

\subsubsection{下风险有限，注意支撑即可。}

\subsubsection{213}

\clearpage

多空转折一手抓

6出增合当（日） 220.52205.92310.59 1000 2指安限0.00 2444.80 0.000

\section{2444.80假突破2444.80}

\subsection{2334.33}

\subsubsection{天量区高点12/4 2312.93}

\subsection{2260.87}

\subsection{2198.85 2181.05}

\subsection{2120 2136.4}

\section{2092.87 2114.54}

个2078.99

/102048.03

\section{长线支撑区破底击}

\subsection{1982.67}

\subsection{1949 1916.16}

1849.65

\section{488600}

2013/1

\subsubsection{（三）上证日线}

上证日线曾在2013年12月14日微幅突破颈线，可惜未突破大量区高点2270 点做确立，且在五个交易日后跌落线下，主要是中国大陆钱荒的利空做再度压盘

\section{动作，而中国大陆钱荒是官方整顿太过猖的地下金融的动作，人为的干预也得}

以控制技术面不再破1949点低点支撑。由这个支撑及2013年7月低档震荡的相 对整理区间，得以画出1949点至2070点为日线支撑区，一旦支撑不破而突破颈 线，则中国大陆股市将是各路英雄大显身手的地方。

至于这次上证日线结构在1995年至1996年初的台股也有类似的结构，记得 1988年9月24日当时“财政部长”郭婉容宣布复征证所税，台股无量崩跌19天， 但当时是由波段低点2241点涨数倍后达8813点涨高的结果，1996年1月4日

\section{“立法院”三读通过复征证所税，如图示台股这次只跌1天且并非无量也未再破}

214

\clearpage

\subsubsection{附录}

底，因为这次是自7180高点长期盘跌后的打底结构完全不同，只是将原来约五 个月的双底结构，再次压盘多打底约三个月，延长为约八个月的三重底，这次利 空测试的底部完成后，走了一年多的多头直奔10256高点才回头。

\section{0000天2回彩指6191.16}

\subsubsection{6191.16-6191.16}

6001.12

\subsubsection{带量长红突破颈线5811.07}

\subsection{5670.50底部完成}

\section{大陆试射飞弹确立多头讯号、 5621.03}

5430.98

5237.72

\subsection{5047.67}

4857.63

\section{预估的双底4734.55 4672.67 4667.594619.27}

\section{复征证所税利空}

\subsubsection{4530.96跳空大跌后扩底4474.32}

\section{1125134 8438507}

底部都是利空测试出来的，以上证A股的结构而言，无论是否会再破底，非 常适合现金在底部买入长期持有，再定期定额续买至未来喷出为止，若是已熟读 技术分析者便可运用所长，预估将来中国大陆股市将是技术分析的天堂。

215

\clearpage

多空转折一手抓

\section{四、关于A股的看法}

很高兴也很荣幸本书受广东经济出版社的青得以在中国大陆发行。记得 《多空转折一手抓》2014年3月在台湾地区出版时，我特地赶了一篇“崛起的大 国中国”作为结尾，预测A股当年为6124高点反转以来的第七年长线转折，A股 也如期在当年喷出。 除了掌握长线喷出转折，2015年四月以后也利用波段涨幅满足掌握高点成 功逃顶，而初升段结束大幅回档后，接下来就是主升段的进行，就长线的结构来 看，6124的高点十年内会再来一次甚至突破，今年2016年已是第九年，因此A 股低档布局风险有限，又一次人弃我取的进场时机，相对至明年底前应会有相当 惊人的报酬率，正如书的结尾所提“中国大陆股市将是技术分析的天堂”。 所谓看不懂不做，行情并非随时都能看得懂，必须等待型态完成结构完整 才容易做后续的判断，“等待是一种艺术”，大波段行情都是等出来的，当型 态完成后，量与价的变化会告诉你进场时机，因为主力作价是真金白银的进场 操作，所以量必然会有变化，量的变化就会牵动价的波动，因此研判量价变化 的结构便能判读主力动向与意图，配合此书在大陆上市编制陆股的一些案例， 由于时间关系仅是信手捻来提供投资人参考，入门者多研读历史K线量价关系 多画线便能逐渐了解，已有基础者对本书应不难上手。预祝大家在中国经济发 展逐渐成熟后，股市进入难得的黄金十年都能利用最有效用的技术分析掌握多 空大波段转折契机。

216

\clearpage

\subsubsection{附录}

上证指数（周线前复权）设置均线

\subsection{笔者截稿时（2016/07/07）上证正达对称时间波第九周的变盘3600.00}

转折，这是逾半年底型态的长线转折，往上变盘后完成底部

\subsection{进入主升段起涨的机率相当高，只要注意不再跌破2822点至3400.00}

\subsubsection{2922点重要的支撑区即可，低档风险有限，往上空间无限}

3200.00

\subsubsection{上证周线}

\subsection{9周9周3000.00}

2800.00

\subsection{时间波对称}

\section{[VOL(5.10)VOLUME:868568786.000↑MA/OL1:732725867.400MAVOL2:703760197100↑263.30改梦数加指标换指标区20.00}

\subsection{0.00}

\section{2016 02 06}

亚更股份（日线前震权）

\section{波段涨幅满足32.22}

\section{亚夏股份日线32.00}

\subsection{28.00}

\subsection{26.00}

24:00 22.00 20.00

\subsection{18.00}

\subsection{14.00}

\subsection{MAVOL1:101266.200MAVOL2100995.200梦数加指标换指标× 113.00}

\subsection{37.30}

\section{Douhnuont 0.00}

\section{07}

\subsubsection{217}

\clearpage

\subsubsection{多空转折一手抓}

上证50（日线不复权） 3494

\subsubsection{波段涨幅满足3439}

\subsubsection{上证50日线}

\section{2800.00}

\subsection{2400.00}

\section{2200.00}

\subsubsection{收敛三角形整理000002}

1800.00

\section{IAVOL1:23717661.800M4/DL2-22001088300改参数加搭标换指标600t02.69}

133

\section{0.00}

上证指数（日线前复权）设置均线

\section{上证日线时间波对称}

第12天变盘305000 3000.00 2950.00

\section{12天12天}

2700.00

\section{VOL(5.10)VOLUME:182166748.000MAVOL1198601095.200.MA/0L2186209711.900改梦载加指标换指2650003.52}

\subsection{0:00}

\subsubsection{218}

\clearpage

\subsubsection{附录}

山东黄金（周线航复权） 67.22消股通融资融券设置均战

\section{三尊头}

\subsubsection{山东黄金周线头肩顶=三尊头}

\subsection{三尊头不=头肩顶（因为头在中间最高） 60.00}

\subsection{以头部高点至颈线的距离的方式做计算5000}

\subsubsection{40.2颈线40.00}

30.00

\subsection{跌幅满足16.98 20.00}

10.00

\section{VOL(5.10)VOLUME799837.000MAVOL1:575207.200↑MAVOL2:428631.600改梦数加指标换指标360.00}

\subsection{119.00 0.00}

中金黄金（日线前复权）沪股通融资融券设置均线

\subsection{17 M头18.00}

\subsubsection{中金黄金日线}

\subsection{14.00 12:00 10.00}

8.00

6.00

跌幅满足3.2 4.00 VOL(5.10)VOLUME:81797.000↑MAVOL1:75490.000MAVOL2:53278.500↑ 3.14 25:50

\subsection{8.41}

\section{11 0.00}

219

\clearpage

多空转折一手抓

中航三鑫（日线前权）

\section{中航三鑫日线15.4 M头15.00}

14.00

\subsection{13.00}

12.00

\subsection{11.4颈线t1.00}

10.00

\subsection{跌幅满足7. 8.00}

7.00

\section{2VOL（5.10）VOLUME:411586.000MA/OL1:441204.600↑M/OL2340835.700改梦数加指标换指标区64.90}

\subsection{28.00 0.00}

\section{天泽信幕（日线前权） 51.7设置均线}

\subsubsection{头肩顶50.00}

\subsubsection{天津信息日线}

45.00

\subsubsection{右肩}

\subsection{左肩40.00}

35.00

\subsection{33.2颈线35.3}

\subsection{25.00 20:00}

\subsection{跌幅满足16.8 15.00}

\section{[VOL(5.10)VOLUME:187590.000MAVOL1:195486.600↑MA/OL2131693.400改参数加指标换指标21.60}

\subsection{7.13}

miooth 0.00

\section{0}

220

\clearpage

\subsection{附录}

乐普医疗（日线不复权）

\section{乐普医疗日线假突破}

55.00

\subsubsection{波段涨幅满足50.05 50.00}

45:00

\subsection{40.00}

35.00

\subsubsection{33.1}

30.00

25.00

\section{VOL(5.10)VOLUME:132572.000 22.51改参数加指标换指标× 23.50}

\subsection{7.76}

\section{02 05 06 07 0.00}

乐普医疗（日线前复权）融海融券设置均线

\subsection{乐普医疗日线 42.00 40.00 38:00}

\section{35破线36.00}

\subsection{32.00 30.00}

\subsection{27.3 28.00}

\section{VOL（5.10)VOLUME:23659.000↓MAVOL1:24676.600MAVOL2:30546.500↑改萝数加指标换指标区] 24.30}

\subsection{8.00}

\section{0.00》}

\subsubsection{221}

\clearpage

\subsubsection{多空转折一手抓}

\subsubsection{宁波东力（日线耐复权） 20.75}

\subsection{头肩顶20:00}

\subsection{宁波东力日线}

18.00

\section{右肩15.00}

\section{左肩13.6 14.00}

\subsection{12.00}

\section{颈线}

10.00

\subsubsection{跌幅满足5.45 6.00}

\section{VOL(5.10)VOLLME:377568.000MAVOL1:327355.600MAVOL2:298946.400改参载加指标换指标区58.10}

\subsection{19.10 0.00 2015}

长江电力（日线前菱权）泸联通随资随账设置均线·

收敛三角形头部14.50 长江电力日线 14:00

\subsection{13:00}

12.50

12.00

\section{跌幅满足11.4破底翻11.50}

\section{VOL(5.10)VOLUME:254238.000MA/OL1:197992.400MAVOL2:157927.40011\_02 11.13改参数加指标换指标513.00}

\subsection{169.00 0.00}

\section{2016 02 04}

222

\clearpage

\section{附录}

\section{上证指骤（日线不发权）假突破上证日线设均}

0169-517 5000.00

4500.00

4000.00

涨幅满足3600

\subsubsection{破线3310 500.00}

\section{门+03190破线}

\subsection{VOL(5.10)VOLUME133439291.000↓MAVOL1:169613201.000MA/OL2 2900收敛三角形改梦数加指标换指标× 8.55}

4.22

\section{2016 0.00}

6日和15分钟30分钟60分钟日周线月线更多周期免加画线离约照开4

\section{创业板指（日线前发权）创业板指日线2915头肩顶设置均线}

右肩

涨幅满足-2805 2800.00

\section{左肩2600.00}

\section{2535颈线}

2400-00

\subsubsection{颈线2320}

\section{2135 2200.00}

跌幅满足1940 2000.00

\section{1780一W底1800.00}

\section{VOL(5.10）VOLUME:36587329.000MAVOL1:30837613.600MVOL232356508700改参数加指标换指标7416.00}

657:00 0.00

\section{02}

223

\clearpage

\subsubsection{多空转折一手抓}

\section{泸深300（日线前发权） 2780设置均线}

\section{沪深300日线收敛三角形头部}

\subsection{2600.00}

\section{2400.00}

\subsubsection{2380}

\subsection{2200.00}

\section{波段跌幅满足2050 2023,17 2100.00}

\section{2VOL(5.10)VOLUME157783247.000 12270改参款加指标换指标1.67}

8346.00

\section{20 03 05 06 0.00回}

\section{沪300（周线前权） 4850 M头量均线}

\subsection{沪深300周线}

\subsection{5060.00}

\section{4500.00}

4000.00

\section{3500.00}

3000.00

\subsection{波段跌幅满足2000.00}

\section{VOL(5.10)VOLUME:438049182.000↓MA/OL1:454175896.400↑MA/OL2400659435.600改梦数加指标换指标文6.54}

3.27

\section{09 10 12 2008 02 03 04 05 07 DB 09 10 12 2009Z0 03 04 0.00}

224

\clearpage

\subsection{附录}

泸第300（周线前发权）设置均线

\section{沪深300周线假突破}

\section{波段涨幅满足4897}

\subsection{口4500.00}

\section{00-0001}

宁500.00

\section{3285下飘旗3000.00}

2500.00

\subsubsection{改萝款加指标换指标29.10}

14.60

\section{口EO 04 05 06 07 08 09 12 2015 02 03 04 05 06 07. 08 0.00}

元证券独资驰券设置均统

28.00

\section{国元证券日线第二波涨幅满足26.3}

\section{中点口26.00}

\section{中十中口}

24.00

\subsubsection{第一波涨幅满足22 22.00}

\subsection{20.00}

18.00

16.00

\subsection{VOL(5.10)VOLUME:402951.000MA/OL1:564739.000↓A/OL2:730622.400收敛三角形改梦数加指标换指标区162.00}

105:00 53.40 0.00

225

\clearpage

\section{www.guziyuan.cn}

多空转折一手抓

分时5日5分钟15分钟30分钟60分钟周线月线更多周期加画线简约银开 罗牛山（日线前复校）设营均线

\subsection{梦牛山日线}

\subsubsection{涨幅满足8.4 098}

\subsection{7.50}

\section{6.8颈线6.6 7.00}

\subsection{中中中口6.50}

6.00

\section{右肩}

\section{左肩5.50}

\section{2VOL(5.10)VOLUME:345257.0005.06改梦数加指标换指标区185.00}

\section{头肩底120.00}

61.10

\section{L.1 2016 0.00}

深圳惠程（日线前萱权） 24.00

\section{深圳惠程日线假突破}

22:00

波段涨幅满足20.97

\subsection{18.00 16.16 16.00 14.00}

\subsection{下飘旗型13.09+7.88 12.00}

10.00

VOL(5.10)VOLUME:508482.000MA/OL1:640242.200MA/OL2613760.1002改参数加指标换指标区170.00 8.28 110.00 56.00

\subsubsection{226}

\clearpage

\subsubsection{附录}

有色（日线前复权）融货路务设置均线

\section{铜陵有色日线收敛三角形头部4:40}

\section{4.00 3.80}

\subsection{3.6 3.60}

\subsection{3.20 DOE}

\subsection{跌幅满足2.7 2.60}

\section{VOL(5.10)VOLUME2085880.000MAVOL1:3556926.400MAVOL2改美数加指标换指标区2082.00}

1029.00

\section{10 2016 02 03 0.00}

黑艺（日线不权） 25.55

黑芝麻日线

\subsection{22.00}

\subsection{上飘旗20.00}

\section{/}

18.00

\subsection{16.00}

\subsection{12:00}

波段跌幅满足8.7 10.00 VOL(5.10)VOLUME:73107.000↓MA/OL1:109208.600MA/OL2:96545000改参数加指标换指标50.60

32:80 16.70

\subsubsection{0.0020}

227

\clearpage

\subsubsection{多空转折一手抓}

\subsubsection{XD陆天高（日线前复权） 8883383}

\section{楚天高速日线M头10.1 10.00}

\subsection{8.00}

\subsubsection{颈线7.00}

009

\subsubsection{跌幅满足4.8 5.00}

\subsection{3.00}

\section{7VOL(5.10)VOLUME:306377.000MAVOL1:488617.60024改数加指标换指标249.00}

162:00 8210

沪深300（周线前复权）

\section{波段涨幅满足3792跌破前面三周低点}

\section{沪深300周线3500.00}

\section{3000.00}

\section{2500.00}

\section{2200多头波向上喷出时可配合三周低点}

\section{（日线则三日低点）是否跌破作为强弱势的参考}

2000.00

\subsection{左肩右肩}

\section{2VOL(5.10)VOLUME:405203196.000 1606头肩底改参数加指标换指标7.20}

3,60

\section{80 09 10 12 2009 02 03 04 05 06 07 08 09 10 11 0.00}

\subsubsection{228}

\clearpage

\subsection{附录}

\section{300（周线的权） 3550设置均线}

\section{收敛三角形头部沪深300周线}

3400.00

3200.00

\section{3000.00}

\section{2800.00}

2500.00

\subsection{波段跌幅满足2380 2400.00}

\section{VOL(5,10)VOLUME:457713699.000↓M/OL1:458108599.800MAVOL2402525887.300改梦款加指标换指标区8.04}

\subsection{000}

\section{21 2011 02 03 0-4 05 06 07 LI 2012 02 03}

\subsubsection{229}

\clearpage

多空转折一手抓

\subsection{后记}

我40岁退休，两年后生活渐感无聊，在朋友的劝说下又复出两三年，进入 一家竞争力颇强的公司，享有千万业绩分析师的光环。随着投顾生态的改变，以 及对事业已无企图心，很快地两三年后又再度退休。

2007年开始撰写部落格（部落格名称：蔡森一随势而为一技术分析），当时 用意是要让当初信任我的投资人，了解我对盘势的看法。每日点阅人数从开始的 一两百人成长到目前五六千人次以上，其成长主因是电视媒体邀约曝光的关系， 这几年每周约一天固定接东森财经“57金钱爆”通告，偶尔也有演讲邀约，加 上自已使用的软件程式对外销售，会不定期办说明会及技术分析课程，目前为止 非常适应这种半退休式的生活。2014年开始积极介入A股市场，希望未来三五年 内能在A股股市再创一次高峰。

以前股票数量不多，每个交易日可从头检视全部个股技术面一遍再挑股票。 自1996年，也就是完成4474低点后的底部那波大多头行情，两次到达万点以上 之后，股票数量大增。2001年初，不知是股票数量太多，还是一心想退休，第 一次发生挑股票到打瞌睡的情况，我便萌生将自已选股逻辑转成软件的念头，而 且考量当时股票软件已日渐成熟，遂请一家软件公司按照我的需求，写了一套程 式。让我可以利用程式中先作一次多空选股后，再作精选。虽然经过软件程式的 帮助确实省却相当多的时间，但还是要靠自已的经验才能做出较正确的判断，本 书是我多年的操作心得毫无保留，内容分成两部分，第一部分讲解12招的多空 判断方法，第二部分，则是第一部分12招多空判断方法的延伸，包括实战运用 及操作策略的说明。企盼透过本书的内容，可以帮助更多的投资者走出股市的迷 雾森林，开创一条属于自已的赚钱明路。

230

\clearpage

\subsubsection{跟大户主力一起进场，}

\subsubsection{轻松赚饱荷包.}

本书特点： 波段随势操作，严守停损一操作不需盯盘 图解12招多空判定法→简单易懂易学 实战线图搭配操作说明一尽书大师智慧 停损停利的规划解释一获利诀窍大公开

面对瞬息万变的金融市场，观念和时机是获 利的关键，因此建立正确的观念，以及准确掌握时 机，才能在实战中小赔大赚，累积出可观获利。本 书淬炼出蔡森29年的投资经验中，最重要的观念 与操作时机的判断方法。不教困难的指标，更不使 用复杂的交易法则，仅用量价关系配合型态做判 断，蔡森大方公开自己12大招获利法，以及实战 的操作策略，让你认识方法后，马上可以在实战的 例题中得到验证，进而轻松掌握未来行情。

\section{ISBN978-7-5454-4764-4}

\section{定价：68.00元}

\end{document}

